\documentclass[11pt,a4paper]{article}

% --- Packages ---
\usepackage[margin=2.5cm]{geometry}
\setlength{\headheight}{14pt}
\usepackage{hyperref}
\usepackage{titlesec}
\usepackage{enumitem}
\usepackage{booktabs}
\usepackage{longtable}
\usepackage{xcolor}
\usepackage{fancyhdr}
\usepackage{tikz}
\usetikzlibrary{shapes.geometric, arrows.meta, positioning, fit, backgrounds}
\usepackage{listings}
\usepackage{parskip}
\usepackage{microtype}
\usepackage{amsmath}
\usepackage{amssymb}
\usepackage{mathtools}
\usepackage{amsthm}
\usepackage{array}
\newcolumntype{L}[1]{>{\raggedright\arraybackslash}p{#1}}

% --- Style ---
\hypersetup{colorlinks=true, linkcolor=blue!60!black, urlcolor=blue!60!black, citecolor=blue!60!black}
\titleformat{\section}{\large\bfseries}{}{0em}{}[\titlerule]
\titleformat{\subsection}{\normalsize\bfseries}{}{0em}{}
\titleformat{\subsubsection}{\normalsize\itshape}{}{0em}{}
\pagestyle{fancy}
\fancyhf{}
\rhead{\textit{Fleet Management --- SRS}}
\lhead{\textit{v1.0 $\cdot$ \today}}
\cfoot{\thepage}

\definecolor{todo}{RGB}{180,0,0}
\newcommand{\TODO}[1]{\textcolor{todo}{\textbf{[TO DO: #1]}}}

\lstset{
  basicstyle=\ttfamily\small,
  breaklines=true,
  frame=single,
  backgroundcolor=\color{gray!8},
  rulecolor=\color{gray!40}
}

% --- Statements ---
\theoremstyle{plain}
\newtheorem{theorem}{Theorem}[section]
\newtheorem{proposition}[theorem]{Proposition}
\newtheorem{lemma}[theorem]{Lemma}
\newtheorem{corollary}[theorem]{Corollary}
\theoremstyle{definition}
\newtheorem{definition}[theorem]{Definition}

% --- Notation shared by every section ---
\newcommand{\ARDinf}{\mathrm{ARD}_{\infty}}
\newcommand{\ARDone}{\mathrm{ARD}_{1}}
\newcommand{\inc}{\mathrm{inc}}
\newcommand{\new}{\mathrm{new}}
\newcommand{\tr}{\mathrm{tr}}

% --- Table rows (input and output tables) ---
\renewcommand{\_}{\textunderscore\allowbreak{}}
\providecommand{\keyrow}[5]{{\raggedright #1\par} & {\raggedright #2\par} & {\raggedright #3\par} & {\raggedright #4\par} & {\raggedright #5\par} \\}
\providecommand{\optrow}[4]{{\raggedright #1\par} & {\raggedright #2\par} & {\raggedright #3\par} & {\raggedright #4\par} \\}
\providecommand{\rejrow}[2]{{\raggedright #1\par} & {\raggedright #2\par} \\}
\providecommand{\typerow}[3]{{\raggedright #1\par} & {\raggedright #2\par} & {\raggedright #3\par} \\}

\begin{document}

% --- Title page ---
\begin{titlepage}
  \centering
  \vspace*{3cm}
  {\Huge\bfseries Fleet Management\par}
  \vspace{0.5cm}
  {\Large Software Requirements Specification\par}
  \vspace{0.3cm}
  {\large Multi-Component Stochastic Degradation MILP\par}
  \vspace{1cm}
  {\large Version 1.0 $\cdot$ \today\par}
  \vfill
  {\small Prepared with spec-creator $\cdot$ Intended for Claude Code implementation}
\end{titlepage}

\tableofcontents
\newpage

% ================================================================
\section{Project Overview \& Goals}
\label{sec:overview}
% ================================================================

The Python package \texttt{fleet\_management} schedules a fleet of $F$ trains over a
horizon of $T$ days. Every day, $M$ missions must be flown. Each train flies at most
one mission, spends the day in the depot, or stays idle. Each train has $L$
components, and each component degrades stochastically. On a depot day a component
may be repaired (imperfect repair) or replaced. The schedule is a mixed-integer
program solved with Gurobi.

\textbf{Goal.} Minimise the cost of depot days, repairs and replacements, plus a
penalty on the largest train damage (Section~\ref{sec:objective}), under a
reliability requirement: for every component and every day,
$\Pr(D_{i\ell k} > \tau_{i\ell}) \leq \varepsilon_{i\ell}$, with $D_{i\ell k}$ the
degradation and $\tau_{i\ell}$ the failure threshold. The horizon has two phases. The
transitory phase ($H_1$ days) starts from the current state of the fleet. The
operating phase ($H_2$ days) is the schedule to repeat. The requirement holds on every
day of the horizon and of every later repetition of the operating phase
(Section~\ref{sec:loop}).

\textbf{Degradation models.} Each cell (component $\ell$ of train $i$) has its own
model, chosen by the input \texttt{model}:
\begin{itemize}[nosep,leftmargin=1.5em]
  \item \textbf{Gaussian}: a Gaussian random walk; tracks the mean and the variance,
        and a repair reduces both (Section~\ref{sec:dyn_gaussian}).
  \item \textbf{Inverse Gaussian (IG)}: tracks the mean; reliability by a cap on the
        mean (Section~\ref{sec:dyn_ig}).
  \item \textbf{Wiener}: Brownian motion with drift; tracks the mean and the
        variance. The variance grows by $\sigma^2$ per mission day; a repair acts
        on the mean only, and only a replacement lowers the variance
        (Section~\ref{sec:dyn_wiener}).
  \item \textbf{Gamma}: Gamma increments with mission-dependent rates; tracks the
        mean; reliability by a cap that is safe for every mix of rates
        (Section~\ref{sec:dyn_gamma}).
  \item \textbf{Rainflow}: distribution-free. The increment law is measured by
        rainflow counting of load histories and enters only through summaries
        (mean, variance, support width, cumulant generating function). Reliability
        is certified by one of five tail bounds: Markov, Cantelli, Hoeffding,
        Bernstein or Chernoff (Section~\ref{sec:dyn_rainflow}).
\end{itemize}
Both the Gaussian and the Wiener model use Gaussian increments. The Gaussian model
tracks a variance that a repair reduces. The Wiener model fits components whose
uncertainty does not depend on the route and is not reduced by a repair. When in
doubt, use the Gaussian model. The rainflow model makes no distributional
assumption: its certificate holds for every increment law with the measured
summaries, at the price of a looser bound than an exact distribution function.

\textbf{Repair and replacement.} Each cell has a repair law, $\ARDinf$ (the default)
or $\ARDone$ (Section~\ref{sec:repair_laws}); every model accepts both. Each cell may
be replaced (Section~\ref{sec:replacement}).
The loop row forces at least one replacement in the operating phase in two cases.
This is not an input check; a schedule without the replacement is infeasible.
\begin{itemize}[nosep,leftmargin=1.5em]
  \item A Wiener cell whose train flies at least one mission in the operating phase:
        its variance falls only by replacement.
  \item An $\ARDone$ cell with $\rho_{i\ell} < 1$ that accumulates mean damage in the
        operating phase: a repair never removes the share $1 - \rho$ of the damage
        since the last intervention (Proposition~\ref{prop:a_ard1_replacement}).
\end{itemize}

\textbf{Problem class.} The problem is a mixed-integer linear program (MILP) when every
reliability row is linear and the damage penalty is \texttt{inf\_norm}. Otherwise it
is a mixed-integer quadratically constrained program (MIQCP). It is nonconvex when a
quadratic reliability row uses the \texttt{exact} form
(Section~\ref{sec:linear_forms}); it is convex when the only quadratic row is the
\texttt{quadratic} damage penalty (Section~\ref{sec:damage_reg}).

\textbf{Scope of the package.} The package solves one schedule (\texttt{solve}),
sweeps the horizon lengths (Section~\ref{sec:sweeps}), starts a solve from a
previous solution or a heuristic (Section~\ref{sec:warm_starts}), plots a schedule
(Section~\ref{sec:plotting}), and checks a result against its input
(Section~\ref{sec:validation}). It reads inputs in YAML, JSON or HDF5
(Section~\ref{sec:schema}) and writes the outputs of Section~\ref{sec:output}.

\textbf{Authority.} This specification is authoritative for the formulation and for
the tools. The code follows it: every change is written here first.

\textbf{Non-goals.}
\begin{itemize}[nosep,leftmargin=1.5em]
  \item Routing decisions are not optimised: the missions are inputs.
  \item Continuous-time formulations are out of scope: all dynamics are
        discrete-time, with a step of one day.
\end{itemize}

% ================================================================
\section{Index Sets and Notation}
\label{sec:notation}
% ================================================================

\begingroup\small
\begin{longtable}{@{}p{4.6cm}p{8.2cm}p{2.2cm}@{}}
\toprule
Symbol & Meaning & Defined in \\
\midrule
\endhead
\multicolumn{3}{@{}l}{\emph{Indices and horizon}} \\
$i \in \{1, \ldots, F\}$ & train & \\
$\ell \in \{1, \ldots, L\}$ & component of a train & \\
$(i, \ell)$ & a \emph{cell}: component $\ell$ of train $i$ & \\
$j \in \{0, 1, \ldots, M\}$ & activity: $j = 0$ is a depot day, $j \geq 1$ a mission & \\
$k \in \{1, \ldots, T\}$ & day; index $k = 0$ holds the initial state & \\
$H = [H_1, H_2]$, $T = H_1 + H_2$ & horizon input; a scalar $H$ means $H_1 = H_2 = H$ &
  \ref{sec:input_scalars} \\
days $1, \ldots, H_1$ & transitory phase & \\
days $H_1 + 1, \ldots, T$ & operating phase, repeated without end in the endless
  schedule & \ref{sec:loop} \\
$p \in \{1, \ldots, P\}$ & segment of the \texttt{pwl} form, $P = \texttt{pwl\_points}$ &
  \ref{sec:linear_forms} \\
\midrule
\multicolumn{3}{@{}l}{\emph{Parameters, increments and states}} \\
$\tau_{i\ell}$, $\varepsilon_{i\ell}$, $\rho_{i\ell}$ & failure threshold,
  reliability level, repair efficiency of a cell; a model section may drop the
  subscript $i\ell$ after saying so & \ref{sec:inputs} \\
$C_M$; $C_{R,i\ell}$, $C_{\mathrm{rep},i\ell}$; $C_D$ & costs: depot day; repair,
  replacement; damage penalty & \ref{sec:objective} \\
$\mu^{\inc}_{i\ell jk}$, $v^{\inc}_{i\ell jk}$, $b^{\inc}_{i\ell jk}$,
  $K^{\inc}_{i\ell jk}$, $\beta_{i\ell jk}$ & mean, variance, support width and
  cumulant generating function (CGF) of the increment of cell $(i,\ell)$ when train $i$
  flies mission $j$ on day $k$, after the phase lookup; Gamma rate &
  \ref{sec:input_increments} \\
$\delta\mu_{i\ell k} = \sum_{j=1}^{M} x_{ijk}\, \mu^{\inc}_{i\ell jk}$ & day increment
  of the mean; $\delta v$, $\delta R$ (from $(b^{\inc})^2$) and $\delta K$ likewise; $0$
  on depot and idle days & \\
$\mu_{i\ell k}$, $v_{i\ell k}$, $R_{i\ell k}$, $K_{i\ell k}$ & tracked mean,
  variance, Hoeffding budget, Chernoff quantity & \ref{sec:variables} \\
$g^{\mu}_{i\ell k}$, $g^{v}_{i\ell k}$, $g^{R}_{i\ell k}$, $g^{K}_{i\ell k}$ &
  anchors (law $\ARDone$) & \ref{sec:repair_laws} \\
$\mu_{0,i\ell}$, $v_{0,i\ell}$, $R_{0,i\ell}$, $K_{0,i\ell}$ & initial values &
  \ref{sec:inputs} \\
$\mu^{\new}_{i\ell}$, $v^{\new}_{i\ell}$, $R^{\new}_{i\ell}$, $K^{\new}_{i\ell}$ &
  replacement values & \ref{sec:replacement} \\
$b^{0}_{i\ell}$, $b^{\new}_{i\ell}$ & support widths of the initial and replacement
  damage & \ref{sec:dyn_rainflow} \\
$\ARDinf$, $\ARDone$; $a$; $\psi$ & repair laws; repair factor; repair map &
  \ref{sec:repair_laws} \\
$\xi_k$, $\zeta_k$, $\varphi_k$, $\Phi$, $\mathcal{C}_{i\ell}$ & full tracked state,
  reference trajectory, day map, map of the operating phase, certificate set &
  \ref{sec:encoding}, \ref{sec:loop} \\
\midrule
\multicolumn{3}{@{}l}{\emph{Derived constants}} \\
$M_s$ & natural range of a tracked quantity $s$ (big-M constant) &
  \ref{sec:encoding} \\
$\kappa_G = \Phi^{-1}(1 - \varepsilon)$ & standard normal quantile & \ref{sec:dyn_gaussian} \\
$V_{\max} = \tau^2 / \kappa_G^2$ & variance range, Gaussian and
  Wiener & \ref{sec:dyn_gaussian} \\
$\eta_{i\ell}$; $Z \sim \mathrm{IG}(1, 1/\eta)$; $q$; $\bar\mu = \tau /
  \mathbb{E}[Z \mid Z > q]$ & inverse Gaussian dispersion; unit law; its
  $(1-\varepsilon)$-quantile; mean cap & \ref{sec:dyn_ig} \\
$\beta^*_{i\ell}$; $t = \beta^*\tau$; $A_1(t, \varepsilon)$; $Q(a, x)$ & smallest
  Gamma rate among the terms of the cell with a positive mean; scaled threshold;
  shape cap; regularised upper incomplete gamma function & \ref{sec:dyn_gamma} \\
$\kappa_C = \sqrt{(1-\varepsilon)/\varepsilon}$ & Cantelli coefficient &
  \ref{sec:dyn_rainflow} \\
$\kappa = \ln(1/\varepsilon)$ & exponent constant (Hoeffding, Bernstein) &
  \ref{sec:dyn_rainflow} \\
$b_{i\ell}$, $c_{i\ell} = \kappa\, b_{i\ell}/3$ & Bernstein support constant and margin &
  \ref{sec:dyn_rainflow} \\
$s_{i\ell}$ & Chernoff parameter (input \texttt{s\_chernoff}) & \ref{sec:dyn_rainflow} \\
$V^{\mathrm{C}}$, $V^{\mathrm{B}}$ & variance ranges, rainflow Cantelli and Bernstein &
  \ref{sec:dyn_rainflow} \\
$\mu_r = \texttt{tangent\_ref}\cdot\tau$; $c_2$, $c_1$ & tangent point; constants of a
  quadratic row & \ref{sec:linear_forms} \\
$E$ & evaluation horizon (input \texttt{evaluation\_horizon}) & \ref{sec:objective} \\
\bottomrule
\end{longtable}
\endgroup
\section{Input Specification}\label{sec:inputs}

An input is one mapping from key names to values.
This section lists every key, in groups, with its shape, its default and the cells that read it.
Section~\ref{sec:schema} gives the file formats and the keys that each cell type needs.

\textbf{Cells and abbreviations.}
A \emph{cell} is a pair $(i,\ell)$ of a train and a component.
Each cell has its own degradation model, repair law, reliability bound, linear form and encoding.
The column \emph{Applies to} names the cells that read a key.
It uses these abbreviations: G (Gaussian), IG (inverse Gaussian), W (Wiener), Ga (Gamma), RF (rainflow).
For a rainflow cell it names the tail bound: Mk (Markov), Ca (Cantelli), Ho (Hoeffding), Be (Bernstein), Ch (Chernoff).
The word \emph{all} means every cell.

\textbf{Shapes.}
A \emph{cell value} is a scalar, a list of $L$ numbers, or an $F \times L$ array.
A scalar is copied to every cell.
A list of $L$ numbers holds one value per component and is copied to every train.
A \emph{selector} has the same three shapes, with strings in place of numbers.
A \emph{profile} holds one value per cell, mission and day; its shapes are given in Section~\ref{sec:input_increments}.
A \emph{flag} is \texttt{true} or \texttt{false}.
No other broadcasting is accepted.

\textbf{Absent entries.}
At a cell where a key does not apply, the entry of an array may be \texttt{null}, and the key may be omitted.
At a cell where a key applies, \texttt{null} is an error, also when the key has a default.
A broadcast value (a scalar, a list, or a profile without the axis $F$) is read only where the key applies; it gives no warning.
A non-\texttt{null} entry of an array that has one entry per cell (an $F \times L$ array, or a profile with the axes $F$ and $L$), at a cell where the key does not apply, is ignored with a \texttt{UserWarning}.

\textbf{Keys.}
The names in the tables are the only keys that an input may hold.
Names are case sensitive.
Any other key raises an error (Section~\ref{sec:input_rejected}).

\subsection{Sizes and Horizon}\label{sec:input_scalars}

\begingroup\small
\begin{longtable}{@{}p{3.2cm}p{2.2cm}p{1.8cm}p{1.9cm}p{4.5cm}@{}}
\toprule
Key & Shape & Default & Applies to & Meaning \\
\midrule
\endhead
\keyrow{\texttt{F}}{integer}{required}{all}{Number of trains. Must exceed $M$}
\keyrow{\texttt{M}}{integer}{required}{all}{Number of missions. Every mission is flown by one train on every day}
\keyrow{\texttt{L}}{integer}{inferred}{all}{Number of components per train. When \texttt{model} is a list or an array, $L$ is the length of its component axis, and a given \texttt{L} must equal it. Required when \texttt{model} is one string and $L > 1$}
\keyrow{\texttt{H}}{integer or \texttt{[H1, H2]}}{required}{all}{Horizon (below). A scalar $H$ means $H_1 = H_2 = H$}
\bottomrule
\end{longtable}
\endgroup

The horizon has $T = H_1 + H_2$ days.
Days $1,\ldots,H_1$ form the \emph{transitory phase}.
Days $H_1+1,\ldots,T$ form the \emph{operating phase}.
Both $H_1$ and $H_2$ are positive integers.
The operating phase is meant to repeat without end; the loop constraint of Section~\ref{sec:loop} makes every repetition safe.
The profiles of Section~\ref{sec:input_increments} have $H_2$ entries for the operating phase and, optionally, $H_1$ entries for the transitory phase.
Index $k = 0$ holds the initial state.

\subsection{Selectors}\label{sec:input_selectors}

\begingroup\small
\begin{longtable}{@{}p{3.2cm}p{2.2cm}p{1.8cm}p{1.9cm}p{4.5cm}@{}}
\toprule
Key & Shape & Default & Applies to & Meaning \\
\midrule
\endhead
\keyrow{\texttt{model}}{selector}{required}{all}{Degradation model of the cell: \texttt{gaussian}, \texttt{inverse\_gaussian}, \texttt{wiener}, \texttt{gamma} or \texttt{rainflow}}
\keyrow{\texttt{repair\_model}}{selector}{\texttt{ardinf}}{all}{Repair law of the cell (Section~\ref{sec:repair_laws}): \texttt{ardinf} (the repair reduces the whole level) or \texttt{ard1} (it reduces the level accumulated since the last intervention). Every model accepts both}
\keyrow{\texttt{bound\_method}}{selector}{\texttt{cantelli}}{RF}{Tail bound of the reliability row (Section~\ref{sec:dyn_rainflow}): \texttt{markov}, \texttt{cantelli}, \texttt{hoeffding}, \texttt{bernstein} or \texttt{chernoff}}
\keyrow{\texttt{reliability\_impl}}{selector}{\texttt{exact}}{G, W, RF (Ca, Be, Ho)}{Form of the reliability row (Section~\ref{sec:linear_forms}): \texttt{exact}, \texttt{tangent} or \texttt{pwl}. A cell whose row is linear (IG, Ga, Mk, Ch) is built in the form \texttt{exact}, and the value is ignored}
\keyrow{\texttt{encoding}}{selector}{\texttt{bigm}}{all}{Encoding of the dynamics rows (Section~\ref{sec:encoding}): \texttt{bigm} or \texttt{indicator}}
\keyrow{\texttt{component\_names}}{list of $L$ strings}{\texttt{Component} $n$ for component $n$}{all}{Name of each component, shown by the plots. A name is not empty after blanks are removed}
\bottomrule
\end{longtable}
\endgroup

\subsection{Reliability}\label{sec:input_reliability}

\begingroup\small
\begin{longtable}{@{}p{3.2cm}p{2.2cm}p{1.8cm}p{1.9cm}p{4.5cm}@{}}
\toprule
Key & Shape & Default & Applies to & Meaning \\
\midrule
\endhead
\keyrow{\texttt{tau}}{cell value}{required}{all}{Degradation threshold $\tau_{i\ell} > 0$}
\keyrow{\texttt{epsilon}}{cell value}{required}{all}{Reliability level $\varepsilon_{i\ell} \in (0,1)$: the probability that the degradation exceeds $\tau_{i\ell}$ must not exceed it on any day. At G and W cells $\varepsilon_{i\ell} < 0.5$}
\keyrow{\texttt{tangent\_ref}}{cell value}{$0.5$}{cells with form \texttt{tangent}}{Relative tangent point in $[0,1)$: the tangent is taken at the mean $\mu_r = \texttt{tangent\_ref}\cdot\tau$}
\keyrow{\texttt{pwl\_points}}{integer}{$8$}{cells with form \texttt{pwl}}{Number of segments of the piecewise linear form}
\bottomrule
\end{longtable}
\endgroup

\emph{Remark (range of $\varepsilon$).}
The Gaussian and Wiener rows contain $\bigl(\Phi^{-1}(1-\varepsilon)\bigr)^2$.
This number is positive only for $\varepsilon < 0.5$.
The variance range $V_{\max} = \tau^2/\bigl(\Phi^{-1}(1-\varepsilon)\bigr)^2$ is the largest variance that the row admits.
It is the natural range of the variance, and the big-M constant of the variance rows (Section~\ref{sec:encoding}).
The ratio $V_{\max}/\tau^2$ is $0.18$ at $\varepsilon = 0.01$, $0.61$ at $0.1$, $3.6$ at $0.3$ and about $1600$ at $0.49$.
The constant is therefore moderate unless $\varepsilon$ is close to $0.5$.
The other models and bounds have no limit below $\varepsilon = 1$: for example, the variance cap $\tau^2\varepsilon/(1-\varepsilon)$ of the Cantelli row is finite for every $\varepsilon < 1$.

\subsection{Repair and Starting Anchors}\label{sec:input_repair}

\begingroup\small
\begin{longtable}{@{}p{3.2cm}p{2.2cm}p{1.8cm}p{1.9cm}p{4.5cm}@{}}
\toprule
Key & Shape & Default & Applies to & Meaning \\
\midrule
\endhead
\keyrow{\texttt{rho}}{cell value}{required}{all}{Repair efficiency $\rho_{i\ell} \in (0,1]$}
\keyrow{\texttt{anchor\_mu\_0}}{cell value}{$\mu_0$}{\texttt{ard1} cells}{Starting anchor $g^{\mu}_{i\ell 0}$ of the mean}
\keyrow{\texttt{anchor\_v\_0}}{cell value}{$v_0$}{\texttt{ard1} cells at G, RF (Ca, Be)}{Starting anchor $g^{v}_{i\ell 0}$ of the variance}
\keyrow{\texttt{anchor\_R\_0}}{cell value}{$R_0$}{\texttt{ard1} cells at RF (Ho)}{Starting anchor $g^{R}_{i\ell 0}$ of the Hoeffding budget}
\keyrow{\texttt{anchor\_K\_0}}{cell value}{$K_0$}{\texttt{ard1} cells at RF (Ch)}{Starting anchor $g^{K}_{i\ell 0}$ of the Chernoff quantity}
\bottomrule
\end{longtable}
\endgroup

Under the law \texttt{ard1}, a repair removes the fraction $\rho$ of the level accumulated since the last intervention (Section~\ref{sec:repair_laws}).
Before the first intervention, the reference level is the starting anchor of the state.
The default of each starting anchor is the initial value of the same state: $\mu_{0}$, $v_{0}$, $R_{0}$ and $K_{0}$ (Section~\ref{sec:input_initial}).
The largest allowed value is this default, because damage only grows between interventions (Section~\ref{sec:consistency_checks}).
A larger anchor gives a weaker repair and a higher state, so the default never overstates what a repair achieves.
For example, let $\mu_0 = 0.5$ and $\rho = 0.5$, and let a repair take place on day~1 before any mission.
With the starting anchor $0$ the mean falls to $0.25$; with the default anchor $0.5$ it stays at $0.5$.
The input is exact when the maintenance history of the component is known.

\subsection{Initial State}\label{sec:input_initial}

\begingroup\small
\begin{longtable}{@{}p{3.2cm}p{2.2cm}p{1.8cm}p{1.9cm}p{4.5cm}@{}}
\toprule
Key & Shape & Default & Applies to & Meaning \\
\midrule
\endhead
\keyrow{\texttt{mu\_0}}{cell value}{required}{all}{Initial mean degradation $\mu_{0,i\ell}$, with $0 \le \mu_{0,i\ell} < \tau_{i\ell}$}
\keyrow{\texttt{v\_0}}{cell value}{$0$}{G, W, RF (Ca, Be)}{Initial variance $v_{0,i\ell} \ge 0$. The value $0$ means a known initial damage}
\keyrow{\texttt{support\_0}}{cell value}{$0$}{RF (Be)}{Width $b^{0}_{i\ell}$ of the range of the initial damage. The value $0$ means a known initial damage}
\keyrow{\texttt{R\_0}}{cell value}{$0$}{RF (Ho)}{Initial value $R_{0,i\ell}$ of the Hoeffding budget: the square of the width of the range of the initial damage}
\keyrow{\texttt{cgf\_0}}{cell value}{$s\,\mu_0$}{RF (Ch)}{Initial value $K_{0,i\ell}$ of the Chernoff quantity: the cumulant generating function of the initial damage at $s_{i\ell} = \texttt{s\_chernoff}$. It is at least $s\,\mu_0$}
\bottomrule
\end{longtable}
\endgroup

The defaults describe a known initial damage: the mean $\mu_0$ with no spread, so that $v_0 = b^0 = R_0 = 0$ and $K_0 = s\,\mu_0$.
The optional keys describe a random initial damage.
Section~\ref{sec:dyn_rainflow} states why these values keep the certificate valid.

\subsection{Replacement}\label{sec:input_replacement}

\begingroup\small
\begin{longtable}{@{}p{3.2cm}p{2.2cm}p{1.8cm}p{1.9cm}p{4.5cm}@{}}
\toprule
Key & Shape & Default & Applies to & Meaning \\
\midrule
\endhead
\keyrow{\texttt{replacement\_mu}}{cell value}{$0$}{all}{Mean degradation $\mu^{\new}_{i\ell}$ after a replacement}
\keyrow{\texttt{replacement\_v}}{cell value}{$0$}{G, W, RF (Ca, Be)}{Variance $v^{\new}_{i\ell}$ after a replacement}
\keyrow{\texttt{support\_new}}{cell value}{$0$}{RF (Be)}{Width $b^{\new}_{i\ell}$ of the range of the damage after a replacement}
\keyrow{\texttt{R\_new}}{cell value}{$0$}{RF (Ho)}{Hoeffding budget $R^{\new}_{i\ell}$ after a replacement}
\keyrow{\texttt{cgf\_new}}{cell value}{$s\,\mu^{\new}$}{RF (Ch)}{Chernoff quantity $K^{\new}_{i\ell}$ after a replacement; at least $s\,\mu^{\new}$}
\keyrow{\texttt{replacement\_as\_new}}{flag}{\texttt{false}}{all}{If \texttt{true}, a replacement gives a component in the zero state, and every replacement input of this table must be $0$}
\keyrow{\texttt{allow\_replacement}}{flag}{\texttt{true}}{all}{If \texttt{false}, no replacement is possible: the variables $x^r$ and their rows are not built}
\bottomrule
\end{longtable}
\endgroup

The state after a replacement comes from these keys for every model.
The defaults give a part in the zero state.
Section~\ref{sec:replacement} states the replacement rows; Section~\ref{sec:dyn_rainflow} states the reset of $R$ and $K$.
For Gamma, the rate of the state after a replacement is \texttt{gamma\_beta\_new} (Section~\ref{sec:input_constants}).

\subsection{Increment Profiles and the Phase Lookup}\label{sec:input_increments}

A \emph{profile} gives the parameters of the damage increment when train $i$ flies mission $j$ on a day, for component $\ell$.
Its canonical shape is $(F, L, M, H_*)$: train, component, mission, day of the phase.
The last axis has $H_2$ entries for an operating profile and $H_1$ entries for a transitory profile.
The parser accepts these shapes and expands each of them to $(F, L, M, H_*)$ before it builds any row (\emph{eager broadcasting}):
\begin{center}
scalar, \quad $(M)$, \quad $(L, M)$, \quad $(L, M, H_*)$, \quad $(F, L, M)$, \quad $(F, L, M, H_*)$.
\end{center}
A missing axis is repeated along that axis.
A three-dimensional array whose shape fits both $(L, M, H_*)$ and $(F, L, M)$ is read as $(L, M, H_*)$; this happens only when $F = L = M = H_*$.
The parser cannot detect an array whose mission axis and component axis are swapped when $M = L$.
Missions have index $j \in \{1,\ldots,M\}$ in a profile; a depot day or an idle day adds no increment.

\textbf{Phase lookup.}
Let $p$ be an operating profile and $p^{\tr}$ the transitory profile of the same key, if it is given.
After the phase lookup, the value used on day $k \in \{1,\ldots,T\}$ is
\begin{equation}\label{eq:c_phase}
p^{\inc}_{i\ell jk} =
\begin{cases}
p^{\tr}_{i\ell jk} & k \le H_1 \text{ and } p^{\tr} \text{ is given},\\[2pt]
p_{i\ell j,\,((k-1) \bmod H_2) + 1} & k \le H_1 \text{ and } p^{\tr} \text{ is not given},\\[2pt]
p_{i\ell j,\,((k-H_1-1) \bmod H_2) + 1} & k > H_1.
\end{cases}
\end{equation}
Every profile key is looked up on its own.
A transitory profile may be given for some keys of a cell and not for others.
For a scalar $H$ without transitory profile, both phases read the same $H$ entries, one after the other.
For $H = [2, 3]$ without transitory profile, days $1,\ldots,5$ read the operating entries $1, 2, 1, 2, 3$.
With a transitory profile, days $1, 2$ read its entries $1, 2$ and days $3, 4, 5$ read the operating entries $1, 2, 3$.
The symbols $\mu^{\inc}_{i\ell jk}$, $v^{\inc}_{i\ell jk}$, $b^{\inc}_{i\ell jk}$ and $K^{\inc}_{i\ell jk}$ of the formulation denote the values after this lookup.

\begingroup\small
\begin{longtable}{@{}p{3.2cm}p{2.2cm}p{1.8cm}p{1.9cm}p{4.5cm}@{}}
\toprule
Key & Shape & Default & Applies to & Meaning \\
\midrule
\endhead
\keyrow{\texttt{mu}}{profile}{required}{all}{Mean $\mu^{\inc}_{i\ell jk} \ge 0$ of the increment. For Ga the increment has shape $\mu\,\beta$ and rate $\beta$, so that its mean is $\mu$}
\keyrow{\texttt{v}}{profile}{required}{G, RF (Ca, Be)}{Variance $v^{\inc}_{i\ell jk} \ge 0$ of the increment}
\keyrow{\texttt{support}}{profile}{required}{RF (Ho, Be)}{Width $b^{\inc}_{i\ell jk} \ge 0$ of the range of the increment: the increment lies in an interval of this length}
\keyrow{\texttt{cgf}}{profile}{required}{RF (Ch)}{Cumulant generating function $K^{\inc}_{i\ell jk} = \ln \mathbb{E}\,e^{s X}$ of the increment $X$ at $s = s_{i\ell}$}
\keyrow{\texttt{s\_chernoff}}{cell value}{required}{RF (Ch)}{Chernoff parameter $s_{i\ell}$, with $s_{i\ell} > \ln(1/\varepsilon_{i\ell})/\tau_{i\ell}$}
\keyrow{\texttt{gamma\_beta}}{scalar, $(L)$, $(F, L)$, $(F, L, M)$ or $(F, L, M, H_2)$}{required}{Ga}{Rate $\beta_{i\ell jk} > 0$ of the Gamma increment. A shape without a day axis is repeated along the days}
\keyrow{\texttt{mu\_trans}, \texttt{v\_trans}, \texttt{support\_trans}, \texttt{cgf\_trans}}{transitory profile}{none}{as the key without the suffix}{Increment parameters of the transitory phase. A transitory profile requires its operating profile}
\keyrow{\texttt{gamma\_beta\_trans}}{scalar, $(L)$, $(F, L)$, $(F, L, M)$ or $(F, L, M, H_1)$}{none}{Ga}{Rates of the transitory phase}
\bottomrule
\end{longtable}
\endgroup

The key \texttt{support} keeps one meaning in every bound: the width of the range of the increment.
A width of $0$ means that the increment is known.

\subsection{Model Constants}\label{sec:input_constants}

\begingroup\small
\begin{longtable}{@{}p{3.2cm}p{2.2cm}p{1.8cm}p{1.9cm}p{4.5cm}@{}}
\toprule
Key & Shape & Default & Applies to & Meaning \\
\midrule
\endhead
\keyrow{\texttt{sigma}}{cell value}{required}{W}{Diffusion constant $\sigma_{i\ell} > 0$: every mission adds $\sigma_{i\ell}^2$ to the variance}
\keyrow{\texttt{eta}}{cell value}{required}{IG}{Ratio $\eta_{i\ell} = \mu/\lambda > 0$ of the mean to the shape parameter of the inverse Gaussian law of the state (Section~\ref{sec:dyn_ig})}
\keyrow{\texttt{v\_max\_user}}{cell value}{none}{G, W, RF (Ca, Be)}{Upper bound $V^{\mathrm{user}}_{\max,i\ell} > 0$ on the variance. If it is given, the row $v_{i\ell k} \le V^{\mathrm{user}}_{\max,i\ell}$ holds on every day. If it is not given, only the reliability row bounds the variance}
\keyrow{\texttt{gamma\_beta\_0}}{cell value}{none (known damage)}{Ga}{Rate of the initial damage, which is a Gamma variable with mean $\mu_0$. If it is omitted, the initial damage is the known constant $\mu_0$, and it does not enter the smallest rate}
\keyrow{\texttt{gamma\_beta\_new}}{cell value}{none (known damage)}{Ga}{Rate of the damage after a replacement, which is a Gamma variable with mean $\mu^{\new}$. If it is omitted, the damage after a replacement is the known constant $\mu^{\new}$, and it does not enter the smallest rate}
\bottomrule
\end{longtable}
\endgroup

For Gamma, the formulation computes the smallest rate of each cell from the rates of the terms with a positive mean: the entries of \texttt{gamma\_beta} and \texttt{gamma\_beta\_trans}, and the rates \texttt{gamma\_beta\_0} and \texttt{gamma\_beta\_new} where they are given (Section~\ref{sec:dyn_gamma}).
No input gives this rate.

\subsection{Costs and Objective}\label{sec:input_costs}

\begingroup\small
\begin{longtable}{@{}p{3.2cm}p{2.2cm}p{1.8cm}p{1.9cm}p{4.5cm}@{}}
\toprule
Key & Shape & Default & Applies to & Meaning \\
\midrule
\endhead
\keyrow{\texttt{C\_M}}{scalar}{required}{all}{Cost $C_M \ge 0$ of one depot day of one train, whatever the number of jobs done on that day}
\keyrow{\texttt{C\_R}}{cell value}{required}{all}{Repair cost $C_{R,i\ell} \ge 0$ per unit of degradation removed. At a cell with $C_{R,i\ell} = 0$ the variable $z$ and its rows are not built}
\keyrow{\texttt{C\_rep}}{cell value}{required}{all}{Replacement cost $C_{\mathrm{rep},i\ell} \ge 0$}
\keyrow{\texttt{C\_D}}{scalar}{required}{all}{Damage regularisation coefficient $C_D \ge 0$. If $C_D = 0$, the variable $u$ and its rows are not built, and \texttt{penalty\_type} is ignored}
\keyrow{\texttt{penalty\_type}}{string}{\texttt{inf\_norm}}{all}{Form of the damage regularisation (Section~\ref{sec:damage_reg}): \texttt{inf\_norm} or \texttt{quadratic}. The value \texttt{quadratic} requires the form \texttt{exact} at every cell}
\bottomrule
\end{longtable}
\endgroup

All costs are finite.
The objective and its three modes (\texttt{objective\_mode}, Section~\ref{sec:run_options}) are stated in Section~\ref{sec:objective}.

\subsection{Run Options}\label{sec:run_options}

Run options are keys of the input.
They do not change the mathematical problem, except \texttt{objective\_mode} and \texttt{evaluation\_horizon}, which choose the objective.

\begingroup\small
\begin{longtable}{@{}p{3.3cm}p{1.9cm}p{1.6cm}p{6.0cm}@{}}
\toprule
Key & Shape & Default & Meaning \\
\midrule
\endhead
\optrow{\texttt{verbose}}{integer}{$1$}{Gurobi output flag: $0$ for no log, $1$ for the log}
\optrow{\texttt{mip\_gap}}{number $\ge 0$}{$0.1$}{Relative gap between the best solution and the best bound at which Gurobi stops (the MIP gap; MIP means mixed-integer program). With the default, the status \texttt{optimal} means that the best solution is at most $10\,\%$ above the best bound. The gap reached is reported as \texttt{mip\_gap} in the result}
\optrow{\texttt{time\_limit}}{number $> 0$}{$3600$}{Time limit of Gurobi, in seconds}
\optrow{\texttt{fast}}{flag}{\texttt{false}}{If \texttt{true}, a preset of Gurobi parameters that favours good solutions found early: \texttt{MIPFocus} $= 1$, \texttt{Heuristics} $= 0.5$, \texttt{ImproveStartGap} $= 0.5$ and \texttt{NoRelHeurTime} $= \max(2,\, 0.15\,\texttt{time\_limit})$}
\optrow{\texttt{gurobi\_params}}{mapping}{none}{Gurobi parameters, by name. They are set after every other option, and so override it}
\optrow{\texttt{objective\_mode}}{string}{\texttt{total}}{\texttt{total}: the cost of all $T$ days; \texttt{operating\_average}: the cost per day of the operating phase; \texttt{evaluation\_total}: the projected cost over \texttt{evaluation\_horizon} days. Each mode adds the damage penalty $C_D\,u$ once (Section~\ref{sec:objective})}
\optrow{\texttt{evaluation\_horizon}}{integer}{none}{Evaluation horizon $E$. Required, and used, only if \texttt{objective\_mode} is \texttt{evaluation\_total}; $E > H_1$. If it is given in another mode, the result reports the projected cost for it (Section~\ref{sec:output})}
\optrow{\texttt{progress\_interval\_seconds}}{number $> 0$}{$5.0$}{Minimum time in seconds between two records of the progress report (Section~\ref{sec:output}). A record is also made whenever the number of solutions found changes}
\optrow{\texttt{relaxation\_warm\_start}}{flag}{\texttt{false}}{If \texttt{true}, and no other warm start is applied, a relaxation start is computed before the solve (Section~\ref{sec:warm_starts})}
\optrow{\texttt{relaxation\_time\_limit}}{number $> 0$}{$60.0$}{Time limit in seconds of the relaxation start}
\bottomrule
\end{longtable}
\endgroup

No input key selects a warm start from a previous result; Section~\ref{sec:warm_starts} defines it.
A horizon sweep is not an input option; the sweep functions of Section~\ref{sec:sweeps} read the same input and solve it for several horizons.

\subsection{Rejected Keys}\label{sec:input_rejected}

Each key of the table below names a quantity that has another name, or another meaning, in this specification.
A key of the table raises an error.
The message names the canonical key and, where the meaning differs, states the meaning of the canonical key.
Any other key that is not in the tables of this section raises the error ``unknown key'', with the list of the canonical keys.
The check on keys runs before any array is read.

\begingroup\small
\begin{longtable}{@{}p{3.3cm}p{11.0cm}@{}}
\toprule
Rejected key & Message: canonical key and meaning \\
\midrule
\endhead
\rejrow{\texttt{maintenance\_type}}{Use \texttt{repair\_model}. Its values are \texttt{ardinf} (the repair reduces the whole level) and \texttt{ard1} (it reduces the level accumulated since the last intervention). In a file written with \texttt{maintenance\_type}, the value \texttt{ARD1} corresponds to \texttt{ardinf}, and the value \texttt{ARA1} to \texttt{ard1}}
\rejrow{\texttt{tail\_bound}, \texttt{method}}{Use \texttt{bound\_method}; its values are \texttt{markov}, \texttt{cantelli}, \texttt{hoeffding}, \texttt{bernstein} and \texttt{chernoff}}
\rejrow{\texttt{formulation}}{Use \texttt{reliability\_impl}. The value \texttt{exact} keeps its meaning. The value \texttt{lp} corresponds to \texttt{tangent} with \texttt{tangent\_ref} equal to $0$}
\rejrow{\texttt{beta}}{Use \texttt{gamma\_beta}. It is a rate, and \texttt{beta} is a scale: \texttt{gamma\_beta} equals $1/\texttt{beta}$}
\rejrow{\texttt{alpha\_inc}}{The Gamma increments are given by the mean \texttt{mu} and the rate \texttt{gamma\_beta}. The shape of an increment is \texttt{mu} times \texttt{gamma\_beta}}
\rejrow{\texttt{gamma\_beta\_bound}}{Remove it. The formulation computes the smallest rate of each cell from the rates of the cell (Section~\ref{sec:dyn_gamma})}
\rejrow{\texttt{gamma\_calibration\_method}}{Remove it. The formulation has no calibration step (Section~\ref{sec:dyn_gamma})}
\rejrow{\texttt{b\_inc}}{Use \texttt{support}: the width of the range of the increment}
\rejrow{\texttt{b\_0}}{Use \texttt{support\_0}: the width of the range of the initial damage}
\rejrow{\texttt{b\_new}}{Use \texttt{support\_new}: the width of the range of the damage after a replacement}
\rejrow{\texttt{n\_workers}, \texttt{warm\_start}}{Remove them. Parallel runs and warm starts between horizons are options of the sweep functions (Sections~\ref{sec:sweeps} and~\ref{sec:warm_starts})}
\rejrow{\texttt{depot\_capacity}}{Remove it. The availability of the depot is not limited at fleet level (Section~\ref{sec:depot})}
\rejrow{\texttt{C\_P}}{Remove it. The loop constraint is a hard constraint and has no penalty (Section~\ref{sec:loop})}
\rejrow{\texttt{C\_S}}{Use \texttt{C\_D}, with the same meaning}
\rejrow{\texttt{alpha}}{Use \texttt{tau}, with the same meaning}
\rejrow{\texttt{xi}, \texttt{repair\_rho}}{Use \texttt{rho}, with the same meaning}
\rejrow{\texttt{mu\_new}}{Use \texttt{replacement\_mu}, with the same meaning}
\rejrow{\texttt{v\_new}}{Use \texttt{replacement\_v}, with the same meaning}
\bottomrule
\end{longtable}
\endgroup

\subsection{Consistency Checks}\label{sec:consistency_checks}

The solver must validate all of the following before it builds the model.
A violation raises a \texttt{ValueError}.
Its message names the key and, for a cell value, the cell $(i,\ell)$.
A check of one model runs only at the cells of that model, on the entries where the key applies.
Each item gives the rule and, after the dash, the reason for it.

\textbf{Keys, shapes and entries.}
\begin{itemize}[leftmargin=1.5em]
  \item \emph{Keys.} Every key is in the tables of Section~\ref{sec:inputs}. The keys \texttt{F}, \texttt{M}, \texttt{H}, \texttt{model}, \texttt{tau}, \texttt{epsilon}, \texttt{rho}, \texttt{mu\_0}, \texttt{mu}, \texttt{C\_M}, \texttt{C\_R}, \texttt{C\_rep} and \texttt{C\_D} are present, and so are the keys that each cell type requires (Section~\ref{sec:schema}). --- A misspelled or missing key must not change the problem without a message.
  \item \emph{Shapes.} Each value has one of the shapes of its table row. The last axis of an operating profile has $H_2$ entries, and that of a transitory profile has $H_1$ entries. \texttt{component\_names} has $L$ entries. --- Silent broadcasting by the array library could read a value for another cell.
  \item \emph{Entries.} At a cell where a key applies, every entry is a finite number, or a string for a selector; \texttt{null} is an error. --- A missing value at a cell must not take a default silently.
\end{itemize}

\textbf{Sizes and selectors.}
\begin{itemize}[leftmargin=1.5em]
  \item $F$, $M$ and $L$ are positive integers, and $F > M$. --- Every mission needs a train on every day; with $F = M$ no train could be in the depot, and no repair could take place.
  \item $H$ is an integer $\ge 1$, or a pair of integers $H_1, H_2 \ge 1$. --- Each phase has at least one day.
  \item \texttt{model} entries are \texttt{gaussian}, \texttt{inverse\_gaussian}, \texttt{wiener}, \texttt{gamma} or \texttt{rainflow}. \texttt{repair\_model} entries are \texttt{ard1} or \texttt{ardinf}, for every model. --- Another value names no model or no repair law.
  \item \texttt{bound\_method} entries at rainflow cells are \texttt{markov}, \texttt{cantelli}, \texttt{hoeffding}, \texttt{bernstein} or \texttt{chernoff}. --- A rainflow cell without one of these bounds has no reliability row.
  \item \texttt{reliability\_impl} entries are \texttt{exact}, \texttt{tangent} or \texttt{pwl} at every cell; at a cell with a linear row (IG, Ga, Mk, Ch) a valid value is ignored, and the form \texttt{exact} is built. \texttt{encoding} entries are \texttt{bigm} or \texttt{indicator}. --- An unknown value must raise an error at every cell; reading it as \texttt{exact} would turn a linear model into a nonconvex one without a message.
  \item \texttt{penalty\_type} is \texttt{inf\_norm} or \texttt{quadratic}. If $C_D > 0$ and \texttt{penalty\_type} is \texttt{quadratic}, every cell is built in the form \texttt{exact}; a cell with a linear row always is. --- The quadratic penalty is a quadratic row, and the linear forms exist to give a linear model.
\end{itemize}

\textbf{Reliability, repair and initial state.}
\begin{itemize}[leftmargin=1.5em]
  \item $\tau_{i\ell} > 0$, and finite. --- The threshold is a positive degradation level.
  \item $\varepsilon_{i\ell} \in (0, 1)$, and $\varepsilon_{i\ell} < 0.5$ at G and W cells. --- A probability lies in $(0,1)$; the squared normal quantile of the Gaussian and Wiener rows is positive only below $0.5$ (Section~\ref{sec:input_reliability}).
  \item $\rho_{i\ell} \in (0, 1]$. --- A repair removes a positive fraction of the level, at most all of it.
  \item $0 \le \mu_{0,i\ell} < \tau_{i\ell}$. --- A component at its threshold has failed already.
  \item $v_0$, \texttt{replacement\_mu}, \texttt{replacement\_v}, \texttt{support\_0}, \texttt{support\_new}, $R_0$ and $R^{\new}$ are $\ge 0$. --- They describe a damage, a variance, a width or a squared width.
  \item \texttt{tangent\_ref} $\in [0, 1)$ at cells with form \texttt{tangent}, and \texttt{pwl\_points} is an integer $\ge 1$. --- At $\texttt{tangent\_ref} = 1$ the tangent row reduces to $Q \le 0$, and a value outside $[0,1]$ has no tangent point in $[0,\tau]$.
  \item Starting anchors: $0 \le \texttt{anchor\_mu\_0} \le \mu_0$, $0 \le \texttt{anchor\_v\_0} \le v_0$, $0 \le \texttt{anchor\_R\_0} \le R_0$ and $0 \le \texttt{anchor\_K\_0} \le K_0$, at the cells that read them. --- Damage only grows between interventions, so the level after the last intervention cannot exceed the initial level.
  \item If \texttt{replacement\_as\_new} is \texttt{true}, then \texttt{replacement\_mu}, \texttt{replacement\_v}, \texttt{support\_new}, \texttt{R\_new} and \texttt{cgf\_new} are $0$ at every cell that reads them. The check does not run when the flag is \texttt{false}. --- The shortcut means a new part, and a non-zero value would be discarded without a message.
\end{itemize}

\textbf{Natural ranges and replacement states.}
\begin{itemize}[leftmargin=1.5em]
  \item Every initial value ($\mu_0$, $v_0$, $R_0$, $K_0$) and every replacement value ($\mu^{\new}$, $v^{\new}$, $R^{\new}$, $K^{\new}$, checked only if \texttt{allow\_replacement} is \texttt{true}) lies in $[0, M_s]$, where $M_s$ is the natural range of its state: $\tau$ for a mean; $V_{\max} = \tau^2/\bigl(\Phi^{-1}(1-\varepsilon)\bigr)^2$ for the variance at G and W cells; $\varepsilon\tau^2/(1-\varepsilon)$ for the variance at Ca cells; $\tau(\tau - 2c)/(2\kappa)$ at Be cells, with $c = \kappa b/3$ and $\kappa = \ln(1/\varepsilon)$; $2\tau^2/\kappa$ for $R$ at Ho cells; $s\tau - \kappa$ for $K$ at Ch cells (Sections~\ref{sec:dyn_gaussian} to~\ref{sec:dyn_rainflow}). A starting anchor lies between $0$ and the initial value of its state, so it lies in the range as well. --- The big-M rows take the natural range as their constant, so a value outside it violates a row at once.
  \item If \texttt{allow\_replacement} is \texttt{true}, the replacement state of every cell satisfies the reliability row of the cell, with $(\mu^{\new}, v^{\new}, R^{\new}, K^{\new})$ in place of the tracked states (Sections~\ref{sec:dyn_gaussian} to~\ref{sec:dyn_rainflow}). --- The row holds on the day of a replacement; otherwise a replacement is never feasible.
\end{itemize}

\textbf{Increments.}
\begin{itemize}[leftmargin=1.5em]
  \item $0 \le \mu^{\inc}_{i\ell jk} < \tau_{i\ell}$ for every mission $j$ and every day $k \in \{1,\ldots,T\}$. --- A mission with mean $\tau$ or more fails the component in one day; the mean $0$ is allowed, so that a mission that does not wear a component can be written as $0$.
  \item Where $\mu^{\inc}_{i\ell jk} = 0$, the values $v^{\inc}_{i\ell jk}$, $b^{\inc}_{i\ell jk}$ and $K^{\inc}_{i\ell jk}$ are $0$. --- A non-negative increment with mean $0$ is $0$ almost surely, so it has no variance, no range and a zero cumulant generating function.
  \item The values $v^{\inc}_{i\ell jk}$, $b^{\inc}_{i\ell jk}$ and $K^{\inc}_{i\ell jk}$ are $\ge 0$. A transitory profile has its operating profile. --- A variance and a width are not negative, and the lookup~\eqref{eq:c_phase} needs the operating entries after the transitory phase.
\end{itemize}

\textbf{Gaussian cells.}
\begin{itemize}[leftmargin=1.5em]
  \item $\mu_0 \ge 3\sqrt{v_0}$ and $\mu^{\inc}_{i\ell jk} \ge 3\sqrt{v^{\inc}_{i\ell jk}}$ for every mission $j$ and day $k$. --- A Gaussian increment is then negative with probability at most about $0.13\,\%$, so the damage stays non-negative in practice.
\end{itemize}

\textbf{Inverse Gaussian cells.}
\begin{itemize}[leftmargin=1.5em]
  \item $\eta_{i\ell} > 0$. --- The dispersion of the inverse Gaussian law must be positive.
\end{itemize}

\textbf{Wiener cells.}
\begin{itemize}[leftmargin=1.5em]
  \item $\sigma_{i\ell} > 0$ and $v_{0,i\ell} \ge 0$. --- The diffusion constant is positive; a known initial damage has variance $0$ and is a valid input.
\end{itemize}

\textbf{Gamma cells.}
\begin{itemize}[leftmargin=1.5em]
  \item Every entry of \texttt{gamma\_beta} and of \texttt{gamma\_beta\_trans} is finite and $> 0$. --- The rate of a Gamma law is positive, and the smallest rate sets the cap (Section~\ref{sec:dyn_gamma}).
  \item \texttt{gamma\_beta\_0} and \texttt{gamma\_beta\_new}, where they are given, are finite and $> 0$. --- A rate is positive; an omitted rate means a known damage, which has no Gamma law and no rate.
\end{itemize}

\textbf{Cells with a variance cap.}
\begin{itemize}[leftmargin=1.5em]
  \item If \texttt{v\_max\_user} is given at a G, W or RF (Ca, Be) cell, then $V^{\mathrm{user}}_{\max,i\ell} > 0$ and $V^{\mathrm{user}}_{\max,i\ell} \ge v_{0,i\ell}$. --- A cap below the initial variance makes the problem infeasible.
\end{itemize}

\textbf{Rainflow cells.}
\begin{itemize}[leftmargin=1.5em]
  \item Bernstein: $v^{\inc}_{i\ell jk} \le (b^{\inc}_{i\ell jk})^2/4$ for every mission $j$ and day $k$, $v_{0} \le (b^{0})^2/4$ and $v^{\new} \le (b^{\new})^2/4$. --- Popoviciu's inequality (below): a law on a range of width $w$ has variance at most $w^2/4$, so a larger variance contradicts the width.
  \item Bernstein: $\tau_{i\ell} > 2\kappa\, b_{i\ell}/3$, with $\kappa = \ln(1/\varepsilon_{i\ell})$ and $b_{i\ell}$ the support constant of Section~\ref{sec:dyn_rainflow}. --- Otherwise the certificate admits no positive variance: it is infeasible at zero damage, and at equality only $\mu = v = 0$ survives, so that the big-M rows break.
  \item Bernstein with the form \texttt{tangent}: $\texttt{tangent\_ref}\cdot\tau_{i\ell} < \tau_{i\ell} - c_{i\ell}$, with $c_{i\ell} = \kappa\, b_{i\ell}/3$. --- Otherwise no point with $v \ge 0$ satisfies the tangent row, so the row is empty (Section~\ref{sec:dyn_rainflow}).
  \item Chernoff: \texttt{s\_chernoff} is given and $s_{i\ell} > \ln(1/\varepsilon_{i\ell})/\tau_{i\ell}$. --- The quantity $K$ is not negative, so the row $K - s\tau \le \ln\varepsilon$ needs $s\tau \ge \ln(1/\varepsilon)$; with a smaller $s$ every schedule is infeasible, and at equality only a component without damage is feasible.
  \item Chernoff: $K^{\inc}_{i\ell jk} \ge s_{i\ell}\,\mu^{\inc}_{i\ell jk}$ for every mission $j$ and day $k$, $K_{0} \ge s\,\mu_{0}$ and $K^{\new} \ge s\,\mu^{\new}$. --- Jensen's inequality: the cumulant generating function of a random variable at $s > 0$ is at least $s$ times its mean.
\end{itemize}

\textbf{Run options.}
\begin{itemize}[leftmargin=1.5em]
  \item \texttt{objective\_mode} is \texttt{total}, \texttt{operating\_average} or \texttt{evaluation\_total}. In the last mode, \texttt{evaluation\_horizon} is an integer $E > H_1$. --- The projected cost needs at least one operating day beyond the transitory phase.
  \item \texttt{mip\_gap} $\ge 0$; \texttt{time\_limit}, \texttt{progress\_interval\_seconds} and \texttt{relaxation\_time\_limit} are $> 0$; \texttt{gurobi\_params} is a mapping; \texttt{fast}, \texttt{relaxation\_warm\_start}, \texttt{allow\_replacement} and \texttt{replacement\_as\_new} are flags. --- A value outside its range must raise an error and must not be changed silently.
\end{itemize}

\emph{Remark (Popoviciu's inequality).}
Let $X$ take values in an interval of width $w$, with midpoint $c$.
Then $\operatorname{Var} X \le \mathbb{E}\,(X - c)^2 \le (w/2)^2$.
The width is the only information on the range, and the range may start at any point $a \ge 0$.
The inequality $\operatorname{Var} X \le w^2/4$ holds for every such position, so it is a necessary condition for the pair (variance, width).
For example, increments in an interval of width $0.2$ have variance at most $0.01$; a two-point law on $\{0, 0.2\}$ with probability $1/2$ each has mean $0.1$ and reaches $0.01$.
A condition of the form $v \le \mu\,(b - \mu)$ assumes that the range starts at $0$, so it is not necessary for a width.

% ================================================================
\section{Decision Variables}
\label{sec:variables}
% ================================================================

A variable, and the rows that define it, are built only where a row or an objective
term with a non-zero coefficient uses them. Example: $u$ and its rows are not built
when $C_D = 0$; $z_{i\ell k}$ and its rows are not built at a cell with
$C_{R,i\ell} = 0$. The day index runs over $k = 1, \ldots, T$; index $k = 0$ holds
constants (initial values and starting anchors), not variables.

\begingroup\small
\begin{longtable}{@{}L{2.3cm}L{2.6cm}L{4.2cm}L{5.4cm}@{}}
\toprule
Variable & Shape, type & Built when & Meaning \\
\midrule
\endhead
$x_{ijk}$ & $F \times (M{+}1) \times T$, binary & always &
  1 if train $i$ flies mission $j$ on day $k$; $j = 0$: depot day \\
$x^m_{i\ell k}$ & $F \times L \times T$, binary & always &
  1 if component $\ell$ of train $i$ is repaired on day $k$ \\
$x^r_{i\ell k}$ & $F \times L \times T$, binary & \texttt{allow\_replacement = true} &
  1 if component $\ell$ of train $i$ is replaced on day $k$ \\
$n_{i\ell k}$ & binary & cells with \texttt{encoding = "indicator"} &
  1 on a day without intervention, Eq.~\eqref{eq:a_ind_sum} \\
$\mu_{i\ell k}$ & continuous, $[0, \tau_{i\ell}]$ & every cell &
  tracked mean degradation \\
$v_{i\ell k}$ & continuous, $[0, M_v]$ & cells whose model or bound tracks a variance &
  tracked variance \\
$R_{i\ell k}$ & continuous, $[0, M_R]$ & rainflow cells with the Hoeffding bound &
  tracked Hoeffding budget \\
$K_{i\ell k}$ & continuous, $[0, M_K]$ & rainflow cells with the Chernoff bound &
  tracked Chernoff quantity \\
$g^{\mu}_{i\ell k}$, $g^{v}_{i\ell k}$, $g^{R}_{i\ell k}$, $g^{K}_{i\ell k}$ &
  continuous, range of the quantity & $\ARDone$ cells, one per tracked quantity
  except the Wiener variance & anchors, Eq.~\eqref{eq:a_anchor} \\
$w_{i\ell kp}$ & binary, $p = 1, \ldots, P$ & cells with a quadratic row in the
  \texttt{pwl} form, $P \geq 2$ & segment of the mean, Eq.~\eqref{eq:a_pwl_sum} \\
$z_{i\ell k}$ & continuous, $\geq 0$ & cells with $C_{R,i\ell} > 0$, on days whose cost
  enters the objective & degradation removed by a repair,
  Section~\ref{sec:repair_cost} \\
$u$ & scalar, $\geq 0$ & $C_D > 0$ & damage regularisation,
  Section~\ref{sec:damage_reg} \\
\bottomrule
\end{longtable}
\endgroup

The ranges $M_v$, $M_R$, $M_K$ are the natural ranges of Section~\ref{sec:encoding};
each model section states which quantities it tracks. The output reports the values
of the tracked quantities; they are upper bounds on the reference values
(Section~\ref{sec:encoding}).

% ================================================================
\section{Objective Function}
\label{sec:objective}
% ================================================================

The cost of day $k$ is
\begin{equation}
  c_k = C_M \sum_{i=1}^{F} x_{i0k}
      + \sum_{i=1}^{F}\sum_{\ell=1}^{L} C_{R,i\ell}\, z_{i\ell k}
      + \sum_{i=1}^{F}\sum_{\ell=1}^{L} C_{\mathrm{rep},i\ell}\, x^r_{i\ell k},
  \label{eq:a_daycost}
\end{equation}
where $z_{i\ell k}$ reads $0$ where it is not built, and the last sum is absent when
\texttt{allow\_replacement = false}. The costs of the two phases are
\begin{equation}
  J_{\mathrm{init}} = \sum_{k=1}^{H_1} c_k, \qquad
  J_{\mathrm{op}} = \sum_{k=H_1+1}^{T} c_k .
  \label{eq:a_phases}
\end{equation}
The input \texttt{objective\_mode} (default \texttt{total}) selects the objective:
\begin{equation}
  \min\;
  \begin{cases}
    J_{\mathrm{init}} + J_{\mathrm{op}} + C_D\, u & \texttt{total}, \\[2pt]
    J_{\mathrm{op}} / H_2 + C_D\, u & \texttt{operating\_average}, \\[2pt]
    J_{\mathrm{init}} + \dfrac{E - H_1}{H_2}\, J_{\mathrm{op}} + C_D\, u
      & \texttt{evaluation\_total},
  \end{cases}
  \label{eq:a_objective}
\end{equation}
where $E$ is the input \texttt{evaluation\_horizon}, required in the mode
\texttt{evaluation\_total}, with $E > H_1$. The term $C_D\,u$ is present only when
$C_D > 0$ (Section~\ref{sec:damage_reg}).

\textbf{Cost terms.}
\begin{itemize}[nosep,leftmargin=1.5em]
  \item $C_M \geq 0$, a scalar, is paid once per train per depot day ($x_{i0k} = 1$),
        whatever the number of components repaired or replaced on that day. It
        prices the visit to the depot, for example the train out of service and the
        set-up. An idle day (Section~\ref{sec:assignment}) costs nothing, so the
        number of depot days is a costed decision.
  \item $C_{R,i\ell} \geq 0$ is paid per unit of mean degradation removed by a
        repair (Section~\ref{sec:repair_cost}).
  \item $C_{\mathrm{rep},i\ell} \geq 0$ is paid per replacement. The input is
        required; it has no default.
  \item $C_D \geq 0$ weights the damage penalty $u$
        (Section~\ref{sec:damage_reg}). The term is added once in every mode,
        because $u$ is the maximum over the whole horizon.
\end{itemize}
The shapes of the cost inputs are given in Section~\ref{sec:inputs}.

\textbf{Meaning of the modes.}
\begin{itemize}[nosep,leftmargin=1.5em]
  \item \texttt{total}: the cost of the $T$ modelled days.
  \item \texttt{operating\_average}: the cost per day of the operating phase. It is
        the long-run cost per day of the endless schedule
        (Section~\ref{sec:loop}), in which the cost of the transitory phase spreads
        over infinitely many days. It compares solves with different $H_2$
        (Section~\ref{sec:sweeps}). The transitory costs do not enter the
        objective in this mode: the solver chooses the transitory schedule for
        feasibility only.
  \item \texttt{evaluation\_total}: the projected cost over $E$ days: the transitory
        phase once, then $(E - H_1)/H_2$ repetitions of the operating phase, a
        fraction of a repetition counted pro rata. It also prices the transitory
        phase, so it compares solves with different $H_1$.
\end{itemize}
Example: $H_1 = 5$, $H_2 = 10$, $E = 365$. Then \texttt{evaluation\_total} weights
the operating cost by $36$, and \texttt{operating\_average} by $1/10$.

The objective value is computed from the tracked values. It is therefore an upper
bound on the cost of the schedule computed from the reference values
(Section~\ref{sec:encoding}).

% ================================================================
\section{Constraints}
\label{sec:constraints}
% ================================================================

The constraints have three groups. The fleet rows (assignment and depot days) do not
depend on the degradation models. The common rules of the dynamics (repair laws,
encoding, replacement, linear forms of the reliability rows) hold for every model;
the five model sections apply them and add each model's tracked quantities, day
maps and reliability row. The repair cost, the damage regularisation and the loop
constraint close the problem. All rows hold for every day $k = 1, \ldots, T$ unless a
section says otherwise.

\subsection{Assignment Constraints}
\label{sec:assignment}

Each train performs at most one activity per day:
\begin{equation}
  \sum_{j=0}^{M} x_{ijk} \leq 1 \qquad \forall\, i, k .
  \label{eq:a_one_activity}
\end{equation}
Each mission is flown by exactly one train every day:
\begin{equation}
  \sum_{i=1}^{F} x_{ijk} = 1 \qquad \forall\, j \in \{1, \ldots, M\},\; k .
  \label{eq:a_coverage}
\end{equation}
Together these rows put exactly $M$ trains on missions every day. The other $F - M$
trains are in the depot or idle. No other fleet-availability row is needed.

\emph{Idle days.} A train is idle on day $k$ when its assignment row is all zeros:
$x_{ijk} = 0$ for every $j \in \{0, 1, \ldots, M\}$. An idle train flies no mission
and is not in the depot. Idle days and depot days add no increment, because every
day aggregate $\delta\mu_{i\ell k}$, $\delta v_{i\ell k}$, $\delta R_{i\ell k}$,
$\delta K_{i\ell k}$ sums over missions $j \geq 1$ only. A depot day must be chosen
explicitly ($x_{i0k} = 1$), and repairs and replacements need one
(Section~\ref{sec:depot}). Because Eq.~\eqref{eq:a_one_activity} is an inequality,
the choice between a depot day and an idle day is free, and the term $C_M \sum
x_{i0k}$ of the objective prices it.

\subsection{Depot-Day Compatibility}
\label{sec:depot}

A component may be repaired or replaced only on a depot day of its train, and at most
one of the two on the same day:
\begin{equation}
  x^m_{i\ell k} + x^r_{i\ell k} \leq x_{i0k} \qquad \forall\, i, \ell, k .
  \label{eq:a_depot}
\end{equation}
With \texttt{allow\_replacement = false} the row reads $x^m_{i\ell k} \leq x_{i0k}$.
Several components of one train may receive interventions on the same depot day. A
depot day without an intervention on a component leaves that component unchanged.

\emph{Remark (one row instead of three).} Eq.~\eqref{eq:a_depot} has the same binary
solutions as the three rows $x^m \leq x_{i0k}$, $x^r \leq x_{i0k}$, $x^m + x^r \leq 1$.
Its linear relaxation is tighter: the point $x_{i0k} = x^m = x^r = 0.5$ satisfies the
three rows but not Eq.~\eqref{eq:a_depot}. With $0 \leq x^m, x^r$ and $x_{i0k} \leq 1$,
Eq.~\eqref{eq:a_depot} describes the convex hull of the four binary solutions
$(x_{i0k}, x^m, x^r) \in \{(0,0,0), (1,0,0), (1,1,0), (1,0,1)\}$.

\subsection{Repair Laws}
\label{sec:repair_laws}

An imperfect repair (in short, a repair) of a component on a depot day removes part
of its degradation. The repair efficiency $\rho_{i\ell} \in (0, 1]$ is the fraction
removed. Let $D^-$ be the degradation level before the repair (at the end of day
$k-1$; a depot day adds no increment) and $D^+$ the level after it (day $k$). The
input \texttt{repair\_model} selects one of two laws per cell; the default is
\texttt{ardinf}.

\begin{definition}[Repair laws]
\label{def:a_laws}
Let $\rho = \rho_{i\ell}$. The two repair laws are:
\begin{itemize}[nosep,leftmargin=1.5em]
  \item $\ARDinf$ (input value \texttt{ardinf}) removes the fraction $\rho$ of the
        whole level:
        \begin{equation}
          D^+ = (1-\rho)\, D^- .
          \label{eq:a_ardinf}
        \end{equation}
  \item $\ARDone$ (input value \texttt{ard1}) removes the fraction $\rho$ of the
        degradation accumulated since the last intervention (repair or
        replacement):
        \begin{equation}
          D^+ = D^- - \rho\,\bigl(D^- - D^{\mathrm{last}}\bigr),
          \label{eq:a_ard1}
        \end{equation}
        where $D^{\mathrm{last}}$ is the level just after the last intervention.
        Before the first intervention of the horizon, $D^{\mathrm{last}}$ is the
        level just after the last intervention before the horizon.
\end{itemize}
\end{definition}

\emph{Names and sources.} $\ARDone$, ``arithmetic reduction of deterioration of
order 1'', is the law of \cite[$\S$3.1, eq.~(4)]{mercier2019}. That source measures the
removed degradation from the level just after the last maintenance, as
Eq.~\eqref{eq:a_ard1} does. The name $\ARDinf$ is an analogy, not a citation.
Doyen and Gaudoin define the classes ARA$_\infty$ \cite[eq.~(20), p.~11]{doyen2004}
and ARA$_1$ \cite[eq.~(22), p.~12]{doyen2004}, which reduce a virtual age, not a
degradation level. ARA$_\infty$ reduces the whole virtual age; ARA$_1$ reduces only
the age accumulated since the previous repair. The subscript is the memory of the
class: the largest number of previous failure times that enter the failure
intensity \cite[p.~8]{doyen2004}. $\ARDinf$ applies the rule of ARA$_\infty$ to the
whole degradation level. No law of this specification acts on a virtual age.

\textbf{Repair maps of the tracked quantities.} A model tracks moments or other
summaries of $D$, not $D$ itself. For a tracked quantity $s$ with anchor $g$, a
repair on day $k$ applies the repair map
\begin{equation}
  \psi(s_{k-1}, g_{k-1}) =
  \begin{cases}
    a\, s_{k-1} & \ARDinf, \\
    a\, s_{k-1} + (1-a)\, g_{k-1} & \ARDone,
  \end{cases}
  \label{eq:a_psi}
\end{equation}
with a repair factor $a \in [0, 1]$. For the mean, $a = 1 - \rho_{i\ell}$, because
expectation is linear: Eq.~\eqref{eq:a_ard1} gives $\mathbb{E}[D^+] =
(1-\rho)\,\mathbb{E}[D^-] + \rho\,\mathbb{E}[D^{\mathrm{last}}]$. For a variance and
for $R$ the factor is $(1-\rho_{i\ell})^2$, and for $K$ it is $1-\rho_{i\ell}$; the
model sections derive each factor. The Wiener variance has the factor $a = 1$
under both laws: a repair acts on the Wiener mean only and leaves the variance
unchanged.

\textbf{Anchors.} Under $\ARDone$ each tracked quantity has an anchor $g$: the value
of the quantity just after the last intervention. The Wiener variance is the
exception: with $a = 1$ its map does not use an anchor, and none is built. The
anchor recursion is
\begin{equation}
  g_0 = \text{starting anchor}, \qquad
  g_k = \begin{cases}
    s_k & \text{on a repair or replacement day}, \\
    g_{k-1} & \text{otherwise},
  \end{cases}
  \label{eq:a_anchor}
\end{equation}
built by the rows Eq.~\eqref{eq:a_bigm_anchor} or Eq.~\eqref{eq:a_ind_anchor}. The
starting anchors $g^{\mu}_{i\ell 0}$, $g^{v}_{i\ell 0}$, $g^{R}_{i\ell 0}$,
$g^{K}_{i\ell 0}$ come from the optional inputs \texttt{anchor\_mu\_0},
\texttt{anchor\_v\_0}, \texttt{anchor\_R\_0}, \texttt{anchor\_K\_0}. The default of
each is the initial value of its quantity. Each must satisfy $0 \leq g_0 \leq s_0$
(Section~\ref{sec:consistency_checks}). Example: $\mu_{k-1} = 0.6$, $g_{k-1} = 0.2$,
$\rho = 0.5$. A repair gives $\mu_k = 0.6 - 0.5\,(0.6 - 0.2) = 0.4$ and $g_k = 0.4$.
On the next day the gap $\mu - g$ is $0$, and it grows again with each mission.

\emph{Remark (default of the starting anchor).} Damage only grows between
interventions. So the true level after the last intervention before the horizon lies
between the level of a new part and the initial level. A larger anchor gives a
weaker repair and a higher state, because $\psi$ is non-decreasing in $g$. The
default $g_0 = s_0$ therefore never overstates what a repair achieves; the input
gives the exact value when the maintenance history is known. Example: $\mu_0 =
0.5$, $\rho = 0.5$, and a repair on day~1 before any mission. With the starting
anchor $0$ the mean drops to $0.25$; with the default $\mu_0$ it stays at $0.5$, and
the repair removes only damage accumulated within the horizon.

\textbf{Laws per model.} Every model accepts both laws. Table~\ref{tab:a_laws} gives
the reason for each model.

\begingroup\small
\begin{longtable}{@{}lllp{7.4cm}@{}}
\caption{Repair laws accepted by each model.}
\label{tab:a_laws} \\
\toprule
Model & $\ARDinf$ & $\ARDone$ & Reason or note \\
\midrule
Gaussian & yes & yes & both laws keep the state Gaussian, with exact moments \\
Inverse Gaussian & yes & yes & the cap holds for any repair weights and any
  dependence between the terms (Section~\ref{sec:dyn_ig}) \\
Wiener & yes & yes & on the mean only; the variance changes only by replacement \\
Gamma & yes & yes & Section~\ref{sec:dyn_gamma} \\
Rainflow, every bound & yes & yes & Chernoff under $\ARDone$ uses a conservative
  recursion for $K$ (Section~\ref{sec:dyn_rainflow}) \\
\bottomrule
\end{longtable}
\endgroup

\subsection{Encoding of the Dynamics}
\label{sec:encoding}

Each day, every tracked quantity of a cell follows one rule: a mission adds an
increment, a repair applies the repair map of Section~\ref{sec:repair_laws}, and a
replacement resets the quantity (Section~\ref{sec:replacement}). This subsection
states these rules once, for every model and both repair laws. The model sections
give, for each tracked quantity, its increment, its repair factor and its natural
range.

\textbf{One-sided rows.} Every dynamics row is an inequality ($\geq$), not an
equality. The solver may choose a value above the result of the rule, never a
value below it. So a tracked quantity is an upper bound on the value that the rule
gives. Example: with $\mu_{i\ell,k-1} = 0.3$ and a mission increment of $0.2$, the
rule gives $0.5$; the rows accept $\mu_{i\ell k} = 0.5$ and also $\mu_{i\ell k} =
0.6$. Replacement rows and anchor updates are equalities.

\textbf{Generic rows.} Let $s$ be one tracked quantity of a cell (a mean, a
variance, $R$ or $K$), with day increment $\delta s_k$, natural range $M_s$,
replacement value $s^{\new}$, anchor $g$ (law $\ARDone$ only) and repair map
$\psi$ (Section~\ref{sec:repair_laws}). Write $x^m_k$, $x^r_k$ for the binaries of
the cell and $e_k = x^m_k + x^r_k$. The input \texttt{encoding} selects one of two
forms per cell. Both forms are built for $k = 1, \ldots, T$; index $k = 0$ holds the
initial value and the starting anchor.

\emph{Big-M form} (\texttt{encoding = "bigm"}, the default):
\begin{align}
  &\text{no intervention:} & s_k &\geq s_{k-1} + \delta s_k - M_s\,(x^m_k + x^r_k),
    \label{eq:a_bigm_carry} \\
  &\text{repair:} & s_k &\geq \psi(s_{k-1}, g_{k-1}) - M_s\,(1 - x^m_k),
    \label{eq:a_bigm_repair} \\
  &\text{replacement:} & |s_k - s^{\new}| &\leq M_s\,(1 - x^r_k),
    \label{eq:a_bigm_repl} \\
  &\text{anchor } (\ARDone)\text{:} & |g_k - s_k| &\leq M_s\,(1 - e_k), \qquad
    |g_k - g_{k-1}| \leq M_s\, e_k .
    \label{eq:a_bigm_anchor}
\end{align}
Each row with an absolute value is built as two linear rows.

\emph{Indicator form} (\texttt{encoding = "indicator"}). One extra binary
$n_{i\ell k}$ per cell and day marks a day without intervention:
\begin{equation}
  n_{i\ell k} + x^m_{i\ell k} + x^r_{i\ell k} = 1 .
  \label{eq:a_ind_sum}
\end{equation}
The rows are Gurobi indicator constraints:
\begin{align}
  n_k = 1 &\;\Rightarrow\; s_k \geq s_{k-1} + \delta s_k, &
  x^m_k = 1 &\;\Rightarrow\; s_k \geq \psi(s_{k-1}, g_{k-1}), &
  x^r_k = 1 &\;\Rightarrow\; s_k = s^{\new}, \label{eq:a_ind_state} \\
  n_k = 1 &\;\Rightarrow\; g_k = g_{k-1}, &
  x^m_k = 1 &\;\Rightarrow\; g_k = s_k, &
  x^r_k = 1 &\;\Rightarrow\; g_k = s_k. \label{eq:a_ind_anchor}
\end{align}
The binary $n_{i\ell k}$ is built only at cells in this form.

\emph{Common rules for both forms.}
\begin{itemize}[nosep,leftmargin=1.5em]
  \item The anchor rows are built only under $\ARDone$, one anchor per tracked
        quantity, except the Wiener variance.
  \item The Wiener variance takes the repair map $\psi = s_{k-1}$ (factor $a = 1$)
        under both laws, and has no anchor. A repair day adds no increment, so in the
        big-M form its no-intervention and repair rows together equal the single row
        $s_k \geq s_{k-1} + \delta s_k - M_s\, x^r_k$ (Eq.~\eqref{eq:w_vrow}), which
        may replace them. The variance is never left free below on a repair day.
  \item With \texttt{allow\_replacement = false} the binaries $x^r$ are not built:
        read $x^r_k = 0$ in every row, and drop the replacement rows.
  \item At binary points the two forms have the same feasible set. Every statement
        below holds for both.
\end{itemize}

\textbf{Natural ranges.} The big-M constant $M_s$ is the natural range of $s$: an
upper bound of $s$ at every point that satisfies the reliability row of the cell.
It is $\tau_{i\ell}$ for a mean and for its anchor, $V_{\max}$ for a Gaussian or
Wiener variance, $V^{\mathrm{C}}$ or $V^{\mathrm{B}}$ for a rainflow variance
(Cantelli or Bernstein), $2\tau^2/\kappa$ for $R$ and $s_{i\ell}\tau - \kappa$ for
$K$ (Section~\ref{sec:dyn_rainflow}). An anchor has the range of its quantity.

\emph{Remark (validity of the big-M constants).} Every tracked quantity and every
anchor lies in $[0, M_s]$. The initial values, the starting anchors and the
replacement values lie in $[0, M_s]$ too, by the input checks of
Section~\ref{sec:consistency_checks}. An intervention day is a depot day
(Section~\ref{sec:depot}), so $\delta s_k = 0$ on it. A relaxed row then cuts no
point: Eq.~\eqref{eq:a_bigm_carry} with $e_k = 1$ reads
$s_k \geq s_{k-1} - M_s$, whose right-hand side is at most $0$;
Eq.~\eqref{eq:a_bigm_repair} with $x^m_k = 0$ has a right-hand side
$\psi - M_s \leq 0$, because $\psi$ is a convex combination of values in
$[0, M_s]$; Eq.~\eqref{eq:a_bigm_repl} with $x^r_k = 0$ allows
$s_k \in [s^{\new} - M_s,\, s^{\new} + M_s] \supseteq [0, M_s]$; the relaxed anchor
rows allow every difference in $[-M_s, M_s]$.

\emph{Remark (choice of the default form).} Gurobi accepts a binary as integral
within its integrality tolerance \texttt{IntFeasTol} (default $10^{-5}$). A big-M
row can therefore leak by at most $10^{-5} M_s$. Because $M_s$ is the quantity's own
range, this is a relative error of $10^{-5}$. Indicator constraints have no such
leak, and Gurobi usually needs more time to solve them. Which form solves this
model faster has not been measured.

\subsubsection{Safety of the one-sided rows}

Fix a schedule, that is, values of every binary, and one cell. The following
objects are used here and in Section~\ref{sec:loop}.
\begin{itemize}[nosep,leftmargin=1.5em]
  \item The \emph{full tracked state} $\xi_k$ collects the values on day $k$ of every
        tracked quantity of the cell and of every anchor. $\xi_0$ holds the initial
        values and the starting anchors.
  \item The \emph{day map} $\varphi_k$ gives, for each coordinate, the value that the
        rule of day $k$ assigns: $s + \delta s_k$ and $g$ on a day without
        intervention; $\psi(s, g)$ for both $s$ and $g$ on a repair day; $s^{\new}$
        for both on a replacement day.
  \item The \emph{reference trajectory} is $\zeta_0 = \xi_0$,
        $\zeta_k = \varphi_k(\zeta_{k-1})$: the rows with every inequality replaced
        by an equality.
  \item The \emph{true statistics} $\theta_k$ are the true values on day $k$ of the
        tracked quantities (the mean, the variance, $R$ or the CGF of the damage
        $D_{i\ell k}$) and of their anchors (their values just after the last
        intervention).
  \item The \emph{certificate set} $\mathcal{C}_{i\ell}$ is the set of values of the
        tracked quantities at which the reliability row of the cell holds in its
        exact form, side rows included. Anchors do not enter it.
\end{itemize}

Every model section states three facts, labelled (M1)--(M3) there:
\begin{enumerate}[nosep,leftmargin=2em,label=(M\arabic*)]
  \item Each repair factor $a$ lies in $[0,1]$ and each increment
        $\delta s_k$ is non-negative.
  \item On day $0$ the true statistics are at most the initial values and starting
        anchors, $\theta_0 \leq \xi_0$. On every later day the true statistic is at
        most the day map applied to the true statistics of the day before:
        $\theta_k \leq \varphi_k(\theta_{k-1})$. Every map holds with equality
        except the map of the Chernoff quantity $K$, which is an upper bound
        (Section~\ref{sec:dyn_rainflow}).
  \item If values $s \in \mathcal{C}_{i\ell}$ are at least the true statistics of
        day $k$, then $\Pr(D_{i\ell k} > \tau_{i\ell}) \leq \varepsilon_{i\ell}$.
\end{enumerate}
The set $\mathcal{C}_{i\ell}$ is closed, and it is a down-set: if $s \in
\mathcal{C}_{i\ell}$ and $0 \leq s' \leq s$ componentwise, then $s' \in
\mathcal{C}_{i\ell}$. The reason is that every reliability row, written in
chance-bound form, reads ``a continuous function of the tracked quantities,
non-decreasing in each of them, is at most a constant'':

\begingroup\small
\begin{longtable}{@{}lll@{}}
\toprule
Model or bound & Row in chance-bound form & Section \\
\midrule
Gaussian, Wiener & $\mu + \Phi^{-1}(1-\varepsilon)\sqrt{v} \leq \tau$ &
  \ref{sec:dyn_gaussian}, \ref{sec:dyn_wiener} \\
Inverse Gaussian; Gamma & $\mu \leq \bar\mu$; $\mu \leq A_1(\beta^*\tau,\varepsilon)/\beta^*$ &
  \ref{sec:dyn_ig}, \ref{sec:dyn_gamma} \\
Rainflow, Markov & $\mu \leq \varepsilon\,\tau$ & \ref{sec:dyn_rainflow} \\
Rainflow, Cantelli & $\mu + \kappa_C\sqrt{v} \leq \tau$ & \ref{sec:dyn_rainflow} \\
Rainflow, Hoeffding & $\mu + \sqrt{\kappa R/2} \leq \tau$ & \ref{sec:dyn_rainflow} \\
Rainflow, Bernstein & $\mu + c + \sqrt{c^2 + 2\kappa v} \leq \tau$ & \ref{sec:dyn_rainflow} \\
Rainflow, Chernoff & $K \leq s_{i\ell}\,\tau + \ln\varepsilon$ & \ref{sec:dyn_rainflow} \\
\bottomrule
\end{longtable}
\endgroup

\begin{lemma}[Monotone day maps]
\label{lem:a_monotone}
Each day map $\varphi_k$ is non-decreasing in each coordinate of the full tracked
state, anchors included.
\end{lemma}
\begin{proof}
Let $\xi \leq \xi'$ componentwise. Each coordinate of $\varphi_k$ has one of the
forms $s + \delta s_k$, $g$, $s$, $a\,s$, $a\,s + (1-a)\,g$ or $s^{\new}$. By (M1),
$a \geq 0$ and $1 - a \geq 0$. Therefore each coordinate of $\varphi_k(\xi)$ is at
most the same coordinate of $\varphi_k(\xi')$.
\end{proof}

\begin{proposition}[Domination]
\label{prop:a_domination}
Every feasible solution satisfies $\xi_k \geq \varphi_k(\xi_{k-1})$ and
$\theta_k \leq \zeta_k \leq \xi_k$ for $k = 1, \ldots, T$.
\end{proposition}
\begin{proof}
At binary points the rows of day $k$ read: $s_k \geq \varphi_k(\xi_{k-1})_s$ for
every tracked quantity $s$, with equality on a replacement day; $g_k = s_k \geq
\varphi_k(\xi_{k-1})_g$ on an intervention day; $g_k = g_{k-1} =
\varphi_k(\xi_{k-1})_g$ on any other day. So $\xi_k \geq \varphi_k(\xi_{k-1})$.
Let $\theta_{k-1} \leq \zeta_{k-1} \leq \xi_{k-1}$ (true for $k = 1$ by (M2), since
$\zeta_0 = \xi_0$). Then, by (M2) and Lemma~\ref{lem:a_monotone},
\[
  \theta_k \leq \varphi_k(\theta_{k-1}) \leq \varphi_k(\zeta_{k-1}) = \zeta_k
  \leq \varphi_k(\xi_{k-1}) \leq \xi_k .
\]
Therefore $\theta_k \leq \zeta_k \leq \xi_k$ for every $k$ by induction.
\end{proof}

\emph{Consequence on the horizon.} On every day $k = 1, \ldots, T$ the tracked
values satisfy the reliability row in the form chosen for the cell. Every form lies
inside $\mathcal{C}_{i\ell}$ (Section~\ref{sec:linear_forms}), so $\xi_k \in
\mathcal{C}_{i\ell}$. By Proposition~\ref{prop:a_domination}, $\theta_k \leq \xi_k$,
so (M3) gives $\Pr(D_{i\ell k} > \tau_{i\ell}) \leq \varepsilon_{i\ell}$.
Theorem~\ref{thm:a_loop} (Section~\ref{sec:loop}) extends this to every day of
every later repetition of the operating phase.

\textbf{What the one-sided rows accept.} Call a schedule \emph{safe} for a cell when
the reference trajectory of the endless schedule (the transitory phase once, then
the operating phase repeated without end, Section~\ref{sec:loop}) lies in
$\mathcal{C}_{i\ell}$ and in the natural ranges on every day. The one-sided rows
relate to this property as follows.
\begin{itemize}[nosep,leftmargin=1.5em]
  \item \emph{They never accept an unsafe schedule} (Theorem~\ref{thm:a_loop}).
  \item \emph{They accept every schedule that equality rows accept}, and more.
        Equality rows force $\xi_k = \zeta_k$; Proposition~\ref{prop:a_accept}~(a)
        shows that such a schedule satisfies the one-sided rows too.
        The first example below shows a safe schedule that only the one-sided
        rows accept.
  \item \emph{They may reject a safe schedule}, as equality rows may, but only when
        tracked quantities move in opposite directions over the first repetition
        (Proposition~\ref{prop:a_accept}~(b) and the second example below).
\end{itemize}

\begin{proposition}[Acceptance]
\label{prop:a_accept}
Fix a schedule and a cell, use the exact reliability form, and let the schedule be
safe for the cell. Let $\bar\zeta = \zeta_{H_1}$ and let $\Phi$ be the composition
of the day maps of the operating phase (Section~\ref{sec:loop}).
\begin{enumerate}[nosep,label=(\alph*)]
  \item If $\Phi(\bar\zeta) \leq \bar\zeta$, the rows of the cell, the loop row
        included, accept the schedule with $\xi_k = \zeta_k$ for every $k$.
  \item If $\Phi(\bar\zeta) \geq \bar\zeta$, the cell has no anchor (law
        $\ARDinf$), and the cell is not replaced on day $H_1$, the rows of the cell
        accept the schedule.
\end{enumerate}
\end{proposition}
\begin{proof}
(a) The values $\xi_k = \zeta_k$ satisfy every row of day $k$ with equality at
binary points. They lie in $\mathcal{C}_{i\ell}$ because the schedule is safe. The
loop row reads $\zeta_T = \Phi(\bar\zeta) \leq \bar\zeta = \zeta_{H_1}$.

(b) Let $y_n = \Phi^n(\bar\zeta)$, the reference state at the end of the $n$-th
repetition. By Lemma~\ref{lem:a_monotone} and $\Phi(\bar\zeta) \geq \bar\zeta$,
$y_{n+1} = \Phi(y_n) \geq \Phi(y_{n-1}) = y_n$. The sequence is non-decreasing and
lies in the natural ranges, so it converges to some $y^*$ with $y^* \geq
\bar\zeta$. Each day map is affine or constant, so $\Phi$ and its prefixes
$\Phi_h$ are continuous: $\Phi(y^*) = \lim_n y_{n+1} = y^*$ and
$\Phi_h(y^*) = \lim_n \Phi_h(y_n)$. Each $\Phi_h(y_n)$ is a reference value of a
later day, so it lies in the closed set $\mathcal{C}_{i\ell}$; hence
$\Phi_h(y^*) \in \mathcal{C}_{i\ell}$ for $h = 1, \ldots, H_2$, and
$y^* = \Phi_{H_2}(y^*) \in \mathcal{C}_{i\ell}$. Take $\xi_k =
\zeta_k$ for $k < H_1$, $\xi_{H_1} = y^*$ and $\xi_{H_1+h} = \Phi_h(y^*)$. Day $H_1$
is not a replacement day and the cell has no anchor, so $\xi_{H_1} = y^* \geq
\varphi_{H_1}(\xi_{H_1-1})$ satisfies the rows of day $H_1$. The later days satisfy
their rows with equality, and the loop row reads $\xi_T = \Phi(y^*) = y^* =
\xi_{H_1}$.
\end{proof}

A scalar state is always comparable with its image. So a cell with one tracked
quantity under $\ARDinf$, not replaced on day $H_1$, accepts every safe schedule.
Under $\ARDone$ the anchor is a second tracked quantity, and its value on day $H_1$
is fixed by the anchor rows; part (b) does not apply, and acceptance of every safe
schedule is not proven.

\emph{Example (one quantity, accepted only by the one-sided rows).}
One tracked mean, law $\ARDinf$, $\rho = 0.5$. The operating phase is a mission
($\delta\mu = 0.2$) followed by a repair, so $\Phi(y) = 0.5\,(y + 0.2) = 0.5\,y +
0.1$, and $\Phi(y) \leq y$ if and only if $y \geq 0.2$. Let $\bar\zeta = 0.1$. Then
$\Phi(0.1) = 0.15 > 0.1$, so equality rows reject the schedule. The one-sided rows
accept it with $\xi_{H_1} = 0.2$: the tracked cycle is $0.2, 0.4, 0.2$, and the
reliability row is checked at $0.4$. The reference trajectory is $0.1, 0.3, 0.15,
0.35, 0.175, \ldots$; its start-of-cycle values rise to $0.2$ and its peaks rise to
$0.4$. So the schedule is safe whenever $0.4 \in \mathcal{C}_{i\ell}$. The price
is in the reported cost: each repair is charged $\rho \cdot 0.4 = 0.2$, while the
reference removal rises from $0.15$ toward $0.2$.

\emph{Example (two quantities that move in opposite directions, rejected).}
Cantelli bound, $\varepsilon = 0.1$ ($\kappa_C^2 = 9$), $\tau = 1$, $\rho = 0.5$,
law $\ARDinf$. The operating phase is a mission, which adds $0.2$ to the mean and
$0.009$ to the variance, followed by a repair. So $\Phi(\mu, v) = (0.5\,\mu + 0.1,\;
0.25\,v + 0.00225)$, with fixed point $(0.2, 0.003)$. Let $\bar\zeta = (0.5, 0)$.
\begin{itemize}[nosep,leftmargin=1.5em]
  \item Reference: over the repetitions the mean falls to $0.2$ and the variance
        rises to $0.003$. The first peak is $(0.7, 0.009)$, with $9 \cdot 0.009 =
        0.081 \leq 0.3^2 = 0.09$. Every later peak has a mean of at most $0.55$ and a
        variance below $0.012$, with $9 \cdot 0.012 = 0.108 \leq 0.45^2 = 0.2025$.
        The schedule is safe.
  \item One-sided rows: the loop row needs $\xi_{H_1} = y \geq (0.5, 0)$ with
        $\Phi(y) \leq y$, so $y_v \geq 0.003$. The tracked peak is then at least
        $(0.7, 0.012)$, where $9 \cdot 0.012 = 0.108 > 0.09$. Rejected.
  \item Equality rows: $\Phi(\bar\zeta) = (0.35, 0.00225)$ has a larger variance
        than $\bar\zeta$. Rejected as well.
\end{itemize}

\textbf{Meaning of the reported values.} The output reports the tracked values
$\xi_k$, not the reference values. By Proposition~\ref{prop:a_domination} each
reported state and anchor is an upper bound on its reference value, and so on its
true statistic. The costs
computed from them are upper bounds too: the repair cost
(Proposition~\ref{prop:a_repair_cost}) and the damage penalty
(Section~\ref{sec:damage_reg}).

\emph{Remark (slack in the reported values).} The minimisation removes slack only
where slack costs something: a larger mean raises the damage penalty $C_D\,u$ where
$u$ is binding, and raises the repair cost $C_R\,z$ of every later repair. A larger
variance, $R$ or $K$ tightens the reliability row and so restricts the schedule.
Elsewhere the solver may return any value above the reference value. With $C_D = 0$
the damage penalty is absent (Section~\ref{sec:damage_reg}), and the reported states
may carry more slack. Slack never affects safety (Theorem~\ref{thm:a_loop}).

\subsection{Replacement}
\label{sec:replacement}

A replacement installs a new part on a depot day. The input
\texttt{allow\_replacement} (default \texttt{true}) makes replacements available. When
it is \texttt{false}, the binaries $x^r_{i\ell k}$, the replacement rows and the
replacement cost are not built.

\textbf{State after a replacement.} On a replacement day every tracked quantity of
the cell takes its replacement value, and every anchor takes the same value
(Eq.~\eqref{eq:a_anchor}). The values come from the input, for every model:
\begin{itemize}[nosep,leftmargin=1.5em]
  \item mean $\mu^{\new}_{i\ell}$, input \texttt{replacement\_mu}, default $0$;
  \item variance $v^{\new}_{i\ell}$, input \texttt{replacement\_v}, default $0$, where
        the model tracks a variance;
  \item $R^{\new}_{i\ell}$ and $K^{\new}_{i\ell}$ for the rainflow Hoeffding and
        Chernoff bounds, with their defaults and their derivation in
        Section~\ref{sec:dyn_rainflow}.
\end{itemize}
A replacement value describes the damage of the new part: a known damage when its
variance is $0$, a random one otherwise. The rows are the equalities
Eq.~\eqref{eq:a_bigm_repl} or the indicator rows Eq.~\eqref{eq:a_ind_state}. The
input checks of Section~\ref{sec:consistency_checks} require each replacement value
to lie in the natural range of its quantity, and the replacement state to satisfy the
reliability row of the cell; otherwise no replacement of that cell would be
feasible.

\textbf{Shortcut for a new part.} The input \texttt{replacement\_as\_new} (default
\texttt{false}) states that a replaced part is new, with no damage: every
replacement value of the cell is $0$. A non-zero replacement input together with
\texttt{replacement\_as\_new = true} is an input error
(Section~\ref{sec:consistency_checks}).

\textbf{Cost.} A replacement costs $C_{\mathrm{rep},i\ell}$ and nothing else: the
repair cost $z_{i\ell k}$ is $0$ on a replacement day
(Section~\ref{sec:repair_cost}). A cell may be replaced on any depot day of its
train, any number of times. The loop row can require at least one replacement in
the operating phase: for every $\ARDone$ cell with $\rho < 1$ that accumulates damage
(Proposition~\ref{prop:a_ard1_replacement}), and for the Wiener variance
(Section~\ref{sec:dyn_wiener}).

\subsection{Linear Forms of the Reliability Rows}
\label{sec:linear_forms}

Five reliability rows are quadratic: the Gaussian and Wiener rows, and the rainflow
Cantelli, Hoeffding and Bernstein rows. Each has the form
\begin{equation}
  \chi_{i\ell k} \;\leq\; f(\mu_{i\ell k}),
  \qquad f(\mu) = c_2\,(\tau - \mu)^2 + c_1\,(\tau - \mu),
  \qquad \mu_{i\ell k} \leq \bar\mu_{\mathrm{side}},
  \label{eq:a_quad}
\end{equation}
with the constants of Table~\ref{tab:a_quad}. The third relation is the side row; it
selects the branch of the parabola on which the row is a valid certificate. Since
$c_2 > 0$, $f$ is convex in $\mu$, with $f'(\mu) = -2c_2(\tau - \mu) - c_1$.

\begingroup\small
\begin{longtable}{@{}lllll@{}}
\caption{Constants of the quadratic reliability rows, Eq.~\eqref{eq:a_quad}. The
model sections derive each row.}
\label{tab:a_quad} \\
\toprule
Row & $\chi$ & $c_2$ & $c_1$ & $\bar\mu_{\mathrm{side}}$ \\
\midrule
Gaussian, Wiener & $v$ & $1/(\Phi^{-1}(1-\varepsilon))^2$ & $0$ & $\tau$ \\
Rainflow, Cantelli & $v$ & $1/\kappa_C^2 = \varepsilon/(1-\varepsilon)$ & $0$ & $\tau$ \\
Rainflow, Hoeffding & $R$ & $2/\kappa$ & $0$ & $\tau$ \\
Rainflow, Bernstein & $v$ & $1/(2\kappa)$ & $-b_{i\ell}/3$ & $\tau - c_{i\ell}$ \\
\bottomrule
\end{longtable}
\endgroup

The input \texttt{reliability\_impl} selects one of three forms per cell. Every form
keeps the side row.
\begin{itemize}[leftmargin=1.5em]
  \item \texttt{exact}: Eq.~\eqref{eq:a_quad} as written. The feasible set
        $\{(\mu, \chi) : \chi \leq f(\mu)\}$ lies below a convex function, so it is not
        convex: the row is a nonconvex quadratic constraint. Gurobi's parameter
        \texttt{NonConvex} is set to $2$ whenever at least one cell has a quadratic
        row in this form.
  \item \texttt{tangent}: one linear row per day, the tangent of $f$ at the point
        $\mu_r = \texttt{tangent\_ref}_{i\ell}\cdot\tau_{i\ell}$:
        \begin{equation}
          \chi_{i\ell k} \;\leq\; f(\mu_r) + f'(\mu_r)\,(\mu_{i\ell k} - \mu_r).
          \label{eq:a_tangent}
        \end{equation}
        The input \texttt{tangent\_ref} is per cell, in $[0, 1)$, default $0.5$.
        At \texttt{tangent\_ref} $= 0$ the Gaussian row becomes
        $(\Phi^{-1}(1-\varepsilon))^2\, v + 2\tau\mu \leq \tau^2$.
        The value $1$ is excluded: for $c_1 = 0$ the tangent at $\tau$ is
        $\chi \leq 0$. At a Bernstein cell the tangent point must also satisfy
        $\mu_r < \tau - c_{i\ell}$ (Section~\ref{sec:consistency_checks}); otherwise
        no point with $v \geq 0$ satisfies the row (Section~\ref{sec:dyn_rainflow}).
  \item \texttt{pwl}: a piecewise tangent with $P = \texttt{pwl\_points}$ segments
        (default $8$). The breakpoints split $[0, \tau]$ into equal segments
        $[(p-1)\tau/P,\; p\tau/P]$ with midpoints $\hat\mu_p = (p - \tfrac12)\tau/P$,
        $p = 1, \ldots, P$. One binary $w_{i\ell kp}$ per cell, day and segment
        selects the segment that contains the mean:
        \begin{align}
          \textstyle\sum_{p=1}^{P} w_{i\ell kp} &= 1, \label{eq:a_pwl_sum} \\
          w_{i\ell kp} = 1 &\;\Rightarrow\;
            (p-1)\,\tau/P \leq \mu_{i\ell k} \leq p\,\tau/P, \label{eq:a_pwl_seg} \\
          w_{i\ell kp} = 1 &\;\Rightarrow\;
            \chi_{i\ell k} \leq f(\hat\mu_p) + f'(\hat\mu_p)\,(\mu_{i\ell k} - \hat\mu_p).
            \label{eq:a_pwl_cap}
        \end{align}
        The implications are Gurobi indicator constraints, whatever the value of
        \texttt{encoding}. With $P = 1$ the form is the single tangent at $\tau/2$,
        with no binaries.
\end{itemize}
At a cell whose reliability row is linear (Markov, Chernoff, inverse Gaussian, Gamma)
the row is built in its exact form, and the value of \texttt{reliability\_impl} is
ignored. An unknown value of \texttt{reliability\_impl} is an input error at every
cell (Section~\ref{sec:consistency_checks}).

\begin{proposition}[Safety of the linear forms]
\label{prop:a_tangent}
For every $\mu_r$ and every $\mu$,
$f(\mu_r) + f'(\mu_r)\,(\mu - \mu_r) = f(\mu) - c_2\,(\mu - \mu_r)^2 \leq f(\mu)$.
Therefore every point that satisfies the \texttt{tangent} or the \texttt{pwl} form
satisfies the \texttt{exact} form, and both linear forms lie inside the certificate
set $\mathcal{C}_{i\ell}$.
\end{proposition}
\begin{proof}
$f$ is a polynomial of degree $2$ in $\mu$ with leading coefficient $c_2$, so its
Taylor expansion at $\mu_r$ is exact: $f(\mu) = f(\mu_r) + f'(\mu_r)(\mu - \mu_r) +
c_2(\mu - \mu_r)^2$. Since $c_2 > 0$, the tangent lies below $f$. Each linear form
caps $\chi$ by one such tangent and keeps the side row. Therefore it implies
Eq.~\eqref{eq:a_quad}.
\end{proof}

\emph{Remark (accuracy).} The \texttt{tangent} form loses $c_2(\mu - \mu_r)^2$ in
$\chi$: it is exact at $\mu = \mu_r$ and loosest far from it. A tangent point near the
operating mean of the cell is the best choice. In the \texttt{pwl} form the loss is
$c_2(\mu - \hat\mu_p)^2 \leq c_2\,\tau^2/(4P^2)$. The midpoint of the segment that
contains $\mu$ is the nearest tangent point to $\mu$. So the \texttt{pwl} form
accepts exactly the points with $\chi \leq \max_p\, [f(\hat\mu_p) + f'(\hat\mu_p)(\mu -
\hat\mu_p)]$, the largest set that the $P$ tangents allow. The price is $P$ binaries
per cell and day.

\emph{Example (Cantelli, $\varepsilon = 0.1$, $\kappa_C^2 = 9$, $\tau = 1$, $v =
0.004$).} The largest admissible mean is $1 - \sqrt{9 \cdot 0.004} \approx 0.810$
with \texttt{exact}; $0.482$ with \texttt{tangent} at $0$ ($9v + 2\mu \leq 1$);
$0.714$ with \texttt{tangent} at $0.5\,\tau$ ($9v \leq 0.75 - \mu$); and $0.810$
with \texttt{pwl} and $P = 8$. The tangent at $0$ never admits a mean above
$\tau/2$, even with zero variance.

% ----------------------------------------------------------------
\subsection{Gaussian Model}
\label{sec:dyn_gaussian}
% ----------------------------------------------------------------

This section applies to the cells $(i,\ell)$ with \texttt{model} $=$
\texttt{gaussian}. Inside it, the subscript $i\ell$ is dropped from the per-cell
parameters and from the states; $k$ is the day.

\textbf{Model.} A mission increment is Gaussian,
$\mathcal{N}(\mu^{\inc}_{jk}, v^{\inc}_{jk})$ (inputs \texttt{mu} and \texttt{v});
the initial damage is $\mathcal{N}(\mu_0, v_0)$ and each replacement draw is
$\mathcal{N}(\mu^{\new}, v^{\new})$. A known damage has variance 0.
\textbf{Assumption N1 (independence).} The increments, the initial damage and the
replacement draws are mutually independent. A repair acts on the damage path by
the law of the cell (Section~\ref{sec:repair_laws}). The Gaussian law gives
negative damage a small probability; the input checks of
Section~\ref{sec:consistency_checks} ($\mu \geq 3\sqrt{v}$ for the initial
damage and for each increment) keep it at most $\Phi(-3) = 0.135\,\%$ for each of
these terms.

\textbf{Tracked states and repair laws.} The cell tracks the mean $\mu_k$ and the
variance $v_k$, and under $\ARDone$ their anchors $g^{\mu}_k$ and $g^{v}_k$. Both
laws are accepted. The exact day maps are
\begin{align}
  \mu_k &= \begin{cases}
    \mu_{k-1} + \delta\mu_k & \text{no intervention on day } k,\\
    (1-\rho)\,\mu_{k-1} & \text{repair, } \ARDinf,\\
    (1-\rho)\,\mu_{k-1} + \rho\, g^{\mu}_{k-1} & \text{repair, } \ARDone,\\
    \mu^{\new} & \text{replacement,}
  \end{cases}
  \label{eq:g_map_mu} \\
  v_k &= \begin{cases}
    v_{k-1} + \delta v_k & \text{no intervention on day } k,\\
    (1-\rho)^2\, v_{k-1} & \text{repair, } \ARDinf,\\
    (1-\rho)^2\, v_{k-1} + \bigl(1 - (1-\rho)^2\bigr)\, g^{v}_{k-1}
      & \text{repair, } \ARDone,\\
    v^{\new} & \text{replacement,}
  \end{cases}
  \label{eq:g_map_v}
\end{align}
with $\delta\mu_k = \sum_{j=1}^{M} x_{ijk}\,\mu^{\inc}_{jk}$ and
$\delta v_k = \sum_{j=1}^{M} x_{ijk}\,v^{\inc}_{jk}$.

\begin{proposition}[The state stays Gaussian]
\label{prop:g_exact}
Under N1, on every day of a schedule $D_k \sim \mathcal{N}(m_k, V_k)$, where
$(m_k, V_k)$ follow the maps~\eqref{eq:g_map_mu} and~\eqref{eq:g_map_v}, with
anchors equal to the mean and the variance just after the last intervention.
\end{proposition}

\begin{proof}
A mission adds an independent Gaussian term, and $\ARDinf$ scales the damage by
$1-\rho$; both keep the law Gaussian and give the first two cases. Under
$\ARDone$, $D^- = G + \Delta$, with $G$ the damage just after the last
intervention and $\Delta$ the increments added since; $G$ and $\Delta$ are
independent Gaussian variables by N1. Then $D^+ = G + (1-\rho)\Delta$ is Gaussian,
with mean $m_G + (1-\rho)(m^- - m_G)$ and variance
\[
  V_G + (1-\rho)^2\, V_\Delta = V_G + (1-\rho)^2\,(V^- - V_G)
  = (1-\rho)^2\, V^- + \bigl(1 - (1-\rho)^2\bigr)\, V_G .
\]
A replacement starts from the replacement draw. Therefore the law is Gaussian on
every day, with the moments of the maps.
\end{proof}

So $\ARDone$ is exact for this model. Its variance map uses the factor
$(1-\rho)^2$ and the variance anchor $g^{v}$. The rows are the generic
rows of Section~\ref{sec:encoding}, with the data below;
$V_{\max} = \tau^2/\kappa_G^2$, with $\kappa_G = \Phi^{-1}(1-\varepsilon)$.

\begingroup\small
\begin{longtable}{@{}llllll@{}}
\toprule
State & Day increment & Factor $a$ & Natural range & Initial value & Replacement value \\
\midrule
$\mu$; anchor $g^{\mu}$ & $\delta\mu_k$ & $1-\rho$ & $\tau$ & $\mu_0$ & $\mu^{\new}$ \\
$v$; anchor $g^{v}$ & $\delta v_k$ & $(1-\rho)^2$ & $V_{\max}$ & $v_0$ & $v^{\new}$ \\
\bottomrule
\end{longtable}
\endgroup

\subsubsection{Reliability row}

Gaussian cells need $\varepsilon < 0.5$, so that $\kappa_G = \Phi^{-1}(1-\varepsilon) > 0$.
For $D \sim \mathcal{N}(m, V)$,
$\Pr(D > \tau) = 1 - \Phi\bigl((\tau - m)/\sqrt{V}\bigr) \leq \varepsilon$ exactly
when $m + \kappa_G\sqrt{V} \leq \tau$. With $\kappa_G > 0$ and $\tau - \mu \geq 0$ this is
equivalent to the squared row, imposed on every day $k = 1, \ldots, T$:
\begin{equation}
  \kappa_G^2\, v_k \leq (\tau - \mu_k)^2, \qquad \mu_k \leq \tau.
  \label{eq:g_row}
\end{equation}

\emph{Validity.} Let $(m_k, V_k) \leq (\mu_k, v_k)$ be the true moments
(Proposition~\ref{prop:g_exact} and Section~\ref{sec:encoding}). If $V_k = 0$,
then $D_k = m_k \leq \tau$. Otherwise $\mu_k < \tau$, because $\mu_k = \tau$
forces $v_k = 0$, and
$\Pr(D_k > \tau) = 1 - \Phi\bigl((\tau - m_k)/\sqrt{V_k}\bigr)
\leq 1 - \Phi\bigl((\tau - \mu_k)/\sqrt{v_k}\bigr) \leq \varepsilon$, because
$(\tau - \mu)/\sqrt{v}$ decreases in $\mu$ and in $v$.

\textbf{Forms.} In the form \texttt{exact}, Eq.~\eqref{eq:g_row} is a
\emph{nonconvex} quadratic constraint: $(\mu, v) = (\tau, 0)$ and
$(0, V_{\max})$ satisfy it, and their midpoint $(\tau/2, V_{\max}/2)$ does not,
since $\tau^2/2 > \tau^2/4$. The solver handles it through Gurobi's nonconvex
quadratic support (parameter \texttt{NonConvex} $= 2$). The row has the form
$v_k \leq f(\mu_k)$ with the convex function
\[
  f(\mu) = \frac{(\tau - \mu)^2}{\kappa_G^2}, \qquad \text{side row } \mu \leq \tau,
\]
so the forms \texttt{tangent} and \texttt{pwl} of
Section~\ref{sec:linear_forms} apply, and both keep the side row.

\emph{Remark (the level $\varepsilon$ and the range $V_{\max}$).}
$V_{\max} = f(0) = \tau^2/\kappa_G^2$ is the largest variance that the row admits, at
$\mu = 0$, so it is the natural range of $v$ and the big-M constant of its rows;
the input checks of Section~\ref{sec:consistency_checks} keep $v_0$, its starting
anchor and $v^{\new}$ inside it.
$V_{\max}/\tau^2$ is $0.18$ at $\varepsilon = 0.01$, $0.61$ at $0.1$, $3.6$ at
$0.3$ and about $1600$ at $0.49$. It grows without bound as $\varepsilon \to 0.5$,
and at $\varepsilon \geq 0.5$ the squared row does not describe the reliability
requirement, because $\kappa_G \leq 0$. Hence the condition $\varepsilon < 0.5$
(Section~\ref{sec:consistency_checks}); the range is large only close to 0.5.

\textbf{Optional row.} If \texttt{v\_max\_user} is given for the cell, the row
$v_k \leq V^{\mathrm{user}}$ is added on every day.

\textbf{Facts used by the encoding and loop arguments.}
(M1) The day maps are nondecreasing in $(\mu, v, g^{\mu}, g^{v})$: their
coefficients lie in $[0, 1]$.
(M2) The true moments follow the maps exactly (Proposition~\ref{prop:g_exact}).
(M3) If the tracked states satisfy Eq.~\eqref{eq:g_row} and the true moments are
not larger, then $\Pr(D_k > \tau) \leq \varepsilon$ (the validity paragraph).
The loop row of Section~\ref{sec:loop} covers $\mu$, $v$ and, under $\ARDone$,
$g^{\mu}$ and $g^{v}$.

% ----------------------------------------------------------------
\subsection{Inverse Gaussian Model}
\label{sec:dyn_ig}
% ----------------------------------------------------------------

This section applies to the cells $(i,\ell)$ with \texttt{model} $=$
\texttt{inverse\_gaussian} (IG). Inside it, the subscript $i\ell$ is dropped from
the per-cell parameters and from the states; $k$ is the day.

\textbf{Model.} $\mathrm{IG}(\mu, \lambda)$ is the inverse Gaussian law with mean
$\mu > 0$ and shape $\lambda > 0$. Its density is
\[
  f(x; \mu, \lambda) = \sqrt{\frac{\lambda}{2\pi x^3}}\,
  \exp\!\Bigl(-\frac{\lambda\,(x - \mu)^2}{2\mu^2 x}\Bigr), \qquad x > 0,
\]
and its variance is $\mu^3/\lambda$. For $c > 0$,
$c\,\mathrm{IG}(\mu, \lambda) = \mathrm{IG}(c\mu, c\lambda)$. The cell has one
dispersion $\eta > 0$ (input \texttt{eta}), and a term of mean $a$ has the law
$\mathrm{IG}(a, a/\eta) = a\,Z$ with $Z \sim \mathrm{IG}(1, 1/\eta)$. Its density is
$\sqrt{a/(2\pi\eta x^3)}\,\exp\bigl(-(x - a)^2/(2\eta a x)\bigr)$, and $\eta$ is
its squared coefficient of variation. Every term has this law with its own mean:
a mission increment (mean $\mu^{\inc}_{jk}$), the initial damage ($\mu_0$) and
each replacement draw ($\mu^{\new}$). A known damage $a$ is the term $a \cdot 1$.
The terms may depend on each other in any way: this model needs no independence
assumption.

\textbf{Tracked state and repair laws.} The cell tracks the mean $\mu_k$, and
under $\ARDone$ its anchor $g^{\mu}_k$. Both laws are accepted. The exact day map
of the mean is
\begin{equation}
  \mu_k = \begin{cases}
    \mu_{k-1} + \delta\mu_k & \text{no intervention on day } k,\\
    (1-\rho)\,\mu_{k-1} & \text{repair, } \ARDinf,\\
    (1-\rho)\,\mu_{k-1} + \rho\, g^{\mu}_{k-1} & \text{repair, } \ARDone,\\
    \mu^{\new} & \text{replacement,}
  \end{cases}
  \label{eq:ig_map}
\end{equation}
with $\delta\mu_k = \sum_{j=1}^{M} x_{ijk}\,\mu^{\inc}_{jk}$. The rows are the
generic rows of Section~\ref{sec:encoding}, with the data below.

\begingroup\small
\begin{longtable}{@{}llllll@{}}
\toprule
State & Day increment & Factor $a$ & Natural range & Initial value & Replacement value \\
\midrule
$\mu$; anchor $g^{\mu}$ & $\delta\mu_k$ & $1-\rho$ & $\tau$ & $\mu_0$ & $\mu^{\new}$ \\
\bottomrule
\end{longtable}
\endgroup

\subsubsection{Reliability row}

Between replacements, the damage on day $k$ is
$D_k = \sum_t w_t X_t = \sum_t (w_t a_t)\, Z_t$, where $X_t = a_t Z_t$ are the
terms added since the last replacement (or day 0), and each $Z_t$ has the law of
$Z$ or equals 1. Each weight $w_t \in (0, 1]$ is a product of factors $1-\rho$:
$\ARDinf$ multiplies every weight by $1-\rho$, and $\ARDone$, with
$D^+ = G + (1-\rho)\,\Delta$, multiplies only the weights of the terms added since
the last intervention. When the starting anchor $g^{\mu}_0$ is below $\mu_0$, the
initial damage $\mu_0 Z_0$ is read as the two terms $g^{\mu}_0 Z_0$ and
$(\mu_0 - g^{\mu}_0) Z_0$. The mean $\sum_t w_t a_t$ is at most the tracked mean
$\mu_k$ (Section~\ref{sec:encoding}).

\begin{proposition}[Stop-loss order]
\label{prop:ig_order}
Let $D = \sum_t w_t a_t Z_t$ with repair weights $w_t \in (0, 1]$, means
$a_t \geq 0$, $\sum_t w_t a_t \leq \mu$, and each $Z_t$ either with the law of $Z$
or equal to 1, under any joint law. Then
$\mathbb{E}(D - d)_+ \leq \mathbb{E}(\mu Z - d)_+$ for every $d \geq 0$.
\end{proposition}

\begin{proof}
Let $c_t = w_t a_t$ and $A = \sum_t c_t$. If $A = 0$, then $D = 0$. Otherwise, $x \mapsto (x - d)_+$
is convex and nondecreasing, $Z \geq 0$, and
$(A - d)_+ \leq \mathbb{E}(AZ - d)_+$ by Jensen's inequality ($\mathbb{E}Z = 1$),
so
\[
  \mathbb{E}(D - d)_+
  = \mathbb{E}\Bigl(\sum_t \frac{c_t}{A}\,(A Z_t) - d\Bigr)_+
  \leq \sum_t \frac{c_t}{A}\,\mathbb{E}(A Z_t - d)_+
  \leq \mathbb{E}(A Z - d)_+
  \leq \mathbb{E}(\mu Z - d)_+ .
\]
\end{proof}

Let $q$ be the $(1-\varepsilon)$-quantile of $Z$, $\Pr(Z > q) = \varepsilon$,
and let $C = \mathbb{E}[Z \mid Z > q]$, the conditional value at risk (CVaR) of
$Z$ at the level $1 - \varepsilon$. The cap and the row, on every day
$k = 1, \ldots, T$, are
\begin{equation}
  \bar\mu = \frac{\tau}{C}, \qquad \mu_k \leq \bar\mu.
  \label{eq:ig_row}
\end{equation}

\begin{theorem}[CVaR cap]
\label{thm:ig_cap}
Let $D \geq 0$ satisfy $\mathbb{E}(D - d)_+ \leq \mathbb{E}(\mu Z - d)_+$ for every
$d \geq 0$, with $0 \leq \mu \leq \bar\mu$. Then $\Pr(D > \tau) \leq \varepsilon$.
\end{theorem}

\begin{proof}
If $\mu = 0$, then $\mathbb{E}D \leq 0$ and $D = 0$. Let $\mu > 0$ and $d = \mu q$.
Since $C > q$, $d < \mu C \leq \tau$. For every $x$,
$\mathbf{1}\{x > \tau\} \leq (x - d)_+/(\tau - d)$, and
$\mathbb{E}(Z - q)_+ = \varepsilon\,(C - q)$. Then
\[
  \Pr(D > \tau) \leq \frac{\mathbb{E}(D - d)_+}{\tau - d}
  \leq \frac{\mu\,\mathbb{E}(Z - q)_+}{\tau - \mu q}
  = \frac{\varepsilon\,\mu\,(C - q)}{\tau - \mu q} \leq \varepsilon,
\]
where the last step is $\mu C \leq \tau$.
\end{proof}

By Proposition~\ref{prop:ig_order} with $\mu = \mu_k$ and
Theorem~\ref{thm:ig_cap}, the row~\eqref{eq:ig_row} gives
$\Pr(D_k > \tau) \leq \varepsilon$ on every day, under both laws. The row is
linear: it is built in its exact form, and a value of \texttt{reliability\_impl}
at the cell is ignored (Section~\ref{sec:linear_forms}).

\textbf{Computation of $\bar\mu$} (offline, once per cell, before the model is
built). With $r(x) = 1/\sqrt{\eta x}$, the distribution function of $Z$ is
\[
  F(x) = \Phi\bigl(r(x)\,(x - 1)\bigr)
  + e^{2/\eta}\,\Phi\bigl(-r(x)\,(x + 1)\bigr).
\]
The second term overflows for small $\eta$; it is computed as
$\exp\bigl(2/\eta + \ln\Phi(-r(x)(x+1))\bigr)$, with $\ln\Phi$ from a log-space
routine (\texttt{scipy.special.log\_ndtr}).
\begin{enumerate}[leftmargin=1.5em, nosep]
  \item $q$ is the root of $1 - F(q) = \varepsilon$, found by bisection or by
        \texttt{brentq} on $\ln q$. The method \texttt{isf} of
        \texttt{scipy.stats.invgauss}, with \texttt{mu} $= \eta$ and
        \texttt{scale} $= 1/\eta$, gives the same value.
  \item The tail mean has the closed form
        $\mathbb{E}[Z;\, Z > q] = F(1/q)$, so
        $\bar\mu = \tau\,\varepsilon/F(1/q)$. Derivation: substituting into the
        density gives $z\, f_Z(z) = z^{-2} f_Z(1/z)$, and the change of variable
        $u = 1/z$ gives
        $\int_q^\infty z\, f_Z(z)\, dz = \int_0^{1/q} f_Z(u)\, du = F(1/q)$.
\end{enumerate}
Values of $\bar\mu/\tau$:
\begin{center}
\begin{tabular}{@{}lccccc@{}}
\toprule
 & $\eta = 0.05$ & $0.25$ & $1$ & $4$ & $10$ \\
\midrule
$\varepsilon = 0.001$ & 0.491 & 0.242 & 0.101 & 0.0370 & 0.0190 \\
$\varepsilon = 0.01$ & 0.570 & 0.322 & 0.156 & 0.0679 & 0.0401 \\
$\varepsilon = 0.1$ & 0.693 & 0.477 & 0.299 & 0.183 & 0.141 \\
\bottomrule
\end{tabular}
\end{center}

\emph{Remark (the largest cap for this order).} Let $\mu > \bar\mu$. For a tail
probability $\varepsilon' > \varepsilon$ close to $\varepsilon$, with
$(1-\varepsilon')$-quantile $q'$, still $\mu\,\mathbb{E}[Z \mid Z > q'] > \tau$.
The two-point law
$D = \mu\,\mathbb{E}\bigl[Z \mid \mathbf{1}\{Z > q'\}\bigr]$ is below $\mu Z$ in
the stop-loss order, by Jensen's inequality, and $\Pr(D > \tau) \geq \varepsilon'$.
So no larger cap is safe for every law that Proposition~\ref{prop:ig_order}
allows. That class also contains laws that no schedule produces, so the row can
reject some safe schedules.

\emph{Remark (the quantile cap is not the row).} The cap $\tau/q$ makes the model
law $\mu Z$ meet the target exactly: $\Pr(\mu Z > \tau) = \varepsilon$ at
$\mu = \tau/q$. A sum of terms below $\mu Z$ in the stop-loss order can have a
larger tail. Example: $\eta = 10$, $\varepsilon = 0.1$, so $\tau/q = 0.473\,\tau$;
three independent missions of equal mean, with total mean $\tau/q$, give
$\Pr(D > \tau) = 0.118$. The CVaR cap is $0.141\,\tau$ there. At
$\varepsilon = 0.01$ and $\eta$ from 0.05 to 10, $\bar\mu$ is 0.93 to 0.59 times
$\tau/q$.

\textbf{Facts used by the encoding and loop arguments.}
(M1) The day map~\eqref{eq:ig_map} is nondecreasing in $(\mu, g^{\mu})$.
(M2) The true mean follows it exactly.
(M3) If $\mu_k \leq \bar\mu$ and the true mean is not larger, then
$\Pr(D_k > \tau) \leq \varepsilon$ (Proposition~\ref{prop:ig_order},
Theorem~\ref{thm:ig_cap}). The loop row of Section~\ref{sec:loop} covers $\mu$
and, under $\ARDone$, $g^{\mu}$.
Each day of each repetition has its own set of terms; the argument above holds
for each of them, because it needs neither independence nor a fixed set of
terms.

% ----------------------------------------------------------------
\subsection{Wiener Model}
\label{sec:dyn_wiener}
% ----------------------------------------------------------------

This section applies to the cells $(i,\ell)$ with \texttt{model} $=$
\texttt{wiener}. Inside it, the subscript $i\ell$ is dropped from the per-cell
parameters and from the states; $k$ is the day.

\textbf{Model.} The damage is a Wiener process with drift, observed once a day. A
mission day adds a Gaussian increment $\mathcal{N}(\mu^{\inc}_{jk}, \sigma^2)$:
the mean comes from the input \texttt{mu}, and the diffusion constant
$\sigma > 0$ (input \texttt{sigma}) is the same for every mission, because it
describes the material of the component, not the route. The initial damage is
$\mathcal{N}(\mu_0, v_0)$ with $v_0 \geq 0$, and each replacement draw is
$\mathcal{N}(\mu^{\new}, v^{\new})$. \textbf{Assumption W1 (independence).} These
terms are mutually independent.

\textbf{Repairs act on the mean only} (a modelling choice). A repair shifts the
damage by a known amount: $D^+ = D^- - \rho\, m^-$ under $\ARDinf$ and
$D^+ = D^- - \rho\,(m^- - g^{\mu})$ under $\ARDone$, with $m^- = \mathbb{E}D^-$.
The law stays Gaussian, the variance is unchanged, and only a replacement lowers
it.

\textbf{Tracked states and repair laws.} The cell tracks $\mu_k$ and $v_k$, and
under $\ARDone$ the mean anchor $g^{\mu}_k$; the variance has no anchor. Both
laws are accepted. The mean follows the map~\eqref{eq:g_map_mu} of the Gaussian
model. The variance map is
\begin{equation}
  v_k = \begin{cases}
    v_{k-1} + \delta v_k & \text{no intervention on day } k,\\
    v_{k-1} & \text{repair, both laws},\\
    v^{\new} & \text{replacement,}
  \end{cases}
  \qquad \delta v_k = \sigma^2 \sum_{j=1}^{M} x_{ijk}.
  \label{eq:w_map_v}
\end{equation}
Under W1 the damage is $\mathcal{N}(m_k, V_k)$ on every day, with $(m_k, V_k)$
given by these maps: a mission adds an independent Gaussian term, and a repair or
a replacement acts as stated above.

The rows are the generic rows of Section~\ref{sec:encoding}, with the data below.
The variance takes the repair map $\psi = v_{k-1}$, that is the factor $a = 1$
under both laws, and no anchor. A depot day adds no increment, so for $v$ the
generic no-intervention and repair rows together equal the single row
\begin{equation}
  v_k \geq v_{k-1} + \delta v_k - V_{\max}\, x^r_k,
  \label{eq:w_vrow}
\end{equation}
which may replace them in the big-M form.

\begingroup\small
\begin{longtable}{@{}llllll@{}}
\toprule
State & Day increment & Factor $a$ & Natural range & Initial value & Replacement value \\
\midrule
$\mu$; anchor $g^{\mu}$ & $\delta\mu_k$ & $1-\rho$ & $\tau$ & $\mu_0$ & $\mu^{\new}$ \\
$v$; no anchor & $\delta v_k$ & $1$ & $V_{\max}$ & $v_0$ & $v^{\new}$ \\
\bottomrule
\end{longtable}
\endgroup

\emph{Remark (a repair of the path).} A repair of the damage path, as in the
Gaussian model, gives the same mean and the variance $(1-\rho)^2 V$ under
$\ARDinf$, or $(1-\rho)^2 V + (1 - (1-\rho)^2)\,V_G$ under $\ARDone$, with
$V_G \leq V$ the variance just after the last intervention. Both are at most $V$.
The reliability row increases in $v$, so the rows of this section are safe for
that reading as well.

\subsubsection{Reliability row}

The row is the Gaussian row~\eqref{eq:g_row}, with $\kappa_G = \Phi^{-1}(1-\varepsilon)$
and $\varepsilon < 0.5$, on every day $k = 1, \ldots, T$; its validity argument,
its forms (\texttt{exact}, \texttt{tangent}, \texttt{pwl}, with the same
function $f$ and side row) and the remark on $V_{\max} = \tau^2/\kappa_G^2$ are those of
Section~\ref{sec:dyn_gaussian}. If \texttt{v\_max\_user} is given for the cell,
the row $v_k \leq V^{\mathrm{user}}$ is added on every day.

\textbf{Facts used by the encoding and loop arguments.}
(M1) The day maps are nondecreasing in $(\mu, v, g^{\mu})$.
(M2) The true moments follow them exactly.
(M3) As for the Gaussian model. The loop row of Section~\ref{sec:loop} covers
$\mu$, $v$ and, under $\ARDone$, $g^{\mu}$.

\emph{Remark (the loop forces a replacement).} The variance grows by $\sigma^2$
on every mission day and falls only at a replacement. If the train flies $n$
missions in the operating phase and the component is not replaced there, then
$v_T \geq v_{H_1} + n\,\sigma^2 > v_{H_1}$, and the loop row $v_T \leq v_{H_1}$
fails. So every feasible schedule includes a replacement of the component in the
operating phase whenever its train flies a mission there. No input check enforces this. With
\texttt{allow\_replacement} false, the train of a Wiener cell cannot fly in the
operating phase.

% ----------------------------------------------------------------
\subsection{Gamma Model}
\label{sec:dyn_gamma}
% ----------------------------------------------------------------

This section applies to the cells $(i,\ell)$ with \texttt{model} $=$
\texttt{gamma}. Inside it, the subscript $i\ell$ is dropped from the per-cell
parameters and from the states; $k$ is the day.

\textbf{Model.} $\mathrm{Gamma}(a, \beta)$ is the Gamma law with shape $a > 0$
and rate $\beta > 0$; its mean is $a/\beta$. Every term that adds damage is a
Gamma variable with its own rate:
\begin{itemize}[leftmargin=1.5em, nosep]
  \item a mission increment with mean $\mu^{\inc}_{jk} > 0$ (input \texttt{mu})
        and rate $\beta_{jk}$ (inputs \texttt{gamma\_beta} and
        \texttt{gamma\_beta\_trans}) is
        $\mathrm{Gamma}(\beta_{jk}\,\mu^{\inc}_{jk},\, \beta_{jk})$;
  \item the initial damage with mean $\mu_0 > 0$ is
        $\mathrm{Gamma}(\beta_0\,\mu_0,\, \beta_0)$ when its rate $\beta_0$ is
        given (input \texttt{gamma\_beta\_0}); when \texttt{gamma\_beta\_0} is
        omitted, the initial damage is known: the constant $\mu_0$;
  \item each replacement draw with mean $\mu^{\new} > 0$ is
        $\mathrm{Gamma}(\beta^{\new}\mu^{\new},\, \beta^{\new})$ when its rate
        $\beta^{\new}$ is given (input \texttt{gamma\_beta\_new}); when
        \texttt{gamma\_beta\_new} is omitted, it is the known constant
        $\mu^{\new}$.
\end{itemize}
\textbf{Assumption G1 (independence).} All these terms are mutually independent.

A repair acts on the damage path by the law of the cell
(Section~\ref{sec:repair_laws}): $\ARDinf$ gives $D^+ = (1-\rho)\,D^-$, and
$\ARDone$ gives $D^+ = D^{\mathrm{last}} + (1-\rho)\,(D^- - D^{\mathrm{last}})$.
After a repair the damage is, in general, not a Gamma variable: it is a weighted
sum of independent Gamma terms and known constants. The reliability row below
holds for every such sum.

\textbf{Tracked state and repair laws.} The cell tracks the mean $\mu_k$ only,
and under $\ARDone$ its anchor $g^{\mu}_k$. Both laws are accepted. The shape of
the damage is not tracked. The exact day map of the mean is
\begin{equation}
  \mu_k = \begin{cases}
    \mu_{k-1} + \delta\mu_k & \text{no intervention on day } k,\\
    (1-\rho)\,\mu_{k-1} & \text{repair, } \ARDinf,\\
    (1-\rho)\,\mu_{k-1} + \rho\, g^{\mu}_{k-1} & \text{repair, } \ARDone,\\
    \mu^{\new} & \text{replacement,}
  \end{cases}
  \label{eq:gm_map}
\end{equation}
with $\delta\mu_k = \sum_{j=1}^{M} x_{ijk}\,\mu^{\inc}_{jk}$. The map is exact
for the true mean $\mathbb{E}D_k$, because expectation is linear. The rows are
the generic rows of Section~\ref{sec:encoding}, with the data below.

\begingroup\small
\begin{longtable}{@{}llllll@{}}
\toprule
State & Day increment & Factor $a$ & Natural range & Initial value & Replacement value \\
\midrule
$\mu$; anchor $g^{\mu}$ & $\delta\mu_k$ & $1-\rho$ & $\tau$ & $\mu_0$ & $\mu^{\new}$ \\
\bottomrule
\end{longtable}
\endgroup

\textbf{Smallest rate.} The cell's smallest rate is
\begin{equation}
  \beta^* = \min\Bigl(\{\beta_{jk} : \mu^{\inc}_{jk} > 0\}
    \cup \{\beta_0 : \mu_0 > 0\} \cup \{\beta^{\new} : \mu^{\new} > 0\}\Bigr),
  \label{eq:gm_betastar}
\end{equation}
where $j$ runs over the missions and $k$ over the days of both phases, and
$\beta_0$, $\beta^{\new}$ enter only when they are given. So $\beta^*$ is the
smallest rate among the terms with a positive mean: a mission with
$\mu^{\inc}_{jk} = 0$ adds no term, and an omitted rate declares a known damage,
which is a constant and does not enter $\beta^*$. If no term has a rate, the
damage is a known constant, and the row below is $\mu_k \leq \tau$. Write
$t = \beta^*\tau$. Units of $1/\beta^*$ turn every law below into a law of rate
at least 1.

\subsubsection{Class of true states}

\begin{definition}[Class of weighted Gamma sums]
\label{def:gm_class}
Let $m \geq 0$. The class $\mathcal{C}(m)$ holds the laws of the random
variables
\[
  S = \delta + \sum_{i=1}^{n} c_i\, Y_i, \qquad
  Y_i \sim \mathrm{Gamma}(a_i, 1) \text{ independent}, \quad
  c_i \in (0, 1], \quad a_i > 0, \quad \delta \geq 0,
\]
with $n \geq 0$ and mean $\delta + \sum_i c_i a_i = m$. The term $c\,Y$ with
$Y \sim \mathrm{Gamma}(m'/c, 1)$ has mean $m'$; as $c \to 0$ it tends to the
constant $m'$.
\end{definition}

\begin{lemma}[True states belong to the class]
\label{lem:gm_class}
Let $D_k$ be the damage of a Gamma cell on day $k$ of a schedule, under either
law. Then the law of $\beta^* D_k$ belongs to $\mathcal{C}(\beta^*\,\mathbb{E}D_k)$.
\end{lemma}

\begin{proof}
Since the last replacement (or day 0), the damage is
$D_k = \delta + \sum_t w_t X_t$: the $X_t$ are the independent Gamma terms added
since then, $\delta \geq 0$ is the weighted sum of the known constants, and each
weight $w_t \in (0, 1]$ is a product of factors $1-\rho$. A mission adds a term
with weight 1. $\ARDinf$ multiplies every weight, constants included, by
$1-\rho$. $\ARDone$ multiplies by $1-\rho$ the weights of the terms and constants
added since the last intervention, and keeps the others. A replacement starts a
new sum with the replacement draw. When the starting anchor $g^{\mu}_0$ is below
$\mu_0$, the initial damage is read as two independent terms of rate $\beta_0$,
with means $g^{\mu}_0$ and $\mu_0 - g^{\mu}_0$, whose sum has the law
$\mathrm{Gamma}(\beta_0\mu_0, \beta_0)$; a known initial damage is read as the
two constants $g^{\mu}_0$ and $\mu_0 - g^{\mu}_0$. Write $X_t = Y_t/\beta_t$ with
$Y_t \sim \mathrm{Gamma}(a_t, 1)$ and $\beta_t \geq \beta^*$ by
Eq.~\eqref{eq:gm_betastar}. Then
\[
  \beta^* D_k = \beta^*\delta + \sum_t \frac{w_t\,\beta^*}{\beta_t}\, Y_t,
  \qquad \frac{w_t\,\beta^*}{\beta_t} \in (0, 1].
\]
The $Y_t$ are independent by Assumption G1. Therefore the law of $\beta^* D_k$
is in $\mathcal{C}(\beta^*\,\mathbb{E}D_k)$.
\end{proof}

\subsubsection{The largest tail at a given mean}

For $0 < m < t$, let $W(m, t)$ be the supremum of $\Pr(S > t)$ over the random
variables $S$ whose law is in $\mathcal{C}(m)$.
$Q(a, x) = \Pr(\mathrm{Gamma}(a, 1) > x)$ is the regularised upper incomplete
gamma function (\texttt{scipy.special.gammaincc(a, x)}). The one-term member
$c\,\mathrm{Gamma}(m/c, 1)$ has the tail $Q(m/c,\, t/c)$ at $t$.

A member of $\mathcal{C}(m)$ is described by its atoms: the pairs
$(m_i, c_i)$ of a term mean $m_i = c_i a_i$ and a scale $c_i$, and the constant
$\delta$, read as an atom at scale $c = 0$. Terms with equal scales merge into
one term, because a sum of independent Gamma variables of one rate is a Gamma
variable.

\begin{lemma}[First-order rates]
\label{lem:gm_rates}
Let $S$ be a member with at least one term, with density $f_S$. Since
$\delta \leq m < t$, $f_S$ is continuous at $t$. For $c \in (0, 1]$ let
\[
  \varphi_S(c) = \frac{1}{c}\int_0^\infty \Pr(t - y < S \leq t)\,
    \frac{e^{-y/c}}{y}\, dy,
  \qquad \varphi_S(0) = f_S(t).
\]
(a) Adding mean $\eta$ at scale $c$ (a term $c\,\mathrm{Gamma}(\eta/c, 1)$,
or the constant $\eta$ when $c = 0$) raises $\Pr(S > t)$ at the rate
$\varphi_S(c)$ as $\eta \to 0$; removing mean from an atom at scale $c$ lowers
it at the same rate. (b) Moving the scale of an atom $(m_i, c_i)$,
$c_i \in (0,1)$, at fixed mean $m_i$ changes $\Pr(S > t)$ at the rate
$m_i\, \partial\varphi_S/\partial c\,(c_i)$.
\end{lemma}

\begin{proof}
$S$ is infinitely divisible, with drift $\delta$ and Lévy density
$\ell(y) = \sum_i (m_i/c_i)\, e^{-y/c_i}/y$ on $y > 0$. A change of the
parameters at the rates $(\dot\delta, \dot\ell)$ changes the Laplace transform
$\mathbb{E}e^{-\sigma S} = \exp\bigl(-\sigma\delta - \int_0^\infty (1 - e^{-\sigma y})\,
\ell(y)\, dy\bigr)$ at the rate
$\int_0^\infty \mathbb{E}\bigl[e^{-\sigma(S+y)} - e^{-\sigma S}\bigr]\,\dot\ell(y)\, dy
- \sigma\dot\delta\,\mathbb{E}e^{-\sigma S}$. This is the transform of the rate
$\int_0^\infty \Pr(t - y < S \leq t)\,\dot\ell(y)\, dy + \dot\delta\, f_S(t)$ of
$\Pr(S > t)$. Case (a) is $\dot\ell(y) = c^{-1}e^{-y/c}/y$, or $\dot\delta = 1$
when $c = 0$. Case (b) is
$\dot\ell(y) = m_i\,\partial_c\bigl(c^{-1}e^{-y/c}/y\bigr)$ at $c = c_i$.
\end{proof}

\begin{theorem}[One term is the worst case]
\label{thm:gm_oneterm}
Let $0 < m < t$. Then
\begin{equation}
  W(m, t) = \max_{c \in (0, 1]} Q\!\left(\frac{m}{c},\, \frac{t}{c}\right).
  \label{eq:gm_oneterm}
\end{equation}
\end{theorem}

\begin{proof}
Let $\mathcal{C}_K(m)$ be the members with at most $K$ atoms, the constant
included, and let $K \geq 3$.
\begin{enumerate}[leftmargin=1.5em]
  \item \emph{A maximiser exists.} $\mathcal{C}_K(m)$ is the image of the
        compact set of atoms $(m_i, c_i)$ with $m_i \geq 0$, $\sum_i m_i = m$,
        $c_i \in [0, 1]$. The law of $S$ depends continuously on the atoms, and
        only the constant member $S = m$ has an atom, at $m \neq t$. So
        $\Pr(S > t)$ is continuous and attains its maximum at some $S^*$. Since
        $\Pr(m > t) = 0 < Q(m, t)$, $S^*$ has at least one term. Every member
        has finitely many atoms, so
        $W(m, t) = \sup_K \max_{\mathcal{C}_K(m)} \Pr(S > t)$.
  \item \emph{First-order conditions.} Write $\varphi^* = \varphi_{S^*}$. By
        Lemma~\ref{lem:gm_rates}: (i) moving mean between two atoms of positive
        mean must not raise the tail, so $\varphi^*$ takes one value $\kappa^*$
        at every atom of positive mean; (ii) at an atom with
        $c \in (0, 1)$, $\partial\varphi^*/\partial c = 0$; (iii) if $S^*$ has
        fewer than $K$ atoms of positive mean, a new atom can be opened at any
        $c \in [0, 1]$, so $\varphi^*(c) \leq \kappa^*$ on $[0, 1]$.
  \item \emph{A single-peak function.} Put $\lambda = 1/c$ and
        $h(u) = \Pr(t - u < S^* \leq t)/u$. Then
        $\varphi^*(\lambda) = \lambda\int_0^\infty h(u)\, e^{-\lambda u}\, du$,
        and $\varphi^*(\infty) = h(0^+) = f_{S^*}(t)$ is the value at the
        constant atom. $S^*$ is a constant plus a sum of independent Gamma
        variables, so its density is unimodal
        \cite[Theorem~4, p.~11]{roosta2015}, proved in
        \cite[proof of Theorem~4, pp.~194--195]{szekely2003}; the constant
        shifts the mode only. Then $g(v) = f_{S^*}(t - v)$ for $v < t$ and
        $g(v) = 0$ for $v \geq t$ is unimodal on $[0, \infty)$, and so is its
        running average $h(u) = u^{-1}\int_0^u g$: before the mode of $g$,
        $g \geq h$; after it, $u\,(g(u) - h(u)) = u\,g(u) - \int_0^u g$ is
        nonincreasing. Since $h' = (g - h)/u$, $h$ rises and then falls.
  \item \emph{Zero count.} For every constant $\kappa'$, $h - \kappa'$ changes
        sign at most twice. The Laplace transform of a function with at most
        $n$ sign changes has at most $n$ zeros in $\lambda > 0$, counted with
        multiplicity: at a sign change $u_1$, multiply the transform by
        $e^{\lambda u_1}$; its derivative in $\lambda$ is the transform of a
        function with one sign change fewer; Rolle's theorem and induction
        conclude. So $\varphi^* - \kappa' = \lambda\int_0^\infty (h - \kappa')\,
        e^{-\lambda u}\, du$ has at most two zeros in $\lambda > 0$. Also
        $\varphi^*(\lambda) \to 0$ as $\lambda \to 0^+$. Hence $\varphi^*$
        cannot rise above a level $\kappa' > 0$, fall below it and rise above it
        again: with $\varphi^*(0^+) = 0$ that would give three zeros. So
        $\varphi^*$ is quasi-concave. It is analytic in $\lambda > 0$ and not
        constant, so its maximum over $\lambda \in [1, \infty]$ is attained at
        one point only: two maximisers would make $\varphi^*$ constant between
        them.
  \item \emph{One term.} By (i) and (ii), an atom with $c \in (0, 1)$ is a zero
        of $\varphi^* - \kappa^*$ of multiplicity at least 2, and an atom at
        $c = 1$ a zero of multiplicity at least 1; the constant atom is at
        $\lambda = \infty$. By step~4, $S^*$ has at most two atoms of positive
        mean. As $K \geq 3$, (iii) applies, so $\kappa^*$ is the maximum of
        $\varphi^*$ over $[1, \infty]$, attained at every atom of positive mean.
        By step~4 there is only one such atom, and by step~1 it is not the
        constant. So $S^* = c^*\,\mathrm{Gamma}(m/c^*, 1)$ with
        $c^* \in (0, 1]$.
\end{enumerate}
Therefore $\max_{\mathcal{C}_K(m)} \Pr(S > t) = \max_{c \in (0,1]} Q(m/c, t/c)$
for every $K \geq 3$, and Eq.~\eqref{eq:gm_oneterm} follows from step~1.
\end{proof}

\begin{lemma}[Monotone in the mean]
\label{lem:gm_monotone}
$W(m, t)$ is nondecreasing in $m \in (0, t)$.
\end{lemma}

\begin{proof}
Let $m < m' < t$ and let the law of $S$ be in $\mathcal{C}(m)$. Then the law of
$S + (m' - m)$ is in $\mathcal{C}(m')$, and
$\Pr(S + m' - m > t) \geq \Pr(S > t)$. Therefore $W(m, t) \leq W(m', t)$.
\end{proof}

\subsubsection{Reliability row}

\textbf{The cap.} For $A \in [0, t)$ let
$w(A) = \max_{c \in (0, 1]} Q(A/c,\, t/c)$; it is continuous and nondecreasing
in $A$, because $Q(a, x)$ increases in $a$. The \emph{sharp cap}
$A_1(t, \varepsilon)$ is the largest $A \in [0, t)$ with
$w(A) \leq \varepsilon$, computed as described below. The row, on every day
$k = 1, \ldots, T$, is
\begin{equation}
  \mu_k \;\leq\; \frac{A_1(\beta^*\tau,\, \varepsilon)}{\beta^*}.
  \label{eq:gm_row}
\end{equation}
The row is linear: it is built in its exact form, and a value of
\texttt{reliability\_impl} at the cell is ignored (Section~\ref{sec:linear_forms}).

\begin{corollary}[The row is safe]
\label{cor:gm_safe}
Let the tracked mean satisfy Eq.~\eqref{eq:gm_row} on day $k$. Then
$\Pr(D_k > \tau) \leq \varepsilon$.
\end{corollary}

\begin{proof}
Let $m = \beta^*\,\mathbb{E}D_k$ and $A = A_1(\beta^*\tau, \varepsilon)$. The
tracked mean bounds the true mean from above (Section~\ref{sec:encoding}), so
$m \leq \beta^*\mu_k \leq A < t$. If $m = 0$, then $D_k = 0$. Otherwise
\[
  \Pr(D_k > \tau) = \Pr(\beta^* D_k > t)
  \leq W(m, t) \leq W(A, t) = w(A) \leq \varepsilon,
\]
by Lemma~\ref{lem:gm_class}, Lemma~\ref{lem:gm_monotone},
Theorem~\ref{thm:gm_oneterm} and the choice of $A$.
\end{proof}

\textbf{Computation of $A_1$} (offline, once per cell, before the model is built):
\begin{enumerate}[leftmargin=1.5em, nosep]
  \item $t = \beta^*\tau$. $\hat A$ is the root of $Q(\hat A, t) = \varepsilon$,
        found by bisection on $\hat A$ ($Q$ increases in its first argument;
        \texttt{gammaincinv} solves for the second argument and does not apply).
  \item $w(A)$: the function $c \mapsto Q(A/c, t/c)$ tends to 0 as $c \to 0$
        and has a single peak on $(0, 1]$ (checked numerically for $t$ from
        $10^{-3}$ to $100$ and $A/t$ from $0.05$ to $0.999$). Its maximum is found on a
        logarithmic grid of at least 400 values of $c$ in
        $[\min(10^{-4}, t/100),\, 1]$, then refined by a golden-section search
        between the neighbours of the best grid point. If the best grid point is
        the smallest value of $c$, the grid is extended downwards.
  \item If $\hat A < t$ and $w(\hat A) \leq \varepsilon$, then
        $A_1 = \hat A$. Otherwise $A_1$ is found by bisection on $A$ in
        $[\varepsilon t,\, \min(\hat A, t))$, to a relative tolerance of
        $10^{-9}$. The lower end satisfies $w(\varepsilon t) \leq \varepsilon$,
        by Markov's inequality $Q(a, x) \leq a/x$. The bisection returns the
        lower end of its last interval, so $w(A_1) \leq \varepsilon$ holds at
        the returned value.
\end{enumerate}

\emph{Remark (the sharp cap and the exact one-rate cap).} $\hat A$ is the exact
cap of a single Gamma law of rate $\beta^*$, and $A_1 \leq \hat A$, because
$c = 1$ is in the maximum. Let $t_1(\varepsilon)$ be the value of $t$ at which
$\partial_c Q(\hat A/c, t/c)$ at $c = 1$ changes sign: $t_1 = 0.436$, $0.448$
and $0.614$ for $\varepsilon = 0.001$, $0.01$ and $0.1$. For $t \geq t_1$,
$A_1 = \hat A$: on a grid of $t$ in $[t_1, 1000]$, no $c < 1$ gives a larger tail
than $c = 1$. Below $t_1$, $A_1 < \hat A$. For $\varepsilon = 0.01$:
\begin{center}
\begin{tabular}{@{}lccccc@{}}
\toprule
$t = \beta^*\tau$ & 0.01 & 0.1 & 0.3 & 1 & 10 \\
\midrule
$A_1/\hat A$ & 0.142 & 0.642 & 0.961 & 1 & 1 \\
\bottomrule
\end{tabular}
\end{center}
When every term has the same rate and no repair occurs, the true state is a single
Gamma law and $\hat A$ is exact for it. With several rates, or with repairs, the true state is a
weighted sum, and below $t_1$ the exact one-rate cap is not safe. Example:
$\varepsilon = 0.01$ and $\beta^*\tau = 0.1$, so $\hat A = 0.0055012$. A new
part flies one mission of mean $\hat A/\beta^*$ and rate $4.56\,\beta^*$. Its
mean meets the exact one-rate cap, but its tail is $1.565\,\varepsilon$. The
sharp cap is $A_1 = 0.642\,\hat A$ there.

\emph{Remark (below the threshold).} With $r = A/t$ and $s = t/c$,
$Q(A/c, t/c) = Q(r s, s)$, and $c \in (0, 1]$ means $s \geq t$. For
$t \leq t_1(\varepsilon)$ the peak of $s \mapsto Q(rs, s)$ at the cap lies at
$s \geq t$, so the maximum does not depend on $t$, and
$A_1(t, \varepsilon) = r^*(\varepsilon)\, t$. The mean cap is then
$A_1/\beta^* = r^*(\varepsilon)\,\tau$, whatever the rates:
$r^* = 0.00355$, $0.0353$ and $0.331$ for $\varepsilon = 0.001$, $0.01$ and
$0.1$.

\emph{Remark (the choice of $\beta^*$).} A larger $t$ shrinks the range
$s \geq t$, so $A_1(t, \varepsilon)/t$ is nondecreasing in $t$, and the mean cap
$A_1(\beta^*\tau, \varepsilon)/\beta^* = \tau\, A_1(t, \varepsilon)/t$ is
nondecreasing in $\beta^*$. Lemma~\ref{lem:gm_class} holds for every
$\beta^*$ at most the smallest rate. So the smallest rate gives the largest safe
cap of this form.

\emph{Remark (conservatism).} The class $\mathcal{C}(m)$ also contains laws that
no schedule produces. Below $t_1$ the row can therefore reject some safe
schedules. The row keeps each mission's own mean; $\beta^*$ only bounds how
dispersed each term may be.

\textbf{Facts used by the encoding and loop arguments.}
(M1) The day map~\eqref{eq:gm_map} is nondecreasing in $(\mu, g^{\mu})$.
(M2) The true mean follows the map~\eqref{eq:gm_map} exactly.
(M3) If the tracked mean satisfies Eq.~\eqref{eq:gm_row} and the true mean is
not larger, then $\Pr(D_k > \tau) \leq \varepsilon$
(Corollary~\ref{cor:gm_safe}). The loop row of Section~\ref{sec:loop} covers
$\mu$ and, under $\ARDone$, $g^{\mu}$.

% ----------------------------------------------------------------
\subsection{Rainflow Model}
\label{sec:dyn_rainflow}
% ----------------------------------------------------------------

This section applies to the cells $(i,\ell)$ with \texttt{model} $=$
\texttt{rainflow}. Inside it, the subscript $i\ell$ is dropped from the per-cell
parameters and from the states; $k$ is the day.

The rainflow model is \emph{distribution-free}. The law of a mission increment
is measured from recorded load histories by rainflow counting (Palmgren--Miner
damage), and it enters the model only through summary statistics: its mean, and,
as the bound requires, its variance, the width of its range, or its cumulant
generating function (CGF) at a fixed point $s$. No parametric family is assumed.
The reliability requirement $\Pr(D_k > \tau) \leq \varepsilon$ is certified by
one of five tail bounds, chosen per cell by \texttt{bound\_method}:
\texttt{markov}, \texttt{cantelli}, \texttt{hoeffding}, \texttt{bernstein},
\texttt{chernoff}.

\textbf{Modelling assumptions (normative).}
\begin{itemize}[leftmargin=1.5em]
  \item \textbf{A0 (non-negative damage).} Every increment, the initial damage
        and every replacement draw are non-negative.
  \item \textbf{A1 (independence).} The mission increments are mutually
        independent, and independent of the initial damage and of every
        replacement draw. Each replacement installs a new draw, independent of
        everything else. Every bound except Markov needs A1.
  \item \textbf{A2 (bounded ranges; Hoeffding and Bernstein).} Each mission
        increment lies in an interval of width $b^{\inc}_{jk}$ (input
        \texttt{support}). The initial damage is a known part plus independent
        terms, each in an interval; $R_0$ is the sum of their squared widths
        (input \texttt{R\_0}) and $b^0$ the largest width (input
        \texttt{support\_0}). The replacement draw is described in the same way by
        $R^{\new}$ and $b^{\new}$ (inputs \texttt{R\_new}, \texttt{support\_new}).
        A known damage has $R = b = 0$, the defaults.
  \item \textbf{A3 (repetition).} The row of the cell's bound holds on every day
        $k = 1, \ldots, T$, and after day $T$ the operating phase repeats with
        period $H_2$, schedule and increment profile included. This is the
        setting of the loop argument of Section~\ref{sec:loop}.
\end{itemize}

\textbf{Width of a range.} The input \texttt{support} of an increment $X$ is the
width $b$ of an interval $[l, l + b]$ that contains $X$ almost surely. Hoeffding's
inequality uses the widths. Bernstein's inequality needs an upper bound of the
centred term, and the width is one: $X - \mathbb{E}X \leq l + b - l = b$, since
$\mathbb{E}X \geq l$. A width $b$ is consistent with a variance $v$ only if
$v \leq b^2/4$ (Popoviciu's inequality:
$\mathrm{Var}\,X \leq \mathbb{E}(X - l - b/2)^2 \leq (b/2)^2$); the check is in
Section~\ref{sec:consistency_checks}.

\textbf{Tracked states per bound.} Every cell tracks the mean $\mu$. Each bound
adds at most one state, and under $\ARDone$ one anchor per tracked state
($g^{\mu}$ and $g^{v}$, $g^{R}$ or $g^{K}$). Both laws are accepted for every
bound.

\begingroup\small
\begin{longtable}{@{}lllll@{}}
\toprule
\texttt{bound\_method} & States & Increment inputs & Initial; replacement inputs & Forms \\
\midrule
\texttt{markov} & $\mu$ & \texttt{mu} & -- & \texttt{exact} \\
\texttt{cantelli} & $\mu$, $v$ & \texttt{mu}, \texttt{v}
  & \texttt{v\_0}; \texttt{replacement\_v} & all three \\
\texttt{hoeffding} & $\mu$, $R$ & \texttt{mu}, \texttt{support}
  & \texttt{R\_0}; \texttt{R\_new} & all three \\
\texttt{bernstein} & $\mu$, $v$ & \texttt{mu}, \texttt{v}, \texttt{support}
  & \texttt{v\_0}, \texttt{support\_0}; & all three \\
 & & & \texttt{replacement\_v}, \texttt{support\_new} & \\
\texttt{chernoff} & $\mu$, $K$ & \texttt{mu}, \texttt{cgf}, \texttt{s\_chernoff}
  & \texttt{cgf\_0}; \texttt{cgf\_new} & \texttt{exact} \\
\bottomrule
\end{longtable}
\endgroup

The column ``Forms'' lists the forms in which the row can be built; ``all
three'' means \texttt{exact}, \texttt{tangent} and \texttt{pwl}. A linear row is
built in its exact form, and a value of \texttt{reliability\_impl} at its cell is
ignored.
$R$ is the \emph{Hoeffding budget}, the sum of the squared widths of the terms of
the damage. $K$ is the CGF of the damage at the cell's parameter
$s = \texttt{s\_chernoff} > 0$. The day increments are
$\delta\mu_k = \sum_{j=1}^M x_{ijk}\,\mu^{\inc}_{jk}$,
$\delta v_k = \sum_{j} x_{ijk}\,v^{\inc}_{jk}$,
$\delta R_k = \sum_{j} x_{ijk}\,(b^{\inc}_{jk})^2$ and
$\delta K_k = \sum_{j} x_{ijk}\,K^{\inc}_{jk}$, with
$K^{\inc}_{jk} = \ln\mathbb{E}\,e^{s X_{jk}}$ the CGF of the increment at $s$
(input \texttt{cgf}).

\subsubsection{Day maps}

Between replacements the damage is $D_k = \sum_t w_t X_t$: the $X_t$ are the
known part, the independent terms of the initial damage or of the replacement
draw, and the mission increments; each weight $w_t \in (0, 1]$ is a product of
factors $1 - \rho$. Under $\ARDone$, write $D^- = G + \Delta$, with $G$ the damage
just after the last intervention and $\Delta$ the weighted increments added
since. $G$ and $\Delta$ are independent by A1, and a repair gives
$D^+ = G + (1-\rho)\,\Delta$. Before the first intervention, $G$ is the part of
the initial damage that the starting anchors describe
(Section~\ref{sec:repair_laws}).

For a damage $D = \sum_t w_t X_t$, write $m(D) = \mathbb{E}D$,
$V(D) = \mathrm{Var}\,D$, $R(D) = \sum_t w_t^2 r_t$ with $r_t$ the squared width
of $X_t$, and $K(D) = \ln\mathbb{E}\,e^{sD}$.

\begin{proposition}[Day maps of the statistics]
\label{prop:rf_maps}
Let $\theta \in \{m, V, R\}$ with factor $a_\theta = 1-\rho$ for $m$ and
$a_\theta = (1-\rho)^2$ for $V$ and $R$. Under A1:
\begin{enumerate}[leftmargin=1.5em, nosep]
  \item a mission day adds the increment's statistic:
        $\theta(D + X) = \theta(D) + \theta(X)$ and $K(D + X) = K(D) + K(X)$;
  \item $\ARDinf$: $\theta\bigl((1-\rho)D\bigr) = a_\theta\,\theta(D)$ and
        $K\bigl((1-\rho)D\bigr) \leq (1-\rho)\,K(D)$;
  \item $\ARDone$: $\theta(D^+) = a_\theta\,\theta(D^-) + (1 - a_\theta)\,\theta(G)$
        and $K(D^+) \leq (1-\rho)\,K(D^-) + \rho\,K(G)$.
\end{enumerate}
\end{proposition}

\begin{proof}
Item 1: the mean and $R$ are additive by definition, the variance and the CGF by
independence. Item 2: $m$ and $V$ scale by $1-\rho$ and $(1-\rho)^2$; every
weight is multiplied by $1-\rho$, so $R$ scales by $(1-\rho)^2$. A CGF is convex
with value 0 at 0, so
$K\bigl((1-\rho)D\bigr) = K_D\bigl((1-\rho)s\bigr) \leq (1-\rho)\,K_D(s) + \rho\,K_D(0)
= (1-\rho)\,K(D)$, with $K_D$ the CGF of $D$ as a function of its argument.
Item 3: by independence, $\theta(D^-) = \theta(G) + \theta(\Delta)$; for $R$ the
weights of the terms of $G$ are unchanged and those of $\Delta$ are multiplied by
$1-\rho$. Hence
\[
  \theta(D^+) = \theta(G) + a_\theta\,\theta(\Delta)
  = \theta(G) + a_\theta\bigl(\theta(D^-) - \theta(G)\bigr)
  = a_\theta\,\theta(D^-) + (1 - a_\theta)\,\theta(G),
\]
and, with item 2 applied to $\Delta$,
\[
  K(D^+) = K(G) + K_\Delta\bigl((1-\rho)s\bigr)
  \leq K(G) + (1-\rho)\bigl(K(D^-) - K(G)\bigr)
  = (1-\rho)\,K(D^-) + \rho\,K(G).
\]
\end{proof}

So each tracked state follows the generic map of Section~\ref{sec:encoding} with
the factor $a = 1-\rho$ for $\mu$ and $K$, and $a = (1-\rho)^2$ for $v$ and $R$.
Its anchor is its value just after the last intervention. The maps of $\mu$,
$v$ and $R$ are exact; the map of $K$ is an upper bound, which is what the
domination argument of Section~\ref{sec:encoding} needs.

\subsubsection{Starting and replacement values}

The initial and replacement values of $\mu$ and $v$ come from the inputs
(Section~\ref{sec:replacement}). Those of $R$ and $K$ are:
\begin{itemize}[leftmargin=1.5em, nosep]
  \item $R_0$ from \texttt{R\_0} and $R^{\new}$ from \texttt{R\_new}, default 0;
  \item $K_0$ from \texttt{cgf\_0}, default $s\,\mu_0$, and $K^{\new}$ from
        \texttt{cgf\_new}, default $s\,\mu^{\new}$.
\end{itemize}
With \texttt{replacement\_as\_new} all replacement values are 0
(Section~\ref{sec:replacement}).

\emph{Why the defaults.} By Jensen's inequality,
$K(D) = \ln\mathbb{E}\,e^{sD} \geq s\,\mathbb{E}D$, with equality exactly when $D$
is almost surely constant. So $s\mu_0$ is the exact CGF of a known initial
damage $\mu_0$, and a lower bound for a random one; a random initial damage needs
its own CGF in \texttt{cgf\_0}. The input checks of
Section~\ref{sec:consistency_checks} require $\texttt{cgf} \geq s\,\mu^{\inc}$ for
every increment, $K_0 \geq s\mu_0$ and $K^{\new} \geq s\mu^{\new}$. A known damage
has width 0, so $R_0 = 0$ is exact for it, and a random one needs \texttt{R\_0}.
The same holds for the replacement draw. A starting value of $K$ below $s\mu_0$ makes the row optimistic. Example:
$\tau = 1$, $\varepsilon = 0.1$, $s = 5$, a known initial damage $\mu_0 = 0.6$,
and one mission that adds 0 or 0.5 with probability $1/2$ each, so
$K^{\inc} = \ln(0.5 + 0.5\,e^{2.5}) = 1.886$. Then
$\Pr(D > \tau) = 0.5$. The row with $K_0 = s\mu_0 = 3$ gives
$4.886 - 5 > \ln 0.1$ and rejects the plan; a start at $K_0 = 0$ would give
$1.886 - 5 = -3.11 \leq \ln 0.1 = -2.30$ and accept it.

\subsubsection{Rows of the dynamics}

The rows are the generic rows of Section~\ref{sec:encoding}, built for $\mu$
and for the one extra state of the bound, with the data below.
$V^{\mathrm{C}} = \varepsilon\tau^2/(1-\varepsilon)$ and
$V^{\mathrm{B}} = \tau(\tau - 2c)/(2\kappa)$ are the variance ranges of the
Cantelli and Bernstein rows, with $\kappa = \ln(1/\varepsilon)$ and $c$ from
Eq.~\eqref{eq:rf_bsupport}. Each natural range is the largest value that the
bound's exact row admits at $\mu \geq 0$. The big-M rows need every initial
value, starting anchor and replacement value inside the natural range of its
state; the input checks of Section~\ref{sec:consistency_checks} require it.

\begingroup\small
\begin{longtable}{@{}llllll@{}}
\toprule
State & Day increment & Factor $a$ & Natural range & Initial value & Replacement value \\
\midrule
$\mu$; anchor $g^{\mu}$ & $\delta\mu_k$ & $1-\rho$ & $\tau$ & $\mu_0$ & $\mu^{\new}$ \\
$v$; anchor $g^{v}$ & $\delta v_k$ & $(1-\rho)^2$ & $V^{\mathrm{C}}$ or $V^{\mathrm{B}}$
  & $v_0$ & $v^{\new}$ \\
$R$; anchor $g^{R}$ & $\delta R_k$ & $(1-\rho)^2$ & $2\tau^2/\kappa$ & $R_0$ & $R^{\new}$ \\
$K$; anchor $g^{K}$ & $\delta K_k$ & $1-\rho$ & $s\tau - \kappa$ & $K_0$ & $K^{\new}$ \\
\bottomrule
\end{longtable}
\endgroup

The anchors exist only under $\ARDone$, and the extra state only for its bound:
$v$ for Cantelli and Bernstein, $R$ for Hoeffding, $K$ for Chernoff.

\subsubsection{Reliability rows}

Write $d_k = \tau - \mu_k$. Let $(m, V, R^{\mathrm{t}}, K^{\mathrm{t}})$ be the
true statistics of $D_k$. By Proposition~\ref{prop:rf_maps} and
Section~\ref{sec:encoding}, each is at most the tracked state of the same name.
Each row below is imposed on every day $k = 1, \ldots, T$.

\textbf{Markov} (\texttt{markov}). By A0 and Markov's inequality,
$\Pr(D_k > \tau) \leq m/\tau \leq \mu_k/\tau$. The row is
\begin{equation}
  \mu_k \leq \varepsilon\,\tau.
  \label{eq:rf_markov}
\end{equation}
It uses the mean only and does not need A1 or A2.

\textbf{Cantelli} (\texttt{cantelli}). Cantelli's inequality gives
$\Pr(D - m \geq a) \leq V/(V + a^2)$ for $a > 0$. The row is
\begin{equation}
  (1 - \varepsilon)\, v_k \leq \varepsilon\,(\tau - \mu_k)^2, \qquad \mu_k \leq \tau,
  \label{eq:rf_cantelli}
\end{equation}
equivalently $\mu_k + \kappa_C\sqrt{v_k} \leq \tau$ with
$\kappa_C = \sqrt{(1-\varepsilon)/\varepsilon}$. \emph{Validity.} If $d_k > 0$,
\[
  \Pr(D_k > \tau) \leq \Pr(D_k - m \geq \tau - m) \leq \frac{V}{V + (\tau - m)^2}
  \leq \frac{v_k}{v_k + d_k^2} \leq \varepsilon,
\]
because $v/(v + d^2)$ increases in $v$ and decreases in $d$, and $\tau - m \geq d_k$.
If $d_k = 0$, the row gives $v_k = 0$, so $V = 0$ and $D_k = m \leq \tau$.

\textbf{Hoeffding} (\texttt{hoeffding}). $D_k - m$ is a sum of independent
centred terms $w_t(X_t - \mathbb{E}X_t)$ with ranges of widths $w_t\sqrt{r_t}$,
and the sum of their squares is $R^{\mathrm{t}}$. Hoeffding's inequality gives
$\Pr(D_k - m \geq a) \leq \exp(-2a^2/R^{\mathrm{t}})$, and
$\exp(-2d^2/R) \leq \varepsilon$ exactly when $d^2 \geq (\kappa/2)\,R$. The row
is
\begin{equation}
  (\tau - \mu_k)^2 \geq \frac{\kappa}{2}\, R_k, \qquad \mu_k \leq \tau.
  \label{eq:rf_hoeffding}
\end{equation}
\emph{Validity.} If $R^{\mathrm{t}} = 0$, every term is constant and
$D_k = m \leq \tau$; this covers $d_k = 0$, where the row gives $R_k = 0$.
Otherwise $d_k > 0$, and
$\Pr(D_k > \tau) \leq \exp\bigl(-2(\tau - m)^2/R^{\mathrm{t}}\bigr)
\leq \exp(-2d_k^2/R_k) \leq \varepsilon$, because $\exp(-2d^2/R)$ increases in
$R$ and decreases in $d$, and $\tau - m \geq d_k$.

\textbf{Bernstein} (\texttt{bernstein}). The \emph{support constant} of the cell
is
\begin{equation}
  b = \max\Bigl(b^0,\; b^{\new},\; \max_{j,k} b^{\inc}_{jk}\Bigr), \qquad
  c = \frac{\kappa\, b}{3},
  \label{eq:rf_bsupport}
\end{equation}
over the missions $j$ and the days $k$ of both phases. It is a constant, not a
decision variable. Every weighted term satisfies
$w_t(X_t - \mathbb{E}X_t) \leq w_t\, b_t \leq b$, whatever the schedule, because
$w_t \leq 1$. Bernstein's inequality for independent terms bounded above by $b$
gives $\Pr(D - m \geq a) \leq \exp\bigl(-a^2/(2(V + b\,a/3))\bigr)$ for $a > 0$;
the right side is at most $\varepsilon$ exactly when
$a^2 - 2c\,a - 2\kappa V \geq 0$. The row is
\begin{equation}
  (\tau - c - \mu_k)^2 \geq c^2 + 2\kappa\, v_k, \qquad \mu_k \leq \tau - c,
  \label{eq:rf_bernstein}
\end{equation}
equivalently $\mu_k + c + \sqrt{c^2 + 2\kappa v_k} \leq \tau$. With $v_k \geq 0$,
the two rows give $d_k \geq 2c$.
\emph{Validity.} $a^2/(2(V + b\,a/3))$ increases in $a > 0$ and decreases in $V$,
and $\tau - m \geq d_k \geq 2c$; so $\Pr(D_k > \tau) \leq \varepsilon$ when
$b > 0$. When $b = 0$, every term is constant and $D_k = m \leq \tau$.

\emph{Remark (the side row).} The quadratic row alone also admits the point
$d_k = 0$, $v_k = 0$: a known damage equal to $\tau$. That point is safe, but it is
cut off from the rest of the feasible set. The side row
$\mu_k \leq \tau - c$ removes it and keeps the branch
$d_k \geq c + \sqrt{c^2 + 2\kappa v_k}$. At $\mu = v = 0$ the row needs
$\tau \geq 2c$; a positive variance needs $\tau > 2\kappa b/3$, the input check
of Section~\ref{sec:consistency_checks}.

\textbf{Chernoff} (\texttt{chernoff}). For $s > 0$, Markov's inequality applied
to $e^{sD}$ gives $\Pr(D_k \geq \tau) \leq \exp(K^{\mathrm{t}} - s\tau)$. The row
is
\begin{equation}
  K_k - s\,\tau \leq \ln\varepsilon.
  \label{eq:rf_chernoff}
\end{equation}
\emph{Validity.} $K^{\mathrm{t}} \leq K_k$, so
$\Pr(D_k > \tau) \leq \exp(K_k - s\tau) \leq \varepsilon$. The row bounds $K$ by
$s\tau - \kappa$, its natural range; $K \geq 0$, because
$K(X) \geq s\,\mathbb{E}X \geq 0$ by Jensen's inequality and A0. A known zero
damage satisfies the row only if $s\tau \geq \kappa$ (input check
$s > \ln(1/\varepsilon)/\tau$, Section~\ref{sec:consistency_checks}). A
replacement day satisfies it only if $K^{\new} \leq s\tau - \kappa$; the input
checks require every replacement state to satisfy the cell's row.

\textbf{Linear forms.} Markov and Chernoff rows are linear: they are built in
their exact form, and a value of \texttt{reliability\_impl} at the cell is
ignored. The Cantelli, Hoeffding and Bernstein rows have the form
$\chi_k \leq f(\mu_k)$ with $f$ convex in $\mu$, and accept the forms
\texttt{exact}, \texttt{tangent} and \texttt{pwl} of
Section~\ref{sec:linear_forms}. Every form keeps the side row.

\begingroup\small
\begin{longtable}{@{}lllll@{}}
\toprule
Bound & $\chi$ & $f(\mu)$ & Side row & $f(0)$ \\
\midrule
Cantelli & $v$ & $\varepsilon(\tau - \mu)^2/(1 - \varepsilon)$ & $\mu \leq \tau$
  & $\varepsilon\tau^2/(1-\varepsilon)$ \\
Hoeffding & $R$ & $2(\tau - \mu)^2/\kappa$ & $\mu \leq \tau$ & $2\tau^2/\kappa$ \\
Bernstein & $v$ & $\bigl((\tau - \mu)^2 - 2c(\tau - \mu)\bigr)/(2\kappa)$
  & $\mu \leq \tau - c$ & $\tau(\tau - 2c)/(2\kappa)$ \\
\bottomrule
\end{longtable}
\endgroup

Each $f$ is a parabola in $\mu$ with a positive leading coefficient, so it is
convex. For Bernstein with $b > 0$, the tangent point must satisfy
$\mu_r < \tau - c$, an input check of Section~\ref{sec:consistency_checks}.
At $\mu_r \geq \tau - c$ the tangent has a nonnegative
slope and lies below $f$, so at every $\mu \leq \tau - c$ it is at most
$f(\tau - c) = -c^2/(2\kappa) < 0$, and no point with $v \geq 0$ satisfies the
row.

\textbf{Optional row.} If \texttt{v\_max\_user} is given for a Cantelli or
Bernstein cell, the row $v_k \leq V^{\mathrm{user}}$ is added on every day.

\textbf{Facts used by the encoding and loop arguments.}
(M1) Every day map is nondecreasing in each state and anchor: its
coefficients $a$ and $1 - a$ lie in $[0, 1]$.
(M2) The true statistics follow the maps with equality for $\mu$, $v$ and $R$,
and with $\leq$ for $K$ (Proposition~\ref{prop:rf_maps}).
(M3) If the tracked states satisfy the row of the cell's bound and the true
statistics are not larger, then $\Pr(D_k > \tau) \leq \varepsilon$ (the validity
paragraphs above). The loop row of Section~\ref{sec:loop} covers $\mu$, the extra
state of the bound and, under $\ARDone$, their anchors.

% ----------------------------------------------------------------
\subsection{Repair Cost Variable}
\label{sec:repair_cost}
% ----------------------------------------------------------------

The variable $z_{i\ell k} \geq 0$ is the degradation removed by a repair of cell
$(i,\ell)$ on day $k$, in mean units, for every model and both laws. The objective
charges $C_{R,i\ell}\, z_{i\ell k}$ (Section~\ref{sec:objective}). The variable and
its rows are built only at cells with $C_{R,i\ell} > 0$, and only on the days whose
cost enters the objective: every day $k = 1, \ldots, T$, except the days of the
transitory phase in the mode \texttt{operating\_average}. The rows are
\begin{equation}
  z_{i\ell k} \;\geq\; \rho_{i\ell}\,\bigl(\mu_{i\ell,k-1} - g^{\mu}_{i\ell,k-1}\bigr)
    - \rho_{i\ell}\,\tau_{i\ell}\,\bigl(1 - x^m_{i\ell k}\bigr),
  \qquad z_{i\ell k} \geq 0,
  \label{eq:a_z}
\end{equation}
with $g^{\mu} \equiv 0$ under $\ARDinf$. On a repair day the first row reads $z_{i\ell
k} \geq \rho\,(\mu_{i\ell,k-1} - g^{\mu}_{i\ell,k-1})$, the removed amount of
Eq.~\eqref{eq:a_psi}. On any other day its right-hand side is at most $0$, because
$0 \leq \mu_{i\ell,k-1} - g^{\mu}_{i\ell,k-1} \leq \mu_{i\ell,k-1} \leq \tau$.

\emph{Replacement days.} A replacement day has $x^m_{i\ell k} = 0$, so $z_{i\ell k} =
0$ at the optimum. A replacement costs $C_{\mathrm{rep},i\ell}$ only. The replaced
part is discarded, so its damage does not change the price of the new part.

\emph{Remark (no upper rows).} The variable $z_{i\ell k}$ appears in no row other
than Eq.~\eqref{eq:a_z}, and in the objective with a positive coefficient. The
minimisation therefore drives it to the largest of its lower bounds:
$z_{i\ell k} = \rho\,(\mu_{i\ell,k-1} - g^{\mu}_{i\ell,k-1})$ on a repair day and
$z_{i\ell k} = 0$ on any other day. The upper rows of a McCormick envelope,
$z \leq \rho\,(\mu - g^{\mu})$ and $z \leq \rho\,\tau\,x^m$, would only remove points
that are never optimal, so they are not built. A solution that is not optimal (for
example, the best one found at the time limit) may carry a larger $z$; its reported
cost is then larger, never smaller.

\begin{proposition}[What $z$ measures]
\label{prop:a_repair_cost}
In every feasible solution, on every repair day $k$ of a cell, under both laws,
\[
  z_{i\ell k} \;\geq\; \rho\,\bigl(\mu_{k-1} - g^{\mu}_{k-1}\bigr)
  \;\geq\; \rho\,\bigl(\zeta^{\mu}_{k-1} - \zeta^{g}_{k-1}\bigr),
\]
where $\zeta^{\mu}$ and $\zeta^{g}$ are the mean and its anchor on the reference
trajectory (Section~\ref{sec:encoding}). When the starting anchor equals the true
mean level after the last intervention before the horizon, the right-hand side is
the expected degradation that the repair removes.
\end{proposition}
\begin{proof}
The first inequality is Eq.~\eqref{eq:a_z} with $x^m_{i\ell k} = 1$. Under $\ARDinf$,
$g^{\mu} = \zeta^{g} = 0$, and $\mu_{k-1} \geq \zeta^{\mu}_{k-1}$ by
Proposition~\ref{prop:a_domination}. Under $\ARDone$, let $t < k$ be the last
intervention day before $k$, or $t = 0$ if there is none. Then $g^{\mu}_{k-1} =
g^{\mu}_t$ and $\zeta^{g}_{k-1} = \zeta^{g}_t$.
\begin{itemize}[nosep,leftmargin=1.5em]
  \item If $t = 0$: $g^{\mu}_0 = \zeta^{g}_0$ (the starting anchor), so
        $g^{\mu}_{k-1} - \zeta^{g}_{k-1} = 0 \leq \mu_{k-1} - \zeta^{\mu}_{k-1}$.
  \item If $t \geq 1$: $g^{\mu}_t = \mu_t$ and $\zeta^{g}_t = \zeta^{\mu}_t$
        (Eq.~\eqref{eq:a_anchor}). On the days $j = t+1, \ldots, k-1$ there is no
        intervention: $\mu_j \geq \mu_{j-1} + \delta\mu_j$ and $\zeta^{\mu}_j =
        \zeta^{\mu}_{j-1} + \delta\mu_j$. So $\mu_j - \zeta^{\mu}_j$ does not
        decrease, and $\mu_{k-1} - \zeta^{\mu}_{k-1} \geq \mu_t - \zeta^{\mu}_t =
        g^{\mu}_{k-1} - \zeta^{g}_{k-1}$.
\end{itemize}
Rearranging gives $\mu_{k-1} - g^{\mu}_{k-1} \geq \zeta^{\mu}_{k-1} -
\zeta^{g}_{k-1}$. The mean recursion is linear (Section~\ref{sec:repair_laws}); with
an exact starting anchor, $\zeta^{\mu}$ and $\zeta^{g}$ are the true means, so
$\rho\,(\zeta^{\mu}_{k-1} - \zeta^{g}_{k-1}) = \mathbb{E}[D^- - D^+]$.
\end{proof}

A starting anchor above the true one, which the default may be
(Section~\ref{sec:repair_laws}), models a weaker repair. For the first repair before
any intervention of the horizon, the charged removal then follows the starting
anchor and may be below the true removal. From the first intervention on, the
reference gap $\zeta^{\mu} - \zeta^{g}$ equals the true gap, the sum of the mean
increments since that intervention, so the bound holds again.

Under both laws, then, the reported repair cost $\sum C_R\, z$ is an upper bound on
the expected repair cost of the schedule, with the exception above for a starting
anchor above the true one. Under $\ARDone$ the tracked mean and the
tracked anchor are both upper bounds, and their difference is still an upper bound
on the reference difference: the anchor carries the slack that the mean had at the
last intervention, and the slack of the mean does not decrease until the next one.
The bound is exact when the tracked mean has no slack.

% ----------------------------------------------------------------
\subsection{Damage Regularisation}
\label{sec:damage_reg}
% ----------------------------------------------------------------

The scalar $u \geq 0$ measures the largest damage of one train on one day, over the
whole horizon. The objective charges $C_D\,u$ (Section~\ref{sec:objective}). The
variable $u$ and its rows are built only when $C_D > 0$. With $C_D = 0$ they are not
built, and the input \texttt{penalty\_type} is ignored. The input check requires
$C_D \geq 0$ (Section~\ref{sec:consistency_checks}): a negative $C_D$ would reward
damage.

The input \texttt{penalty\_type} selects one of two forms. Both have one row per train
and day.
\begin{itemize}[leftmargin=1.5em]
  \item \texttt{inf\_norm} (default), a linear row:
        \begin{equation}
          u \;\geq\; \sum_{\ell=1}^{L} \mu_{i\ell k} \qquad \forall\, i, k .
          \label{eq:a_u_inf}
        \end{equation}
        At the optimum, $u = \max_{i,k} \sum_{\ell} \mu_{i\ell k}$.
  \item \texttt{quadratic}, a convex quadratic row:
        \begin{equation}
          u \;\geq\; \sum_{\ell=1}^{L} \mu_{i\ell k}^2 \qquad \forall\, i, k .
          \label{eq:a_u_quad}
        \end{equation}
        The row is a rotated second-order cone, $u \cdot 1 \geq \|\mu_{i\cdot k}\|_2^2$,
        which Gurobi solves without the parameter \texttt{NonConvex}. It requires
        \texttt{reliability\_impl = "exact"} at every cell
        (Section~\ref{sec:consistency_checks}): the linear forms are chosen for a
        model without quadratic rows, and this row is quadratic.
\end{itemize}

\emph{Role.} The term prices the worst train-level damage, so the minimisation
spreads damage across the fleet. It also removes slack from the tracked means where
$u$ is binding (Section~\ref{sec:encoding}). Since $\mu_{i\ell k} \geq
\zeta^{\mu}_{i\ell k} \geq 0$ (Proposition~\ref{prop:a_domination}), both forms of $u$
are upper bounds on the same quantity computed from the reference values.

% ----------------------------------------------------------------
\subsection{Loop Constraint}
\label{sec:loop}
% ----------------------------------------------------------------

The schedule of the operating phase is meant to repeat without end. One hard row per
tracked quantity and per anchor makes every repetition safe:
\begin{equation}
  s_{i\ell T} \;\leq\; s_{i\ell H_1}
  \qquad \text{for every cell } (i,\ell) \text{ and every coordinate } s
  \text{ of its full tracked state.}
  \label{eq:a_loop}
\end{equation}
The rows cover every tracked quantity of the cell (the mean, and the variance, $R$
or $K$ where the model tracks them) and, under $\ARDone$, every anchor. There is no
soft loop penalty in the objective.

\textbf{Endless schedule.} Repeat the operating phase without end: day $T + n H_2 +
h$, for $n \geq 0$ and $1 \leq h \leq H_2$, carries the binaries of day $H_1 + h$.
It also carries the increments of day $H_1 + h$, because the operating profile is
read cyclically with period $H_2$ (Section~\ref{sec:input_increments}). So every
repetition applies the same day maps $\varphi_{H_1+1}, \ldots, \varphi_T$. Write
$\Phi_h = \varphi_{H_1+h} \circ \cdots \circ \varphi_{H_1+1}$ for the first $h$
days of the operating phase and $\Phi = \Phi_{H_2}$ for the whole phase.

\begin{theorem}[Perpetual safety]
\label{thm:a_loop}
Let a solution satisfy every row of the problem. Then, for every cell and every day
$k \geq 1$ of the endless schedule, the true statistics $\theta_k$ and the reference
values $\zeta_k$ are at most a tracked value of the horizon that lies in
$\mathcal{C}_{i\ell}$. Therefore $\zeta_k \in \mathcal{C}_{i\ell}$ and
$\Pr(D_{i\ell k} > \tau_{i\ell}) \leq \varepsilon_{i\ell}$ on every day $k \geq 1$.
\end{theorem}
\begin{proof}
Fix a cell. Let $\xi_k$ be its tracked full state, and $\theta_k$ and $\zeta_k$ the
true statistics and the reference trajectory of the endless schedule. By (M2) and
Lemma~\ref{lem:a_monotone}, the induction of Proposition~\ref{prop:a_domination}
gives $\theta_k \leq \zeta_k$ on every day $k \geq 1$.
\begin{itemize}[nosep,leftmargin=1.5em]
  \item Days $k \leq T$: $\zeta_k \leq \xi_k$ (Proposition~\ref{prop:a_domination})
        and $\xi_k \in \mathcal{C}_{i\ell}$ (consequence on the horizon,
        Section~\ref{sec:encoding}).
  \item By Proposition~\ref{prop:a_domination} and Lemma~\ref{lem:a_monotone},
        $\xi_{H_1+h} \geq \varphi_{H_1+h}(\xi_{H_1+h-1}) \geq \cdots \geq
        \Phi_h(\xi_{H_1})$ for $h = 1, \ldots, H_2$. With Eq.~\eqref{eq:a_loop},
        $\Phi(\xi_{H_1}) \leq \xi_T \leq \xi_{H_1}$.
  \item Let $\Phi^n(\xi_{H_1}) \leq \xi_{H_1}$ (true for $n = 0$). Then
        $\Phi^{n+1}(\xi_{H_1}) = \Phi(\Phi^n(\xi_{H_1})) \leq \Phi(\xi_{H_1}) \leq
        \xi_{H_1}$ by Lemma~\ref{lem:a_monotone}. So $\Phi^n(\xi_{H_1}) \leq
        \xi_{H_1}$ for every $n \geq 0$.
  \item Day $k = H_1 + n H_2 + h$ with $n \geq 0$, $1 \leq h \leq H_2$:
        \[
          \zeta_k = \Phi_h\bigl(\Phi^n(\zeta_{H_1})\bigr)
          \leq \Phi_h\bigl(\Phi^n(\xi_{H_1})\bigr)
          \leq \Phi_h(\xi_{H_1})
          \leq \xi_{H_1+h} \in \mathcal{C}_{i\ell}.
        \]
\end{itemize}
Therefore, on every day $k \geq 1$, $\theta_k \leq \zeta_k \leq s$ for a tracked value
$s \in \mathcal{C}_{i\ell}$. Since $\mathcal{C}_{i\ell}$ is a down-set,
$\zeta_k \in \mathcal{C}_{i\ell}$, and (M3) gives $\Pr(D_{i\ell k} > \tau_{i\ell})
\leq \varepsilon_{i\ell}$.
\end{proof}

The theorem holds for every model and both repair laws, because it uses only
the facts (M1)--(M3) and Lemma~\ref{lem:a_monotone}. It holds in both encodings and
in every reliability form. It needs no relation between $H_1$ and $H_2$: the
transitory phase may be shorter or longer than the operating phase, because
$H_1$ only marks the day at which the repeated block starts. Slack in the tracked
values does not weaken it: the proof compares the reference trajectory with the
tracked values, whatever slack they carry.

\emph{Remark (why the loop row covers the full state).} The third step of the proof
needs $\Phi(\xi_{H_1}) \leq \xi_{H_1}$ in every coordinate, anchors included.
Example: law $\ARDone$, $\rho = 0.5$, the first phase ends at $(\mu, g) = (0.4,
0.1)$, and the operating phase is a mission ($+0.1$), a repair, a mission
($+0.05$). The phase ends at $(0.35, 0.3)$, since the repair gives $0.5 \cdot 0.5 +
0.5 \cdot 0.1 = 0.3$. The row on the mean holds ($0.35 \leq 0.4$); the row on the
anchor fails ($0.3 > 0.1$). Each repetition raises the state after the repair by
$0.075$, so the mean grows without bound: the third repetition already peaks at
$0.525$, above the checked peak $0.5$. A loop row on the mean alone would accept
this unsafe schedule.

\emph{Remark (no loop row on a scalar bound).} The loop row must compare the tracked
quantities one by one. A single comparison of a scalar function of them, such as the
failure-probability bound $U(\mu, v)$, does not chain: $U$ takes the same value at
many states, and its level sets are not invariant under the day maps. Example
(Cantelli, $U = v / (v + (\tau - \mu)^2)$, $\tau = 10$, $\rho = 0.5$, law
$\ARDinf$; the operating phase is one repair day followed by mission days that add
$3$ to the mean and $3.99$ to the variance in total): from $(\mu, v) = (4, 9)$,
$U = 0.2000$ at the end of the first phase and $U = 0.1997$ at the end of the
operating phase, so a scalar row would accept. The next repetition gives $U =
0.2151 > 0.2$, and $U$ tends to $0.2495$. Eq.~\eqref{eq:a_loop} rejects the
schedule, since the mean rises from $4$ to $5$. The loop row states an order between
tracked quantities; it is not a stochastic order between damage laws.

\begin{proposition}[Replacement under $\ARDone$]
\label{prop:a_ard1_replacement}
Let a cell follow $\ARDone$ with $\rho_{i\ell} < 1$, and let its train fly, in the
operating phase, at least one mission with a positive mean increment for the cell.
Then every feasible solution has at least one replacement of the cell in the
operating phase.
\end{proposition}
\begin{proof}
Let the cell have no replacement in the operating phase. By the anchor rows,
$g_k = g_{k-1}$ on a day without intervention and $g_k = \mu_k$ on a repair day.
Every day satisfies $\mu_k \geq g_k$: on an intervention day with equality, and on
other days because $\mu$ does not decrease and $g$ stays constant (with $g_0 \leq
\mu_0$, Section~\ref{sec:repair_laws}). On a repair day $k$ of the operating phase,
\[
  g_k = \mu_k \geq (1-\rho)\,\mu_{k-1} + \rho\, g_{k-1}
  = g_{k-1} + (1-\rho)\,(\mu_{k-1} - g_{k-1}).
\]
So $g_T \geq g_{H_1} + (1-\rho) \sum (\mu_{k-1} - g_{k-1})$, the sum over the
repair days of the operating phase, each term non-negative. The anchor row $g_T
\leq g_{H_1}$ and $\rho < 1$ force every term to $0$. Then every repair day gives
$\mu_k \geq \mu_{k-1}$, every other day gives $\mu_k \geq \mu_{k-1} + \delta\mu_k$,
and
\[
  \mu_T \geq \mu_{H_1} + \textstyle\sum_{k = H_1+1}^{T} \delta\mu_k > \mu_{H_1},
\]
which violates the row $\mu_T \leq \mu_{H_1}$. Therefore the cell is replaced at
least once in the operating phase.
\end{proof}

The requirement is not an artefact of the rows. Without a replacement, each repair
leaves the share $1-\rho$ of the damage since the last intervention in place, so the
reference mean grows without bound over the repetitions, as in the first remark
above. With $\rho = 1$ a repair returns the state to the anchor, and no replacement
is needed. The Wiener model needs a replacement in the operating phase for another
reason, its variance (Section~\ref{sec:dyn_wiener}).

% ================================================================
\section{Full Optimization Problem}
\label{sec:full_problem}
% ================================================================

This section collects the whole problem in one place. Each row is stated once, in the
section that the reference names; this section does not repeat the rows.

\textbf{Objective.} Minimise Eq.~\eqref{eq:a_objective}, the mode chosen by
\texttt{objective\_mode}, with the day cost Eq.~\eqref{eq:a_daycost}.

\textbf{Fleet rows.} For every train $i$ and day $k$: one activity,
Eq.~\eqref{eq:a_one_activity}. For every mission $j \geq 1$ and day $k$: coverage,
Eq.~\eqref{eq:a_coverage}. For every cell and day: depot-day compatibility,
Eq.~\eqref{eq:a_depot}.

\textbf{Rows of a cell.} For every cell $(i,\ell)$, with the model, repair law,
tail bound, encoding and reliability form of the cell:
\begin{enumerate}[nosep,leftmargin=2em]
  \item for every tracked quantity $s$ of Table~\ref{tab:full_cells}: the generic rows
        Eqs.~\eqref{eq:a_bigm_carry}--\eqref{eq:a_bigm_repl} (\texttt{encoding = "bigm"})
        or Eqs.~\eqref{eq:a_ind_sum}--\eqref{eq:a_ind_state}
        (\texttt{encoding = "indicator"}), with the repair map
        Eq.~\eqref{eq:a_psi} and the factor $a$ of the table; the Wiener variance may
        use the single row Eq.~\eqref{eq:w_vrow};
  \item under $\ARDone$, for every tracked quantity except the Wiener variance: the
        anchor rows Eq.~\eqref{eq:a_bigm_anchor} or Eq.~\eqref{eq:a_ind_anchor};
  \item the reliability row of the table, on every day $k = 1, \ldots, T$, in the form
        \texttt{exact}, \texttt{tangent} (Eq.~\eqref{eq:a_tangent}) or \texttt{pwl}
        (Eqs.~\eqref{eq:a_pwl_sum}--\eqref{eq:a_pwl_cap}) for a quadratic row, and in
        its exact form for a linear row;
  \item the loop rows Eq.~\eqref{eq:a_loop}, one per tracked quantity and per anchor;
  \item where $C_{R,i\ell} > 0$: the repair-cost rows Eq.~\eqref{eq:a_z}.
\end{enumerate}

\textbf{Damage regularisation.} Where $C_D > 0$: Eq.~\eqref{eq:a_u_inf}
(\texttt{inf\_norm}) or Eq.~\eqref{eq:a_u_quad} (\texttt{quadratic}), for every train
and day.

\textbf{Domains.} The variables and their domains are those of
Section~\ref{sec:variables}; a variable that no row uses is not built.

\begingroup\small
\begin{longtable}{@{}L{3.1cm}L{4.4cm}L{3.0cm}L{4.0cm}@{}}
\caption{Tracked quantities and reliability row of each kind of cell. The factor $a$
is the repair factor of Eq.~\eqref{eq:a_psi}; $M_s$ is the natural range of
Section~\ref{sec:encoding}.}
\label{tab:full_cells} \\
\toprule
Cell & Tracked quantities (factor $a$) & Natural ranges $M_s$ & Reliability row \\
\midrule
\endfirsthead
\toprule
Cell & Tracked quantities (factor $a$) & Natural ranges $M_s$ & Reliability row \\
\midrule
\endhead
Gaussian & $\mu$ ($1-\rho$), $v$ ($(1-\rho)^2$) & $\tau$, $V_{\max}$ &
  Eq.~\eqref{eq:g_row}, quadratic \\
Inverse Gaussian & $\mu$ ($1-\rho$) & $\tau$ & Eq.~\eqref{eq:ig_row}, linear \\
Wiener & $\mu$ ($1-\rho$), $v$ ($1$, no anchor) & $\tau$, $V_{\max}$ &
  Eq.~\eqref{eq:g_row}, quadratic \\
Gamma & $\mu$ ($1-\rho$) & $\tau$ & Eq.~\eqref{eq:gm_row}, linear \\
Rainflow, Markov & $\mu$ ($1-\rho$) & $\tau$ & Eq.~\eqref{eq:rf_markov}, linear \\
Rainflow, Cantelli & $\mu$ ($1-\rho$), $v$ ($(1-\rho)^2$) & $\tau$, $V^{\mathrm{C}}$ &
  Eq.~\eqref{eq:rf_cantelli}, quadratic \\
Rainflow, Hoeffding & $\mu$ ($1-\rho$), $R$ ($(1-\rho)^2$) & $\tau$, $2\tau^2/\kappa$ &
  Eq.~\eqref{eq:rf_hoeffding}, quadratic \\
Rainflow, Bernstein & $\mu$ ($1-\rho$), $v$ ($(1-\rho)^2$) & $\tau$, $V^{\mathrm{B}}$ &
  Eq.~\eqref{eq:rf_bernstein}, quadratic \\
Rainflow, Chernoff & $\mu$ ($1-\rho$), $K$ ($1-\rho$) & $\tau$, $s\,\tau - \kappa$ &
  Eq.~\eqref{eq:rf_chernoff}, linear \\
\bottomrule
\end{longtable}
\endgroup

\textbf{Problem class.} The problem is a mixed-integer linear program when every
quadratic reliability row uses the \texttt{tangent} or \texttt{pwl} form and the
damage penalty is \texttt{inf\_norm} or absent. A quadratic row in the \texttt{exact}
form makes it a nonconvex mixed-integer quadratically constrained program, which
Gurobi solves with \texttt{NonConvex} $= 2$. The \texttt{quadratic} damage penalty
alone makes it a convex one (Section~\ref{sec:damage_reg}).

\section{Output Specification}\label{sec:output}

The function \texttt{solve()} returns one result dictionary.
If its argument \texttt{results\_path} is not \texttt{None}, the function also writes the same content to a file.
Every call solves one problem, with the horizon $H$ of the input; a solve for several horizons calls \texttt{solve()} once per horizon (Section~\ref{sec:sweeps}).

\subsection{Result Dictionary}

The array values are \texttt{numpy.ndarray} objects.
The result holds no Gurobi object.
In the column \emph{Shape}, a dash marks a scalar or a string.

\begingroup\small
\begin{longtable}{@{}p{3.3cm}p{2.9cm}p{8.0cm}@{}}
\toprule
Key & Shape & Description \\
\midrule
\endhead
\multicolumn{3}{@{}l}{\emph{Status and quality}} \\
\typerow{\texttt{status}}{---}{Status of the solve: \texttt{optimal}, \texttt{time\_limit}, \texttt{suboptimal}, \texttt{interrupted}, \texttt{solution\_limit}, \texttt{node\_limit}, \texttt{iteration\_limit}, \texttt{infeasible}, \texttt{inf\_or\_unbounded}, \texttt{unbounded}, or \texttt{gurobi\_status\_}$n$ for any other Gurobi status code $n$}
\typerow{\texttt{objective}}{--- or \texttt{null}}{Objective value of the best solution found; \texttt{null} if there is none}
\typerow{\texttt{mip\_gap}}{--- or \texttt{null}}{Relative MIP gap of the best solution at the end of the solve; \texttt{null} if there is no solution. The input key \texttt{mip\_gap} is the gap that was requested}
\typerow{\texttt{bound}}{--- or \texttt{null}}{Best bound of Gurobi, if it is available}
\midrule
\multicolumn{3}{@{}l}{\emph{Problem}} \\
\typerow{\texttt{F}, \texttt{M}, \texttt{L}}{---}{Numbers of trains, missions and components}
\typerow{\texttt{H}}{integer or $[H_1, H_2]$}{Horizon as written in the input}
\typerow{\texttt{H1}, \texttt{H2}, \texttt{T}}{---}{Lengths $H_1$ and $H_2$ of the two phases, and $T = H_1 + H_2$}
\typerow{\texttt{model}}{$F \times L$ list of strings}{Degradation model of each cell, as in the input}
\typerow{\texttt{component\_names}}{list of $L$ strings}{Names of the components}
\typerow{\texttt{tau}, \texttt{mu\_0}, \texttt{v\_0}}{$F \times L$}{Threshold, initial mean and initial variance of each cell, after broadcasting}
\typerow{\texttt{bound\_method}, \texttt{repair\_model}, \texttt{reliability\_impl}, \texttt{encoding}}{string, or $F \times L$ list of strings}{The selector of each cell after broadcasting: one string if all cells agree, else a list. \texttt{bound\_method} is \texttt{null} at a cell that is not rainflow}
\typerow{\texttt{objective\_mode}, \texttt{evaluation\_horizon}}{---}{The objective mode, and $E$ if it was given}
\midrule
\multicolumn{3}{@{}l}{\emph{Solution (\texttt{null} if no solution was found)}} \\
\typerow{\texttt{x}}{$F \times (M{+}1) \times T$}{Assignment: $x_{ijk} = 1$ if train $i$ flies mission $j \ge 1$ on day $k$, and $x_{i0k} = 1$ if it is in the depot. A train with $x_{ijk} = 0$ for every $j \in \{0,\ldots,M\}$ is idle}
\typerow{\texttt{m}}{$F \times L \times T$}{Repair: $x^m_{i\ell k}$}
\typerow{\texttt{r}}{$F \times L \times T$}{Replacement: $x^r_{i\ell k}$; zero if \texttt{allow\_replacement} is \texttt{false}}
\typerow{\texttt{idle}}{$F \times L \times T$}{$\text{idle}_{i\ell k} = x_{i0k} - x^m_{i\ell k} - x^r_{i\ell k}$: component $\ell$ is in the depot without an intervention}
\typerow{\texttt{mu}}{$F \times L \times T$}{Tracked mean degradation}
\typerow{\texttt{v}}{$F \times L \times T$}{Tracked variance. \texttt{NaN} in memory and \texttt{null} in a file at cells that track no variance (IG, Ga, RF with Mk, Ho or Ch)}
\typerow{\texttt{z}}{$F \times L \times T$}{Repair-cost variable. \texttt{NaN} in memory and \texttt{null} in a file at cells with $C_{R,i\ell} = 0$, where the variable is not built. On the transitory days of the mode \texttt{operating\_average} the value is computed after the solve (below)}
\typerow{\texttt{R}, \texttt{K}}{$F \times L \times T$}{Tracked Hoeffding budget $R$ and Chernoff quantity $K$. \texttt{NaN} in memory and \texttt{null} in a file at the cells that do not track the quantity: $R$ is tracked at RF (Ho) cells, $K$ at RF (Ch) cells}
\typerow{\texttt{g\_mu}, \texttt{g\_v}, \texttt{g\_R}, \texttt{g\_K}}{$F \times L \times T$}{Tracked anchors $g^{\mu}_{i\ell k}$, $g^{v}_{i\ell k}$, $g^{R}_{i\ell k}$ and $g^{K}_{i\ell k}$: the value of the state after the last intervention up to day $k$. \texttt{NaN} in memory and \texttt{null} in a file at the cells that do not track the anchor: only cells with \texttt{repair\_model} \texttt{ard1} track one, \texttt{g\_mu} for every model, \texttt{g\_v} at G and RF (Ca, Be) cells, \texttt{g\_R} at RF (Ho) cells, \texttt{g\_K} at RF (Ch) cells}
\typerow{\texttt{u}}{---}{Damage regularisation value (below)}
\midrule
\multicolumn{3}{@{}l}{\emph{Cost report (\texttt{null} if no solution was found)}} \\
\typerow{\texttt{step\_costs}}{$T$}{Cost $c_k$ of each day, Equation~\eqref{eq:c_step}}
\typerow{\texttt{J\_initialization}, \texttt{J\_op}, \texttt{J\_op\_average}}{---}{$J_{\mathrm{init}}$, $J_{\mathrm{op}}$ and $J_{\mathrm{op}}/H_2$, Equation~\eqref{eq:c_phases}}
\typerow{\texttt{peak\_damage}, \texttt{damage\_penalty}}{---}{$u$ and $C_D\,u$}
\typerow{\texttt{operating\_objective}, \texttt{J\_total}}{---}{$J_{\mathrm{op}}/H_2 + C_D\,u$ and $J_{\mathrm{init}} + J_{\mathrm{op}} + C_D\,u$}
\typerow{\texttt{projected\_evaluation\_cost}, \texttt{evaluation\_objective}}{--- or \texttt{null}}{$J_{\mathrm{init}} + (E - H_1)\,J_{\mathrm{op}}/H_2$ and the same plus $C_D\,u$; \texttt{null} if \texttt{evaluation\_horizon} is not given}
\typerow{\texttt{component\_costs}}{mapping}{Cost coefficients used: \texttt{C\_M} and \texttt{C\_D} (scalars), \texttt{C\_R} and \texttt{C\_rep} ($F \times L$)}
\midrule
\multicolumn{3}{@{}l}{\emph{Reports}} \\
\typerow{\texttt{performance}}{mapping}{Sizes and times of the solve (Section~\ref{sec:output_reports})}
\typerow{\texttt{optimization\_progress}}{list of mappings}{Progress of the incumbent and the bound (Section~\ref{sec:output_reports})}
\typerow{\texttt{warm\_start}, \texttt{relaxation\_warm\_start}}{mappings}{What each warm start did (Section~\ref{sec:output_reports})}
\bottomrule
\end{longtable}
\endgroup

\subsection{Rules for the Values}

\textbf{Solutions and status.}
The solution arrays are present whenever Gurobi has found at least one solution, and then \texttt{status} tells how the solve ended.
For example, a solution found before the time limit is reported with the status \texttt{time\_limit}.
If no solution was found, the solution arrays and the cost report are \texttt{null}, and \texttt{objective} and \texttt{mip\_gap} are \texttt{null}.
The four reports of the last block of the table are present in both cases.

\textbf{Tracked values.}
The arrays \texttt{mu}, \texttt{v}, \texttt{R}, \texttt{K}, \texttt{g\_mu}, \texttt{g\_v}, \texttt{g\_R} and \texttt{g\_K} hold the values of the tracked variables of the model, and \texttt{z} holds the repair-cost variable.
The result carries every tracked state and anchor, so that a validator can check every loop row of Section~\ref{sec:loop}.
The dynamics rows are one-sided (Section~\ref{sec:encoding}), so these values, and the costs computed from them, are upper bounds, not exact values.

\textbf{Repair cost on transitory days.}
In the mode \texttt{operating\_average}, the variable $z$ is not built on the days of the transitory phase, because their cost does not enter the objective (Section~\ref{sec:repair_cost}).
For these days, at a cell with $C_{R,i\ell} > 0$, the result holds the value computed after the solve from the arrays \texttt{mu} and \texttt{g\_mu}:
\begin{equation}\label{eq:c_zpost}
  z_{i\ell k} = \begin{cases}
    \rho_{i\ell}\,\bigl(\mu_{i\ell,k-1} - g^{\mu}_{i\ell,k-1}\bigr) & \text{if } x^m_{i\ell k} = 1,\\
    0 & \text{otherwise,}
  \end{cases}
  \qquad k = 1, \ldots, H_1,
\end{equation}
with $\mu_{i\ell 0} = \mu_{0,i\ell}$, $g^{\mu}_{i\ell 0}$ the starting anchor, and $g^{\mu} \equiv 0$ under \texttt{ardinf}.
It is the amount that the repair removes, in mean units, computed on the tracked values.
The cost report uses these values.

\textbf{Damage regularisation.}
The value $u$ is recomputed from the arrays \texttt{mu},
\begin{equation}\label{eq:c_u}
  u = \max_{i,k} \sum_{\ell=1}^{L} \mu_{i\ell k} \quad (\texttt{inf\_norm}),
  \qquad
  u = \max_{i,k} \sum_{\ell=1}^{L} \mu_{i\ell k}^2 \quad (\texttt{quadratic}).
\end{equation}
It is reported also when $C_D = 0$, where the variable $u$ is not built.

\textbf{Depot, idle and intervention.}
A row of zeros $x_{i1k} = \cdots = x_{iMk} = x_{i0k} = 0$ is an idle day of train $i$: no mission and no depot.
On a depot day ($x_{i0k} = 1$) each component is repaired, replaced or in the depot without an intervention; the last case is \texttt{idle}.
Hence $x^m_{i\ell k} + x^r_{i\ell k} + \text{idle}_{i\ell k} = x_{i0k}$.

\textbf{Cost report.}
The cost of day $k$ and the phase costs are
\begin{align}
  c_k &= C_M \sum_{i=1}^{F} x_{i0k}
       + \sum_{i=1}^{F}\sum_{\ell=1}^{L} C_{R,i\ell}\, z_{i\ell k}
       + \sum_{i=1}^{F}\sum_{\ell=1}^{L} C_{\mathrm{rep},i\ell}\, x^r_{i\ell k}, \label{eq:c_step}\\
  J_{\mathrm{init}} &= \sum_{k=1}^{H_1} c_k, \qquad
  J_{\mathrm{op}} = \sum_{k=H_1+1}^{T} c_k. \label{eq:c_phases}
\end{align}
The value \texttt{objective} is the value of the model objective.
Up to the tolerance of the solver, it equals \texttt{J\_total} if \texttt{objective\_mode} is \texttt{total}, \texttt{operating\_objective} if it is \texttt{operating\_average}, and \texttt{evaluation\_objective} if it is \texttt{evaluation\_total}.

\subsection{Reports}\label{sec:output_reports}

\textbf{\texttt{performance}.}
The mapping has these entries; an entry that Gurobi does not give is \texttt{null}.

\begingroup\small
\begin{longtable}{@{}p{3.4cm}p{10.7cm}@{}}
\toprule
Group & Entries \\
\midrule
\endhead
\rejrow{Size of the model}{\texttt{gurobi\_version}, \texttt{variables}, \texttt{continuous\_variables}, \texttt{integer\_variables}, \texttt{binary\_variables}, \texttt{linear\_constraints}, \texttt{general\_constraints}, \texttt{indicator\_constraints}, \texttt{quadratic\_constraints}, \texttt{nonzeros}}
\rejrow{Work of Gurobi}{\texttt{solutions\_found}, \texttt{gurobi\_runtime\_seconds}, \texttt{branch\_and\_bound\_nodes}, \texttt{simplex\_iterations}, \texttt{barrier\_iterations}, \texttt{work\_units}}
\rejrow{Wall-clock times, in seconds}{\texttt{model\_construction\_seconds}, \texttt{optimizer\_call\_seconds}, \texttt{solution\_extraction\_seconds}, \texttt{backend\_wall\_seconds}}
\rejrow{Settings}{\texttt{requested\_time\_limit\_seconds}, \texttt{requested\_mip\_gap}, \texttt{threads}, \texttt{seed}}
\bottomrule
\end{longtable}
\endgroup

\textbf{\texttt{optimization\_progress}.}
A list of records with the entries \texttt{runtime\_seconds}, \texttt{incumbent}, \texttt{best\_bound}, \texttt{relative\_gap}, \texttt{nodes} and \texttt{solutions}.
The incumbent, the bound and the gap are \texttt{null} while Gurobi has none.
A record is made at the first progress report of Gurobi, then at least \texttt{progress\_interval\_seconds} after the previous record, and whenever the number of solutions changes; the last record is the final state of the solve.

\textbf{\texttt{warm\_start}.}
A mapping with the entries \texttt{applied} (\texttt{true} when at least one value was set), \texttt{values} (the number of values set), \texttt{mode} (\texttt{"start"}, \texttt{"hint"} or \texttt{"both"}, and \texttt{null} without a source), \texttt{source\_horizon} (the pair $H_1$, $H_2$ of the source, or \texttt{null}) and \texttt{completion}.
The entry \texttt{completion} is a mapping with \texttt{requested}, \texttt{status}, \texttt{fixed\_integer\_variables}, \texttt{values\_copied}, \texttt{runtime\_seconds} and \texttt{iis\_constraints}.
The value of \texttt{status} is \texttt{"not\_requested"}, \texttt{"completed"}, \texttt{"infeasible"}, \texttt{"infeasible\_or\_unbounded"}, \texttt{"time\_limit"}, or \texttt{"status\_}$q$\texttt{"} with the Gurobi status code $q$.
The entry \texttt{iis\_constraints} is a list of at most $50$ names of rows of an irreducible infeasible subset; it is empty unless the status is \texttt{"infeasible"}.
Section~\ref{sec:warm_starts} defines the mechanisms.

\textbf{\texttt{relaxation\_warm\_start}.}
A mapping with the entries \texttt{enabled}; \texttt{status} (\texttt{"disabled"} when the option is false, \texttt{"skipped"} when it is true and a start or a hint was applied, otherwise the Gurobi status of the relaxed copy as a string); \texttt{relaxation\_objective}; \texttt{values\_copied}; and \texttt{runtime\_seconds}.

\subsection{Result Files}

The function \texttt{solve} writes the result to the file \texttt{results\_path}.
\begin{itemize}[leftmargin=1.5em]
  \item The default value of \texttt{results\_path} is \texttt{"output.yaml"}. The value \texttt{None} writes no file and only returns the dictionary.
  \item A path without extension gets the extension \texttt{.yaml}. The extensions \texttt{.yaml} and \texttt{.yml} give YAML, \texttt{.json} gives JSON, and \texttt{.h5} and \texttt{.hdf5} give HDF5. A path with any other extension raises a \texttt{ValueError}.
  \item The extension, and the directory of the path, are checked before the model is solved, so that a long solve is not lost. The directory must exist and be writable.
  \item YAML and JSON files hold plain data: lists, numbers, strings and \texttt{null}. \texttt{NaN} is written as \texttt{null}. The files load with \texttt{yaml.safe\_load} and \texttt{json.load}.
  \item A key whose value is \texttt{None} is not written, except \texttt{objective}, which is written as \texttt{null}.
  \item The layout of an HDF5 file is that of Section~\ref{sec:schema_hdf5}.
\end{itemize}

% ================================================================
\section{Horizon Sweeps}
\label{sec:sweeps}
% ================================================================

A single call of \texttt{solve()} uses one horizon pair $H = [H_1, H_2]$ (Section~\ref{sec:notation}).
A \emph{horizon sweep} solves one scenario for several pairs and compares the solved cases.
Two functions implement the sweeps: \texttt{sweep\_operating\_horizons} varies $H_2$ for one $H_1$, and
\texttt{sweep\_horizon\_grid} varies $H_1$ and $H_2$ on a grid.
No input key starts a sweep, and \texttt{solve()} never runs one.
Both functions are exported by the package (Section~\ref{sec:api}).

\subsection{Cost Measures and the Ranking Criteria}
\label{sec:sweep_criteria}

Let $x$, $z$, $x^r$ be the solution of a case.
The cost of day $k$ is
\begin{equation}
  c_k = C_M \sum_{i=1}^{F} x_{i0k}
      + \sum_{i=1}^{F}\sum_{\ell=1}^{L} \bigl( C_{R,i\ell}\, z_{i\ell k} + C_{\mathrm{rep},i\ell}\, x^r_{i\ell k} \bigr),
  \label{eq:d_daycost}
\end{equation}
where $z_{i\ell k} = 0$ for a cell with $C_{R,i\ell} = 0$ (no variable $z$ is built, Section~\ref{sec:repair_cost}).
The cost of the transitory phase and the cost of the operating phase are
\begin{equation}
  J_{\mathrm{init}} = \sum_{k=1}^{H_1} c_k, \qquad
  J_{\mathrm{op}} = \sum_{k=H_1+1}^{T} c_k, \qquad
  \bar J_{\mathrm{op}} = J_{\mathrm{op}} / H_2 .
  \label{eq:d_phasecosts}
\end{equation}
The peak damage is $u = \max_{i,k} \sum_{\ell} \mu_{i\ell k}$, and the damage penalty is $C_D\, u$.
The result keys are \texttt{J\_initialization}, \texttt{J\_op}, \texttt{J\_op\_average} ($= \bar J_{\mathrm{op}}$),
\texttt{peak\_damage} and \texttt{damage\_penalty} (Section~\ref{sec:output}).

A solve with \texttt{objective\_mode = "operating\_average"} minimises $\bar J_{\mathrm{op}} + C_D\,u$ (Section~\ref{sec:objective}).
The transitory phase has no cost in this objective, but it is still constrained.
With \texttt{objective\_mode = "evaluation\_total"} and an \emph{evaluation horizon} $E > H_1$ (days),
the solve minimises $P_E + C_D\,u$, where $P_E$ is the \emph{projected cost}
\begin{equation}
  P_E = J_{\mathrm{init}} + (E - H_1)\,\bar J_{\mathrm{op}} .
  \label{eq:d_projected}
\end{equation}
$P_E$ is the cost of the transitory phase, then of the operating phase repeated until day $E$.
Two result keys carry it:
\texttt{projected\_evaluation\_cost} holds $P_E$, and
\texttt{evaluation\_objective} holds $P_E + C_D\,u$.

A case is \emph{feasible} when its result holds a solution and its status is not \texttt{"infeasible"},
not \texttt{"unbounded"} and not \texttt{"inf\_or\_unbounded"}.
The \emph{MIP gap} is the relative gap between the objective of the best solution found and the best bound
(MIP: mixed-integer program).
A case is \emph{proven} when it is feasible and Gurobi reports the status \texttt{"optimal"}
at the MIP gap used for that case, that is, the run option \texttt{mip\_gap} of Section~\ref{sec:run_options}.
A proven case is optimal within that gap, not exactly optimal.
The sweep report records the requested gap of each case, so a reader sees how strong the proof is.

\emph{Ranking criteria.}
\texttt{sweep\_operating\_horizons} ranks cases by $\bar J_{\mathrm{op}}$.
The loop constraint (Section~\ref{sec:loop}) makes the operating phase repeat,
so $\bar J_{\mathrm{op}}$ is the cost per day of the repeating regime.
It is comparable across values of $H_2$, whereas a raw objective over $T$ days is not.
\texttt{sweep\_horizon\_grid} ranks by $\bar J_{\mathrm{op}}$ when no evaluation horizon is given.
With an evaluation horizon it ranks by $P_E$, which also prices the transitory phase,
so it is the criterion that compares different values of $H_1$.
In both functions the \emph{best proven case} is the proven case with the smallest ranking value,
and the \emph{best feasible case} is the feasible case with the smallest ranking value.
A tie goes to the case evaluated first.

\emph{Remark (what the proof certifies).}
The status \texttt{"optimal"} certifies the objective that the solve minimised.
This objective is $\bar J_{\mathrm{op}} + C_D\,u$ in the mode \texttt{operating\_average} and $P_E + C_D\,u$ in the mode
\texttt{evaluation\_total}, while the ranking value is $\bar J_{\mathrm{op}}$ or $P_E$ alone.
The sweep therefore ranks by cost and uses the damage penalty only inside each solve.

\subsection{\texttt{sweep\_operating\_horizons}}
\label{sec:sweep_operating}

\begin{lstlisting}[language=Python]
sweep_operating_horizons(
    input_path, operating_horizons, *, output_path=None,
    stop_on_gradient=False, gradient_tolerance=1e-3,
    flat_gradients_required=2, minimum_cases=3,
    maximum_mip_gap_for_stopping=None,
    warm_start=False, warm_start_mode="start", n_workers=1,
    objective_mode="operating_average", evaluation_horizon=None) -> dict
\end{lstlisting}

\begingroup\small
\begin{longtable}{@{}p{5.8cm}p{9.2cm}@{}}
\toprule
Argument & Meaning \\
\midrule
{\raggedright \texttt{input\_path}\par} & {\raggedright Input file (YAML, JSON or HDF5). $H_1$ is the first entry of the input key \texttt{H},
  or its value when \texttt{H} is a scalar. The second entry of \texttt{H} is ignored.\par} \\
{\raggedright \texttt{operating\_horizons}\par} & {\raggedright Candidate values of $H_2$: positive integers, at least one, no duplicates.
  The function evaluates them in increasing order.\par} \\
{\raggedright \texttt{output\_path}\par} & {\raggedright Report file, with the suffix \texttt{.yaml} or \texttt{.yml}. The function creates missing
  parent folders. \texttt{None}: no file is written.\par} \\
{\raggedright \texttt{stop\_on\_gradient}, \texttt{gradient\_tolerance}, \texttt{flat\_gradients\_required},
  \texttt{minimum\_cases}, \texttt{maximum\_mip\_gap\_for\_stopping}\par} & {\raggedright Early stop (Section~\ref{sec:sweep_stop}).
  Required: $\texttt{gradient\_tolerance} \ge 0$ and finite, $\texttt{flat\_gradients\_required} \ge 1$,
  $\texttt{minimum\_cases} \ge 2$, and $\texttt{maximum\_mip\_gap\_for\_stopping}$ is \texttt{None} or a finite number $\ge 0$.\par} \\
{\raggedright \texttt{warm\_start}, \texttt{warm\_start\_mode}\par} & {\raggedright Warm start between cases (Section~\ref{sec:warm_starts}).\par} \\
{\raggedright \texttt{n\_workers}\par} & {\raggedright Number of cases solved at the same time (Section~\ref{sec:sweep_parallel}).\par} \\
{\raggedright \texttt{objective\_mode}, \texttt{evaluation\_horizon}\par} & {\raggedright Written into every case input, replacing the values of
  the input file. The mode \texttt{evaluation\_total} needs $E > H_1$.\par} \\
\bottomrule
\end{longtable}
\endgroup

\emph{Behaviour.}
\begin{enumerate}[label=\arabic*.]
  \item The function reads the input file with the reader of \texttt{solve()} (Section~\ref{sec:api}).
    It accepts the same formats and raises the same file errors.
    The data must be a mapping (\texttt{TypeError} otherwise).
  \item The function checks every argument of the table above before it solves any case.
    An invalid argument raises \texttt{ValueError}. An argument error is never recorded as a case error.
    It also raises \texttt{ValueError} when \texttt{H} is a list of a length other than two,
    or when \texttt{warm\_start} is true and $\texttt{n\_workers} > 1$,
    or when \texttt{stop\_on\_gradient} is true and $\texttt{n\_workers} > 1$.
  \item For each candidate $H_2$, the case input is the input data with
    \texttt{H} $= [H_1, H_2]$, \texttt{objective\_mode} and (if given) \texttt{evaluation\_horizon}
    replaced. Nothing else changes: the function does not interpolate time-indexed data.
    A profile whose length differs from $H_2$ (and does not broadcast) is rejected by the input checks of
    Section~\ref{sec:consistency_checks}. This rejection is an error of that case only.
  \item The case is solved as \texttt{solve()} solves an input file, with \texttt{results\_path=None}.
    The sweep keeps the summary of the case (Section~\ref{sec:sweep_report}), not its schedule.
  \item After each case the function appends the case row, annotates the gradient
    (Section~\ref{sec:sweep_stop}), tests the stopping rule, and writes a checkpoint when
    \texttt{output\_path} is given (Section~\ref{sec:sweep_report}).
  \item At the end the function writes the final report and returns it.
\end{enumerate}

\subsection{Early Stop}
\label{sec:sweep_stop}

Let $\bar J_n$ and $H_{2,n}$ be the cost per day and the operating horizon of the $n$-th evaluated case.
When cases $n-1$ and $n$ are both feasible, the \emph{relative gradient} of case $n$ is
\begin{equation}
  \hat g_n = \frac{\bar J_n - \bar J_{n-1}}{(H_{2,n} - H_{2,n-1})\,\max\{\,|\bar J_{n-1}|,\ \varepsilon_{\mathrm{mach}}\}},
  \label{eq:d_gradient}
\end{equation}
with $\varepsilon_{\mathrm{mach}}$ the machine precision of double-precision numbers.
Case $n$ is classified \texttt{"increase"} if $\hat g_n > \texttt{gradient\_tolerance}$,
\texttt{"decrease"} if $\hat g_n < -\texttt{gradient\_tolerance}$, and \texttt{"flat"} otherwise.
If one of the two cases is not feasible, or if $n = 1$, the class is \texttt{"unavailable"}.
A case \emph{qualifies} when it reports a finite MIP gap that is at most its gap threshold.
The threshold of a case is \texttt{maximum\_mip\_gap\_for\_stopping} if that argument is a number,
and otherwise the MIP gap used for that case (the entry \texttt{requested\_mip\_gap} of its row).
With the default, a case qualifies when Gurobi has met the gap that was requested for it.
The gradient of case $n$ is \emph{eligible to stop} when both cases $n-1$ and $n$ qualify.
The function annotates every case with these values, also when \texttt{stop\_on\_gradient} is false.

With \texttt{stop\_on\_gradient = True} the function stops after case $n$ when
$n \ge \texttt{minimum\_cases}$, the gradient of case $n$ is eligible to stop, and one of these holds:
\begin{itemize}
  \item the class of case $n$ is \texttt{"increase"} (stopping reason \texttt{"cost\_increase"});
  \item the class of case $n$ is \texttt{"flat"}, and the last \texttt{flat\_gradients\_required} cases up to $n$
    all have the class \texttt{"flat"} and an eligible gradient (stopping reason \texttt{"flat\_gradient"}).
\end{itemize}
The last candidate in \texttt{operating\_horizons} is always a hard limit.
A case that does not qualify, for example one that stopped at the time limit above its requested gap, is recorded,
but it cannot trigger a stop.

\subsection{Report, Checkpoint and Errors}
\label{sec:sweep_report}

The report is a mapping of plain Python types. The file is written with a safe YAML writer, so it
contains no arrays and no Python objects.

\begingroup\small
\begin{longtable}{@{}p{5.8cm}p{9.2cm}@{}}
\toprule
Key & Content \\
\midrule
{\raggedright \texttt{input}\par} & {\raggedright Path of the input file, as given.\par} \\
{\raggedright \texttt{complete}\par} & {\raggedright \texttt{false} in a checkpoint, \texttt{true} in the final report (also after an early stop).\par} \\
{\raggedright \texttt{last\_completed\_H2}\par} & {\raggedright Largest $H_2$ among the evaluated cases; \texttt{null} before the first case.\par} \\
{\raggedright \texttt{objective}, \texttt{varied\_parameter}\par} & {\raggedright Text that names the minimised objective, and the string \texttt{"H2"}.\par} \\
{\raggedright \texttt{fixed\_dimensions}, \texttt{H1}\par} & {\raggedright $F$, $M$, $L$, $H_1$.\par} \\
{\raggedright \texttt{planned\_operating\_horizons}, \texttt{evaluated\_operating\_horizons}\par} & {\raggedright The sorted candidates, and the
  $H_2$ of the cases in \texttt{cases}.\par} \\
{\raggedright \texttt{warm\_start}, \texttt{n\_workers}\par} & {\raggedright \texttt{enabled}, \texttt{mode}; number of workers.\par} \\
{\raggedright \texttt{stopping\_rule}\par} & {\raggedright \texttt{enabled}, \texttt{relative\_gradient\_tolerance\_per\_H2},
  \texttt{flat\_gradients\_required}, \texttt{minimum\_cases}, \texttt{maximum\_mip\_gap\_for\_stopping}
  (\texttt{null} when the default applies), \texttt{reason} (\texttt{"in\_progress"}, \texttt{"candidate\_range\_exhausted"}, \texttt{"cost\_increase"} or
  \texttt{"flat\_gradient"}), \texttt{stopped\_early} (the run is complete and fewer cases than planned were evaluated).\par} \\
{\raggedright \texttt{cases}\par} & {\raggedright One row per evaluated case, ordered by $H_2$ (rows below).\par} \\
{\raggedright \texttt{formulation\_growth}\par} & {\raggedright Between the first and the last case that carry formulation counts:
  \texttt{from\_H2}, \texttt{to\_H2}, \texttt{added\_time\_steps}, and for each count key $q$ of the row entry
  \texttt{formulation}, the entries \texttt{added\_}$q$ and $q$\texttt{\_per\_added\_time\_step}.
  \texttt{null} when fewer than two cases carry counts or both have the same $T$.\par} \\
{\raggedright \texttt{best\_proven\_H2}, \texttt{best\_proven\_J\_op\_average}\par} & {\raggedright Best proven case (\texttt{null} if none).\par} \\
{\raggedright \texttt{best\_feasible\_H2}, \texttt{best\_feasible\_J\_op\_average}, \texttt{best\_feasible\_status}\par} & {\raggedright   Best feasible case (\texttt{null} if none).\par} \\
\bottomrule
\end{longtable}

\begin{longtable}{@{}p{5.8cm}p{9.2cm}@{}}
\toprule
Entry of a case row & Content \\
\midrule
{\raggedright \texttt{H1}, \texttt{H2}, \texttt{T}, \texttt{dimensions}\par} & {\raggedright Horizons and the sizes $F$, $M$, $L$, $H_1$, $H_2$, $T$.\par} \\
{\raggedright \texttt{status}, \texttt{objective\_mode}, \texttt{evaluation\_horizon}\par} & {\raggedright As in the result of the case.\par} \\
{\raggedright \texttt{objective}, \texttt{J\_op}, \texttt{J\_op\_average}, \texttt{J\_initialization}, \texttt{peak\_damage},
  \texttt{damage\_penalty}, \texttt{operating\_objective}, \texttt{evaluation\_objective}, \texttt{J\_total},
  \texttt{projected\_evaluation\_cost}\par} & {\raggedright Result keys of the case (Section~\ref{sec:output}).\par} \\
{\raggedright \texttt{objective\_bound}, \texttt{mip\_gap}, \texttt{absolute\_mip\_gap}\par} & {\raggedright Best bound; relative gap; the difference
  $\max\{0,\ \texttt{objective} - \texttt{objective\_bound}\}$.\par} \\
{\raggedright \texttt{requested\_mip\_gap}, \texttt{time\_limit\_seconds}, \texttt{optimizer\_seconds}, \texttt{solutions\_found}\par} & {\raggedright   The run options used and the optimizer outcome.\par} \\
{\raggedright \texttt{warm\_start}\par} & {\raggedright The warm-start report of the case (Section~\ref{sec:warm_starts}).\par} \\
{\raggedright \texttt{formulation}\par} & {\raggedright Model size from the \texttt{performance} entry of the result: \texttt{variables},
  \texttt{continuous\_variables}, \texttt{integer\_variables}, \texttt{binary\_variables},
  \texttt{linear\_constraints}, \texttt{general\_constraints}, \texttt{indicator\_constraints},
  \texttt{quadratic\_constraints}, \texttt{nonzeros}.\par} \\
{\raggedright \texttt{timing}\par} & {\raggedright \texttt{model\_construction\_seconds}, \texttt{optimizer\_call\_seconds},
  \texttt{solution\_extraction\_seconds}, \texttt{backend\_wall\_seconds}, \texttt{branch\_and\_bound\_nodes},
  \texttt{simplex\_iterations}, \texttt{work\_units}.\par} \\
{\raggedright \texttt{cost\_gradient\_per\_H2}, \texttt{relative\_cost\_gradient\_per\_H2}, \texttt{gradient\_classification},
  \texttt{gradient\_stopping\_eligible}\par} & {\raggedright Equation~\eqref{eq:d_gradient} and its class (Section~\ref{sec:sweep_stop}).\par} \\
\bottomrule
\end{longtable}
\endgroup

\emph{Checkpoint.}
When \texttt{output\_path} is given, the function rewrites the report after every case.
Until the sweep ends, the report has \texttt{complete = false} and the stopping reason \texttt{"in\_progress"}.
If a job ends before the run is complete, the file is a valid report of the cases evaluated so far.
In parallel mode the function recomputes the gradient annotations over the rows ordered by $H_2$ at each rewrite.

\emph{Error capture.}
The function catches every exception (class \texttt{Exception}) that a case raises after the arguments are accepted,
for example a failed input check of the case or a solver error such as a missing licence.
It logs the error at level ERROR with $H_1$ and $H_2$, records a row with the status
\texttt{"error: "} followed by the message and \texttt{null} in every other entry except the sizes,
and continues with the next case.
Such a row is neither feasible nor proven.
An error does not stop the sweep, and a warm start never uses the result of a failed case (Section~\ref{sec:warm_starts}).

\subsection{Parallel Cases}
\label{sec:sweep_parallel}

With $\texttt{n\_workers} > 1$ the function solves up to \texttt{n\_workers} cases at the same time,
each in its own process with its own Gurobi environment.
Each running solve uses one Gurobi licence token, so a licence with one token allows only $\texttt{n\_workers} = 1$.
The cases do not use each other's results, so the option is refused together with warm starts between cases,
and together with the early stop, which needs the preceding cases before it starts the following ones.
Rows arrive in the order in which the cases finish; the report orders them by $H_2$.

\subsection{\texttt{sweep\_horizon\_grid}}
\label{sec:sweep_grid}

\begin{lstlisting}[language=Python]
sweep_horizon_grid(
    input_path, transitory_horizons, operating_horizons, *,
    output_path=None, warm_start=True, evaluation_horizon=None,
    **sweep_options) -> dict
\end{lstlisting}

The function evaluates every pair $(H_1, H_2)$ of the two lists.
Both lists contain positive integers; the function sorts them and removes repeated values.
If \texttt{evaluation\_horizon} is given, it must exceed every candidate $H_1$ (\texttt{ValueError} otherwise).
The keyword options in \texttt{sweep\_options} are those of Section~\ref{sec:sweep_operating} except
\texttt{objective\_mode}, which the grid sets itself.
Because \texttt{warm\_start} is true by default, a call with $\texttt{n\_workers} > 1$ must also pass
\texttt{warm\_start=False}.

For each $H_1$ in increasing order, the function runs \texttt{sweep\_operating\_horizons} on the input data with that $H_1$
and the list of $H_2$.
It sets \texttt{objective\_mode} to \texttt{"evaluation\_total"} when an evaluation horizon is given,
and to \texttt{"operating\_average"} otherwise.
The first case of each $H_1$ starts without warm start; the following cases use the nearest preceding case of the same $H_1$
that has a solution.
The early stop applies within each $H_1$ and uses $\bar J_{\mathrm{op}}$ as in Section~\ref{sec:sweep_stop}.
Case errors are captured as in Section~\ref{sec:sweep_report}.
The function rewrites the grid report after every case of every $H_1$ (with \texttt{complete = false}) and at the end.

\begingroup\small
\begin{longtable}{@{}p{5.8cm}p{9.2cm}@{}}
\toprule
Key of the grid report & Content \\
\midrule
{\raggedright \texttt{input}, \texttt{complete}\par} & {\raggedright As in Section~\ref{sec:sweep_report}.\par} \\
{\raggedright \texttt{objective}, \texttt{comparison\_metric}\par} & {\raggedright Text that names the minimised objective; the ranking key,
  \texttt{"projected\_evaluation\_cost"} ($P_E$) or \texttt{"J\_op\_average"}.\par} \\
{\raggedright \texttt{mip\_certificate\_metric}\par} & {\raggedright The string \texttt{"objective"}: the proof rule refers to the solved objective
  (remark in Section~\ref{sec:sweep_criteria}).\par} \\
{\raggedright \texttt{evaluation\_horizon}\par} & {\raggedright $E$, or \texttt{null}.\par} \\
{\raggedright \texttt{varied\_parameters}, \texttt{fixed\_dimensions}\par} & {\raggedright \texttt{["H1", "H2"]}; $F$, $M$, $L$.\par} \\
{\raggedright \texttt{planned\_transitory\_horizons}, \texttt{planned\_operating\_horizons}\par} & {\raggedright The sorted candidates.\par} \\
{\raggedright \texttt{warm\_start\_enabled}\par} & {\raggedright The value of \texttt{warm\_start}.\par} \\
{\raggedright \texttt{H1\_sweeps}\par} & {\raggedright The report of each operating sweep, one per $H_1$.
  Its key \texttt{input} names the input file of the grid.\par} \\
{\raggedright \texttt{cases}\par} & {\raggedright All case rows, ordered by $H_1$ and then $H_2$.\par} \\
{\raggedright \texttt{best\_proven}, \texttt{best\_feasible}\par} & {\raggedright \texttt{null}, or a mapping with \texttt{H1}, \texttt{H2}, \texttt{T},
  \texttt{J\_initialization}, \texttt{J\_op\_average}, \texttt{peak\_damage}, \texttt{damage\_penalty},
  \texttt{objective}, \texttt{operating\_objective}, \texttt{evaluation\_objective},
  \texttt{projected\_evaluation\_cost}, \texttt{status}, \texttt{mip\_gap} of the selected case.\par} \\
\bottomrule
\end{longtable}
\endgroup

% ================================================================
\section{Warm Starts}
\label{sec:warm_starts}
% ================================================================

A \emph{warm start} gives Gurobi values for the binary variables, so that the search begins from a known schedule.
A warm start changes the time of a solve, not the model: it does not change the feasible set or the objective.
The package has four mechanisms:
\begin{enumerate}[nosep,label=\arabic*.]
  \item a \emph{start} from the schedule of a previous result (Gurobi attribute \texttt{Start}), with an optional completion;
  \item \emph{hints} from the schedule of a previous result (Gurobi attribute \texttt{VarHintVal});
  \item a \emph{relaxation start}, which needs no previous result;
  \item a \emph{greedy heuristic}, which constructs a schedule from the input and returns it in the form of item 1.
\end{enumerate}

\subsection{Interface}
\label{sec:warm_interface}

\texttt{solve()} (Section~\ref{sec:api}) has two keyword arguments for warm starts:
\texttt{warm\_start} (default \texttt{None}) and \texttt{warm\_start\_mode} (default \texttt{"start"}).
The argument \texttt{warm\_start} is a \emph{source}: a dictionary with the keys of a result,
or the dictionary that the greedy heuristic returns (Section~\ref{sec:greedy}).
A source must carry the arrays \texttt{x} (shape $F' \times (M'+1) \times T'$) and the integers \texttt{H1}, \texttt{H2}
with $T' = H_1' + H_2'$; it may carry \texttt{m} and \texttt{r} (shape $F' \times L' \times T'$).
A source without a solution (\texttt{x} is \texttt{null}) is accepted and has no effect.
A result with a solution found before a time limit is a valid source.
The sizes $F'$, $M'$, $L'$ may differ from the problem's sizes: the values are set for the indices that both have.

\texttt{warm\_start\_mode} is \texttt{"start"}, \texttt{"hint"} or \texttt{"both"}; another value raises \texttt{ValueError}.
The relaxation start is selected by the run options \texttt{relaxation\_warm\_start} and \texttt{relaxation\_time\_limit}
(Section~\ref{sec:run_options}).

\subsection{Day Mapping}
\label{sec:warm_map}

The source has the horizons $H_1'$, $H_2'$ and the problem has $H_1$, $H_2$.
Day $k \in \{1,\ldots,T\}$ of the problem takes the values of the source day
\begin{equation}
  \sigma(k) =
  \begin{cases}
    k, & k \leq \min\{H_1, H_1'\}, \\
    H_1' + 1 + \bigl(\max\{0,\ k - H_1 - 1\} \bmod H_2'\bigr), & \text{otherwise.}
  \end{cases}
  \label{eq:d_daymap}
\end{equation}
So the transitory days keep their index, and the operating days repeat the operating phase of the source cyclically.
A shorter operating phase can seed a longer one.
When $H_1 > H_1'$, the days $H_1'+1,\ldots,H_1$ all take the first operating day of the source.
When the source has the same sizes $F$, $M$, $L$, the mapped schedule copies whole days of the source,
so it keeps the one-activity rows, the mission coverage and the depot-day rows of every day.
It need not satisfy the dynamics or the loop rows of the horizons of the problem.
For each day $k$, train $i$ and cell $(i,\ell)$ the mapped values are
$x_{ijk} \leftarrow x'_{ij\sigma(k)}$ for $j = 0,\ldots,M$, $x^m_{i\ell k} \leftarrow m'_{i\ell\sigma(k)}$
(value $0$ if the source has no \texttt{m}), and $x^r_{i\ell k} \leftarrow r'_{i\ell\sigma(k)}$
when replacement is allowed (\texttt{allow\_replacement}, Section~\ref{sec:replacement}) and the source has \texttt{r}.
Every value is rounded to the nearest integer.
Each other binary variable that these values fix through its defining rows
(the selector $n_{i\ell k}$ of the indicator form, Eq.~\eqref{eq:a_ind_sum}) receives that value.

\subsection{Start, Hints and Completion}
\label{sec:warm_start_modes}

\emph{Mode \texttt{"start"}.}
The solve sets the attribute \texttt{Start} of every mapped variable (Section~\ref{sec:warm_map}).
The continuous variables get no value. Gurobi tries to complete the start by itself, with limited effort,
and discards a start that it cannot make feasible.

\emph{Completion.}
If the source has the key \texttt{complete\_binary\_start} with the value \texttt{true}, the solve completes the start
before it optimises. The time limit of the completion is the source key \texttt{completion\_time\_limit} (seconds, default $60$).
The steps are:
\begin{enumerate}[nosep,label=\arabic*.]
  \item Copy the model. In the copy, fix every binary variable that received a \texttt{Start} value to that value.
  \item Optimise the copy. All binary values are fixed, so the copy has only continuous variables.
  \item If the copy has a solution, set the attribute \texttt{Start} of every variable of the original model to the value in the copy.
    The start is then a complete, feasible point of the original model.
  \item Otherwise record the status of the copy. If the copy is infeasible, compute an \emph{irreducible infeasible subset} (IIS)
    of the copy: a set of rows that is infeasible, and that becomes feasible when any one row is removed (Gurobi function \texttt{computeIIS}).
    Record the names of its linear rows and general constraints, at most $50$ names.
\end{enumerate}
The solve never relaxes or changes the original model, and it never stops because a completion failed:
the optimisation starts with the binary values only.
A completion is requested by the source, not by an argument of \texttt{solve()}.
The greedy heuristic requests it.

\emph{Mode \texttt{"hint"}.}
The solve sets the attribute \texttt{VarHintVal} of every mapped variable to its mapped value, and sets no \texttt{Start} and no completion.
Gurobi uses hints to guide its heuristics and its branching during the whole solve.
A hint needs no feasible point, so it also helps when the mapped schedule does not fit the horizons of the problem
and Gurobi discards it as a start. The hint priority \texttt{VarHintPri} is not set.

\emph{Mode \texttt{"both"}.}
The solve sets \texttt{Start} and \texttt{VarHintVal} on the same variables, and performs the completion if the source requests it.

\subsection{Relaxation Start}
\label{sec:warm_relaxation}

If \texttt{relaxation\_warm\_start} is true and no start or hint was applied (the entry \texttt{applied} of \texttt{warm\_start} is false),
the solve computes a start without a source:
\begin{enumerate}[nosep,label=\arabic*.]
  \item Copy the model, and relax its linear rows with Gurobi's function \texttt{feasRelaxS}, which minimises the sum of the absolute violations.
    The relaxation keeps integrality, variable bounds and general constraints.
  \item Optimise the relaxed copy with the time limit $\max\{1,\ \texttt{relaxation\_time\_limit}\}$ seconds and the MIP gap $0.05$.
  \item If the copy has a solution, set \texttt{Start} of every integer variable of the original model to the rounded value of
    the variable of the same name in the copy.
\end{enumerate}
The relaxed solution is not a feasible point of the original model.
It gives integer values only, and the original model must find its own first solution.

\subsection{Warm-Start Reports}
\label{sec:warm_reports}

The result of a solve carries two dictionaries (Section~\ref{sec:output}).

\begingroup\small
\begin{longtable}{@{}p{5.8cm}p{9.2cm}@{}}
\toprule
Key & Content \\
\midrule
{\raggedright \texttt{warm\_start.applied}\par} & {\raggedright \texttt{true} when at least one value was set.\par} \\
{\raggedright \texttt{warm\_start.values}\par} & {\raggedright Number of values set.\par} \\
{\raggedright \texttt{warm\_start.mode}\par} & {\raggedright The mode used, or \texttt{null} without a source.\par} \\
{\raggedright \texttt{warm\_start.source\_horizon}\par} & {\raggedright \texttt{H1} and \texttt{H2} of the source, or \texttt{null}.\par} \\
{\raggedright \texttt{warm\_start.completion}\par} & {\raggedright \texttt{requested}; \texttt{status} (\texttt{"not\_requested"},
  \texttt{"completed"}, \texttt{"infeasible"}, \texttt{"infeasible\_or\_unbounded"}, \texttt{"time\_limit"} or
  \texttt{"status\_}$q$\texttt{"} with the Gurobi status code $q$); \texttt{fixed\_integer\_variables};
  \texttt{values\_copied}; \texttt{runtime\_seconds}; \texttt{iis\_constraints} (list of names, empty unless the status is \texttt{"infeasible"}).\par} \\
{\raggedright \texttt{relaxation\_warm\_start}\par} & {\raggedright \texttt{enabled}; \texttt{status} (\texttt{"disabled"} when the option is false,
  \texttt{"skipped"} when it is true and a start was applied, otherwise the Gurobi status of the relaxed copy as a string);
  \texttt{relaxation\_objective}; \texttt{values\_copied}; \texttt{runtime\_seconds}.\par} \\
\bottomrule
\end{longtable}
\endgroup

\subsection{Warm Starts in Sweeps}
\label{sec:warm_sweeps}

The sweep functions (Section~\ref{sec:sweeps}) with \texttt{warm\_start=True} use, for each case, the result of the nearest preceding case of the same $H_1$
that has a solution as the source of the solve, with the mode \texttt{warm\_start\_mode}.
A case that raised an error or has no solution is skipped.
The first case of a sequence has no source.
Warm starts between cases need the preceding results, so they are refused together with parallel cases (Section~\ref{sec:sweep_parallel}).

\subsection{Greedy Heuristic}
\label{sec:greedy}

\begin{lstlisting}[language=Python]
build_greedy_warm_start(cfg, *, repair_trigger=0.70)
    -> (dict, GreedyWarmStartReport)
\end{lstlisting}

The function is in the module \texttt{fleet\_management.greedy\_warm\_start}; the package does not export it.
The argument \texttt{cfg} is the normalised configuration of an input (Section~\ref{sec:inputs}).
The trigger \texttt{repair\_trigger} must lie in $[0,1]$ (\texttt{ValueError} otherwise).
The function works for every model of this specification.
It builds a schedule that tries to satisfy the rows of the model. It does not try to give a low cost.
It returns a source for \texttt{solve()}: the dictionary with the keys \texttt{H1}, \texttt{H2}, \texttt{x}, \texttt{m}, \texttt{r},
\texttt{complete\_binary\_start} (\texttt{true}) and \texttt{completion\_time\_limit} ($60$),
and a report.

\emph{Construction pass.}
Keep a proxy state $\hat s_{i\ell}$ for each cell, with the initial value $\mu_{0,i\ell}$.
Treat the days $k = 1,\ldots,T$ in order.
\begin{enumerate}[label=\arabic*.]
  \item Order the missions by decreasing difficulty $\max_{i,\ell} \mu^{\inc}_{i\ell jk}/\tau_{i\ell}$.
    A tie goes to the mission with the larger index.
  \item Take the missions in this order. Give each one to the train $i$ that has no mission yet on day $k$ and
    minimises the pair $\bigl(\max_\ell (\hat s_{i\ell} + \mu^{\inc}_{i\ell jk})/\tau_{i\ell},\ \sum_\ell (\hat s_{i\ell} + \mu^{\inc}_{i\ell jk})/\tau_{i\ell}\bigr)$
    in lexicographic order. A tie goes to the smaller index $i$. Increase $\hat s_{i\ell}$ by $\mu^{\inc}_{i\ell jk}$ for the train and day.
    The input check $F > M$ (Section~\ref{sec:consistency_checks}) guarantees a train for every mission.
  \item Let $r$ be the value of \texttt{repair\_trigger}.
    A train without a mission stays idle, except when some component has $\hat s_{i\ell} > r\,\tau_{i\ell}$.
    Then the train goes to the depot ($x_{i0k} = 1$), every such component is repaired ($x^m_{i\ell k} = 1$),
    and $\hat s_{i\ell} \leftarrow (1-\rho_{i\ell})\,\hat s_{i\ell}$ for it.
\end{enumerate}
The pass schedules no replacement. The proxy uses the whole-state repair map for every cell, whatever its repair law.

\emph{Loop pass.}
After the construction pass the missions are fixed, so the cells are independent except through the depot days of a train.
For each cell, compute its reference trajectory $\zeta$ (Section~\ref{sec:encoding}) for the current schedule,
with the starting values, the starting anchors and the replacement values of the inputs.
The \emph{excess} $e$ of the cell is the largest of
\begin{itemize}[nosep]
  \item $\max\{0,\ \zeta^{q}_{T} - \zeta^{q}_{H_1}\}/M_q$ for every tracked quantity and anchor $q$ with natural range $M_q$ (the loop row, Section~\ref{sec:loop});
  \item the violation of the reliability row of the cell in its exact form, on every day, divided by the natural range of the quantity that the row bounds
    (Sections~\ref{sec:dyn_gaussian}--\ref{sec:dyn_rainflow}).
\end{itemize}
While $e > 10^{-10}$, the pass tries every \emph{free day} $k \in \{H_1+1,\ldots,T\}$, a day on which the train has no mission and the cell has no intervention.
It puts a repair of the cell on the free day (the train goes to the depot on that day) and recomputes the excess.
It keeps the repair that gives the smallest excess; a tie goes to the later day.
If no repair lowers $e$ by more than $10^{-10}$ and replacement is allowed, the same trial is made with a replacement.
If no trial lowers $e$, the pass counts one failure for the cell and goes to the next cell.
Repairs and replacements of a cell never use a day on which the train flies a mission.

\emph{Report.}
\texttt{GreedyWarmStartReport} is a record; its method \texttt{as\_dict()} returns the fields below as plain Python types.
\begingroup\small
\begin{longtable}{@{}p{5.8cm}p{9.2cm}@{}}
\toprule
Field & Content \\
\midrule
{\raggedright \texttt{repairs}, \texttt{replacements}\par} & {\raggedright Number of ones in \texttt{m} and in \texttt{r}.\par} \\
{\raggedright \texttt{repeatability\_feasible}\par} & {\raggedright \texttt{true} when the loop pass counted no failure.\par} \\
{\raggedright \texttt{repeatability\_failures}\par} & {\raggedright Number of cells with a failure.\par} \\
{\raggedright \texttt{maximum\_repeatability\_excess}\par} & {\raggedright Largest raw excess (in the units of its quantity) over all cells after the loop pass.\par} \\
{\raggedright \texttt{maximum\_proxy\_ratio}\par} & {\raggedright Largest value of $\hat s_{i\ell}/\tau_{i\ell}$ in the construction pass and of the reference mean over $\tau_{i\ell}$ in the loop pass.\par} \\
{\raggedright \texttt{assignment\_order}\par} & {\raggedright For each day, the list of missions (indices $1,\ldots,M$) in the order that the pass treated them.\par} \\
\bottomrule
\end{longtable}
\endgroup
A start with failures is still returned. Its completion can then end with the status \texttt{"infeasible"} with an IIS.

% ================================================================
\section{File \& Folder Structure}
\label{sec:files}
% ================================================================

The package \texttt{fleet\_management} has one module per task. The shared Gurobi model and the registry of cell builders
are in \texttt{base.py}; each degradation model has its own module.

\begin{lstlisting}
Fleet_Management/
    pyproject.toml                Python >= 3.10; extra "dashboard" (streamlit)
    README.md
    spec/spec.tex, references.bib this specification
    examples/                     example inputs and regression scripts
    tests/                        pytest suite
    euler/                        two Slurm job scripts of the ETH Euler cluster
    src/fleet_management/
        __init__.py               public API (Section "Public API")
        config.py                 input reading, checks and broadcasting: load_config -> FleetConfig
        solver.py                 solve(): file checks, dispatch, serialisation (YAML, JSON, HDF5)
        horizon_sweep.py          sweep_operating_horizons, sweep_horizon_grid
        greedy_warm_start.py      build_greedy_warm_start
        validator_dashboard.py    Streamlit entry point of the dashboards
        degradation_model/
            base.py               shared model, registry, objective, warm starts, result extraction,
                                  performance counters
            gaussian.py           cell builder "gaussian"
            inverse_gaussian.py   cell builder "inverse_gaussian"
            wiener.py             cell builder "wiener"
            gamma.py              cell builder "gamma"; computation of the cap A_1
            rainflow.py           cell builder "rainflow" (five bounds)
            gamma_utils/
                gamma_replay_validator.py       replay validator
                gamma_stochastic_validator.py   Monte Carlo validator
                gamma_tail_validator.py         exact validator
                gamma_tail_bound.py             moschopoulos_tail_probability
        utils/
            plotter.py            plot_management, plot_horizon_sweep
        validation/
            validator.py          validate
            dashboard_inspector.py, dashboard_gamma.py
\end{lstlisting}

\emph{Cluster jobs.}
The folder \texttt{euler/} holds two Slurm job scripts (Slurm is the job scheduler of the Euler cluster of ETH Z\"urich).
Each runs one script of \texttt{examples/regression/}:
\begin{itemize}[nosep]
  \item \texttt{gamma\_horizon\_convergence.sbatch} runs \texttt{run\_horizon\_sweep.py};
  \item \texttt{vbz\_progressive\_feasibility.sbatch} runs \texttt{run\_progressive\_vbz\_feasibility.py}.
\end{itemize}
Both scripts call only the public API of Section~\ref{sec:api}. The behaviour that this specification defines does not depend on them.

\subsection{Public API}
\label{sec:api}

The package exports the functions of the table.

\begingroup\small
\begin{longtable}{@{}p{6.8cm}p{8.2cm}@{}}
\toprule
Function & Defined in \\
\midrule
\texttt{solve} & this section \\
\texttt{plot\_management} & Section~\ref{sec:plotting} \\
\texttt{plot\_horizon\_sweep} & Section~\ref{sec:plotting} \\
\texttt{sweep\_operating\_horizons} & Section~\ref{sec:sweeps} \\
\texttt{sweep\_horizon\_grid} & Section~\ref{sec:sweeps} \\
\texttt{validate} & Section~\ref{sec:validation} \\
\texttt{validate\_gamma\_replay\_schedule} & Section~\ref{sec:validation_gamma} \\
\texttt{validate\_gamma\_replay\_files} & Section~\ref{sec:validation_gamma} \\
\texttt{validate\_gamma\_stochastic\_schedule} & Section~\ref{sec:validation_gamma} \\
\texttt{validate\_gamma\_schedule} & Section~\ref{sec:validation_gamma} \\
\texttt{validate\_gamma\_schedule\_files} & Section~\ref{sec:validation_gamma} \\
\bottomrule
\end{longtable}
\endgroup

\begin{lstlisting}[language=Python]
solve(input_path, results_path="output.yaml", *,
      warm_start=None, warm_start_mode="start") -> dict
\end{lstlisting}

\texttt{solve()} reads one input file, builds one Gurobi model for the whole fleet, solves it, and returns the result dictionary of
Section~\ref{sec:output}. Arrays in the dictionary are NumPy arrays; Gurobi objects are not stored in it.
\begin{enumerate}[nosep,label=\arabic*.]
  \item The input file must exist (\texttt{FileNotFoundError}) and have the extension \texttt{.yaml}, \texttt{.yml}, \texttt{.json},
    \texttt{.h5} or \texttt{.hdf5} (\texttt{ValueError}).
  \item If \texttt{results\_path} is not \texttt{None}, the solve resolves it before it reads the input.
    A path without an extension gets the extension \texttt{.yaml}.
    The extension selects the format of the result file: YAML, JSON or HDF5 (as for inputs); another extension raises \texttt{ValueError}.
    The folder of the result file must exist (\texttt{FileNotFoundError}) and be writable (\texttt{PermissionError}).
    With \texttt{results\_path=None} no file is written.
  \item The \emph{input reader} reads the file by its extension into a mapping of the keys of Section~\ref{sec:schema}.
    The horizon sweeps and \texttt{validate} use the same reader.
  \item \texttt{load\_config} normalises the mapping and runs the checks of Section~\ref{sec:consistency_checks}.
    A failed check raises \texttt{ValueError}, and no model is built.
  \item The solve builds the model (Section~\ref{sec:builders}), applies the warm start (Section~\ref{sec:warm_starts}), optimises,
    extracts the solution, writes the result file if requested, and returns the dictionary.
\end{enumerate}
Every fleet, whatever its models, goes through the same route: no model has a separate solver.
Gurobi errors, such as a missing licence, are not caught.

\subsection{Cell-Builder Interface}
\label{sec:builders}

The shared model of \texttt{base.py} is built once for the whole fleet. Each cell $(i,\ell)$ is then filled in by the
\emph{cell builder} of its model (the entry of \texttt{model} in the input). Builders are registered in a dictionary by model name.

\begin{lstlisting}[language=Python]
register_cell_builder(name, builder)   # CELL_BUILDERS[name] = builder
get_cell_builder(name)                 # NotImplementedError if name is not registered
dispatch_cell(ctx, i, l)               # get_cell_builder(ctx.model_of[i, l]).add_cell(ctx, i, l)
\end{lstlisting}

The names are \texttt{gaussian}, \texttt{inverse\_gaussian}, \texttt{wiener}, \texttt{gamma} and \texttt{rainflow}.
The error of \texttt{get\_cell\_builder} lists the registered names.
A builder is an object with the attribute \texttt{name} and the hooks below.

\begingroup\small
\begin{longtable}{@{}p{3.6cm}p{11.4cm}@{}}
\toprule
Hook & Duty \\
\midrule
{\raggedright \texttt{prepare(ctx, cfg, cells, opts)}\par} & Called once for each model that occurs in the fleet, with the list \texttt{cells} of the cells of
  that model. It stores the per-cell arrays of the model in the context and creates the variables of the model's tracked quantities,
  anchors and selectors, for these cells only (no variable that no row uses). It may tighten variable bounds, and it sets Gurobi's parameter \texttt{NonConvex} to $2$ when some cell of the model keeps an exact quadratic row. \\
{\raggedright \texttt{add\_cell(ctx, i, l)}\par} & Adds every row of one cell that Section~\ref{sec:constraints} assigns to it: the dynamics in the form selected by
  \texttt{encoding} (Section~\ref{sec:encoding}), the replacement and anchor rows, the repair-cost rows (Section~\ref{sec:repair_cost}),
  the reliability rows in the form selected by \texttt{reliability\_impl} (Section~\ref{sec:linear_forms}), and the loop rows of the full tracked
  state (Section~\ref{sec:loop}). \\
{\raggedright \texttt{extract(ctx, cfg, out)}\par} & Optional. Adds model-specific entries to the result dictionary \texttt{out}, after the shared entries
  are extracted. \\
\bottomrule
\end{longtable}
\endgroup

The \emph{context} \texttt{ctx} carries the sizes ($F$, $L$, $M$, $H_1$, $H_2$, $T$), the Gurobi model, the shared variables
($x$ of shape $F \times (M+1) \times T$; $x^m$ and $x^r$ of shape $F \times L \times T$, with $x^r$ absent when replacement is not allowed;
the mean $\mu$ of shape $F \times L \times T$; $z$ at the cells with $C_R > 0$; $u$ when $C_D > 0$), the per-cell arrays
(model, $\tau$, $\varepsilon$, $\rho$, initial and replacement values, starting anchors), the increment accessor, the resolved run options
and a dictionary \texttt{extras} for storage of a model.
The increment accessor returns, for train $i$, mission $j$ (index from $0$ for the depot), cell $\ell$ and day $k$, the increment after the
phase lookup of Section~\ref{sec:input_increments}.

\emph{Building order.}
\begin{enumerate}[nosep,label=\arabic*.]
  \item Create the model, apply the run options, and create the shared variables.
  \item Call \texttt{prepare} once for each model present, in the alphabetical order of the names.
  \item Build the objective (Section~\ref{sec:objective}) and the rows that do not depend on the model:
    the assignment rows (Section~\ref{sec:assignment}), the depot-day rows (Section~\ref{sec:depot}) and the damage-regularisation rows
    (Section~\ref{sec:damage_reg}).
  \item Call \texttt{add\_cell} for every cell, with $i$ in the outer loop and $\ell$ in the inner loop.
\end{enumerate}
Every builder must define the shared variable $\mu_{i\ell k}$ for all $k = 1,\ldots,T$ of each of its cells, because the damage regularisation
and the output read it. After the optimisation, the shared extraction reads the variables, and then each model's \texttt{extract} hook runs.
A model is added by writing a module with its builder and registering it, by adding its input checks to \texttt{config.py}
(Section~\ref{sec:consistency_checks}), and by adding its section to this specification.

\subsection{Plotting API and Layout}
\label{sec:plotting}

The package has one schedule plot and one sweep plot, both in \texttt{utils/plotter.py}.

\begin{lstlisting}[language=Python]
plot_management(input_file_path, plot_file_path=None) -> None
\end{lstlisting}

\begingroup\small
\begin{longtable}{@{}p{3.8cm}p{11.2cm}@{}}
\toprule
Parameter & Description \\
\midrule
\texttt{input\_file\_path} & Result file of \texttt{solve()}, in YAML, JSON or HDF5 format. A missing file raises
  \texttt{FileNotFoundError}; another extension raises \texttt{ValueError}. \\
\texttt{plot\_file\_path} & Image file, \texttt{.png} or \texttt{.pdf}. Default \texttt{"output.png"}. A path without an extension gets
  \texttt{.png}. Another extension raises \texttt{ValueError}; a missing folder, \texttt{FileNotFoundError}; a folder without write
  permission, \texttt{PermissionError}. \\
\bottomrule
\end{longtable}
\endgroup

The function plots any result that holds a solution, whatever its status: an incumbent found before a time limit is plotted, and
the status is shown in the summary panel. It raises \texttt{ValueError} when the result lacks one of the keys \texttt{F}, \texttt{M},
\texttt{L}, \texttt{H1}, \texttt{H2}, \texttt{tau}, \texttt{mu\_0}, \texttt{mu}, \texttt{x} (names the missing keys),
when a shape differs from Section~\ref{sec:output}, or when $H_1 + H_2 \neq T$.
The keys \texttt{m} and \texttt{r} are read when present and taken as all zero otherwise.
The names of the components come from \texttt{component\_names} (default ``Component $\ell$'' with $\ell$ from $1$);
a list of a length other than $L$ raises \texttt{ValueError}.
The model of each cell comes from the key \texttt{model} (Section~\ref{sec:output}).

\subsubsection{Plot Layout}

The figure has a grid on the left and a summary panel on the right.
\begin{itemize}[nosep]
  \item \textbf{Rows.} One row for each cell $(i,\ell)$, ordered by train and then by component: $F\,L$ rows.
    The label of a row has two lines: ``Train $i$ --- \emph{component name}'' and the model name of the cell
    (\texttt{gaussian}, \texttt{inverse\_gaussian}, \texttt{wiener}, \texttt{gamma} or \texttt{rainflow}).
    A heavier horizontal line separates the trains.
  \item \textbf{Columns.} $T + 1$ columns: ``initial'' (the values $\mu_{0,i\ell}$) and the days $1,\ldots,T$.
    A vertical dashed line after day $H_1$ separates the phases. Above the grid, the labels ``transitory'' and ``operating'' name the two parts.
  \item \textbf{Colour.} The background of a field is a green-to-red scale of $\mu_{i\ell k}/\tau_{i\ell}$, clipped to $[0,1]$ (green $0$, red $1$).
    A colour bar is labelled ``tracked mean / failure threshold''. The values are the tracked means, which are upper bounds of the expected degradation
    (Section~\ref{sec:encoding}).
  \item \textbf{Marks.} Each field of a day carries one mark. The first rule that applies decides:
\end{itemize}

\begingroup\small
\begin{longtable}{@{}p{1.6cm}p{13.4cm}@{}}
\toprule
Mark & Rule for cell $(i,\ell)$ on day $k$ \\
\midrule
\texttt{P} & replacement: $x^r_{i\ell k} = 1$ \\
\texttt{R} & repair: $x^m_{i\ell k} = 1$ \\
\texttt{M}$_j$ & train $i$ flies mission $j$: $x_{ijk} = 1$ with $j \geq 1$ (the mark shows $j$) \\
\texttt{D} & depot day without intervention: $x_{i0k} = 1$, $x^m_{i\ell k} = 0$, $x^r_{i\ell k} = 0$ \\
\texttt{I} & idle: the train has no activity on day $k$ (all $x_{ijk} = 0$) \\
\bottomrule
\end{longtable}
\endgroup

A legend below the grid names the five marks, with a distinct colour for each. A value is $1$ when it is above $0.5$.
The column ``initial'' has no mark.

\emph{Summary panel} (title ``Run summary''). Each line is a label and a value; a line is omitted when its value is missing from the result.
Status (coloured by outcome); Objective; Dimensions $F$, $M$, $L$; Horizon $H_1$, $H_2$, $T$; Models (the distinct model names);
Repair (the distinct values of \texttt{repair\_model}); Objective mode and $J_{\mathrm{op}}/H_2$, if \texttt{objective\_mode} is in the result;
MIP gap in percent; the numbers of variables, continuous variables, integer variables and linear rows; optimizer time;
branch-and-bound nodes; repaired components; replaced components.
The model counts are those of Gurobi's presolved model when the result carries them (\texttt{presolved\_performance}),
and those of the original model otherwise; the label says which.
When the result carries no presolved counts and a file \texttt{runner.out} (a Gurobi log) lies in the folder of the result file,
the plot reads the counts from the lines ``Presolved: $r$ rows, $c$ columns'' and ``Variable types: $a$ continuous, $b$ integer''.
The numbers of repaired and replaced components are the numbers of ones in \texttt{m} and \texttt{r}
(components, not depot visits).

\begin{lstlisting}[language=Python]
plot_horizon_sweep(sweep_report_path, plot_file_path=None) -> None
\end{lstlisting}

\texttt{sweep\_report\_path} is one path, or a list of paths, of reports of \texttt{sweep\_operating\_horizons}.
A report is a YAML or JSON file. Another file type raises \texttt{ValueError}, a missing file raises \texttt{FileNotFoundError},
and an empty list raises \texttt{ValueError}.
The default image file is \texttt{"horizon\_sweep.png"}; the rules for \texttt{plot\_file\_path} are those of \texttt{plot\_management}.
Reports in a list must agree in $F$, $M$, $L$ and $H_1$ (\texttt{ValueError} otherwise), and each needs a list \texttt{cases}.
When several cases have the same $H_2$, the plot keeps the one with a solution, then the smaller $J_{\mathrm{op}}/H_2$, then the larger
objective bound. A report without a case raises \texttt{ValueError}.
The figure has three panels over $H_2$ (the cases are joined by lines that are not extrapolations):
\begin{enumerate}[nosep,label=(\alph*)]
  \item $J_{\mathrm{op}}/H_2$ for each case, with a different mark for a proven case (status \texttt{"optimal"}), a feasible case that is not proven,
    and a case without a feasible solution (including an error row); the best feasible case is marked and the values of the proven cases are written;
  \item the relative MIP gap in percent. When the entry \texttt{maximum\_}\allowbreak\texttt{mip\_gap\_for\_stopping}
    of \texttt{stopping\_rule} is a number, a horizontal line shows it; otherwise a mark shows the requested MIP gap of each case;
  \item the numbers of continuous variables, integer variables and linear constraints.
\end{enumerate}
The title names the sweep and the sizes $F$, $M$, $L$, $H_1$.

% ================================================================
\section{Validation Tools}
\label{sec:validation}
% ================================================================

The validation tools check a saved result against its input. They do not call the optimiser and they read no Gurobi object.
There are three groups: the structural validator for every model (Section~\ref{sec:validation_structural}),
the dashboards (Section~\ref{sec:validation_dashboards}), and the validators of the Gamma model (Section~\ref{sec:validation_gamma}).
All tools read the input with the input reader of \texttt{solve()} and normalise it with \texttt{load\_config} (Section~\ref{sec:api}),
and read a result file in the formats that \texttt{solve()} writes. They read the model of each cell from the result key \texttt{model}.

\subsection{Structural Validator}
\label{sec:validation_structural}

\begin{lstlisting}[language=Python]
validate(input_path, results_path, validation_path="validation.yaml",
         tol=1e-6) -> dict
\end{lstlisting}

\texttt{validate} checks that the saved schedule and states satisfy the rows of the formulation that do not need the solver,
and writes and returns a report. The report file has the extension \texttt{.yaml}, \texttt{.yml} or \texttt{.json}
(a path without an extension gets \texttt{.yaml}; another extension raises \texttt{ValueError}).
The function creates the folder of the report; it raises \texttt{PermissionError} if the folder is not writable.
It raises \texttt{FileNotFoundError} for a missing input or result file, and \texttt{ValueError} for the following:
an unsupported extension, a result without a solution (\texttt{x} is \texttt{null}), an input that fails the checks of
Section~\ref{sec:consistency_checks}, and a result whose sizes $F$, $M$, $L$, $H_1$, $H_2$ or whose array shapes differ from those of the input
(the message names the quantity).
The solver status is not a check: a feasible incumbent found before a time limit is validated like an optimal solution,
and the report echoes the status.

A check passes when its \emph{violation} is at most \texttt{tol}. The checks are:

\begingroup\small
\begin{longtable}{@{}p{3.4cm}p{6.2cm}p{5.4cm}@{}}
\toprule
Name & Statement & Violation \\
\midrule
{\raggedright \texttt{binary}\par} & $x$, $x^m$, $x^r$ take only the values $0$ and $1$ &
  largest distance of an entry from $\{0,1\}$ \\
{\raggedright \texttt{one\_activity}\par} & $\sum_{j=0}^{M} x_{ijk} \leq 1$ for all $i$, $k$ (Section~\ref{sec:assignment}) &
  $\max\{0,\ \max_{i,k}(\sum_j x_{ijk} - 1)\}$ \\
{\raggedright \texttt{mission\_coverage}\par} & $\sum_{i} x_{ijk} = 1$ for $j = 1,\ldots,M$ and all $k$; the depot $j = 0$ is not summed &
  $\max_{j \geq 1,k} |\sum_i x_{ijk} - 1|$ \\
{\raggedright \texttt{depot\_gating}\par} & $x^m_{i\ell k} + x^r_{i\ell k} \leq x_{i0k}$ (Section~\ref{sec:depot}); $x^r = 0$ if \texttt{allow\_replacement} is false &
  largest excess of the left over the right side \\
{\raggedright \texttt{idle}\par} & if the result has \texttt{idle}: $\texttt{idle}_{i\ell k} = x_{i0k} - x^m_{i\ell k} - x^r_{i\ell k}$ &
  $\max |\texttt{idle} - (x_{i0k} - x^m - x^r)|$ \\
{\raggedright \texttt{damage\_rows}\par} & if $C_D > 0$: the rows of Section~\ref{sec:damage_reg} hold for the saved $u$ and $\mu$ &
  largest violation of a row \\
{\raggedright \texttt{loop\_}$q$\par} & the loop row (Section~\ref{sec:loop}) $q_{i\ell T} \leq q_{i\ell H_1}$ for each tracked quantity or anchor $q \in \{\texttt{mu}, \texttt{v}, \texttt{R}, \texttt{K}, \texttt{g\_mu}, \texttt{g\_v}, \texttt{g\_R}, \texttt{g\_K}\}$. The result carries every one of them, with \texttt{NaN} (in a file \texttt{null}) at the cells that do not track it; these entries are skipped &
  $\max\{0,\ \max_{i,\ell}(q_{i\ell T} - q_{i\ell H_1})\}$ \\
{\raggedright \texttt{objective}\par} & the objective recomputed from the saved arrays equals the saved \texttt{objective} &
  {\raggedright relative difference of the two values, with the divisor $\max\{1, |\texttt{objective}|\}$\par} \\
\bottomrule
\end{longtable}
\endgroup

The recomputed objective uses the day cost $c_k$ of Eq.~\eqref{eq:d_daycost} with the saved $x$, $z$, $x^r$ and the costs of the input, and the
saved $u$. For \texttt{objective\_mode} equal to \texttt{total} it is $\sum_{k=1}^{T} c_k + C_D\,u$;
for \texttt{operating\_average} it is $J_{\mathrm{op}}/H_2 + C_D\,u$; for \texttt{evaluation\_total} it is the projected cost $P_E$ of
Eq.~\eqref{eq:d_projected} (Section~\ref{sec:objective}).
The check does not test the optimality of the schedule. It tests that the reported objective is the cost of the reported solution.

The loop checks use the saved tracked states, which satisfy the loop row, and not the reference trajectory:
the reference trajectory of a feasible schedule can violate the loop row (Section~\ref{sec:encoding}).

The report is a mapping with the keys \texttt{status}, \texttt{solver\_objective}, \texttt{recomputed\_objective}, \texttt{tolerance},
\texttt{checks} (a list of mappings with \texttt{name}, \texttt{passed}, \texttt{violation}, \texttt{tolerance}), \texttt{passed}
(all checks passed) and \texttt{max\_violation} (the largest violation).

\subsection{Dashboards}
\label{sec:validation_dashboards}

The dashboards are Streamlit pages (Streamlit is a Python library for web dashboards; it is in the extra \texttt{dashboard}).
The entry point is \texttt{validator\_dashboard.py}; the command is
\texttt{python -m streamlit run src/fleet\_management/validator\_dashboard.py}.
A sidebar selects the mode, the input file, the result file and the tolerance (default $10^{-6}$), and a button loads the pair.
The dashboards read only the input and the result: they need no log file.
A failed load shows the error on the page; no page is drawn from partial data.

\emph{Mode ``Result inspector''} (every model). Tabs:
\begin{enumerate}[nosep,label=\arabic*.]
  \item \emph{Overview}: status, objective, the models present, $H_1$, $H_2$, $T$, the number of reliability violations,
    the largest $\mu/\tau$, and the cells with the smallest margin.
  \item \emph{Reliability}: a table of the margin of every (train, component, day), with filters for train, component and violations only.
  \item \emph{Repeatability}: for each cell and each saved tracked state, the value of the loop row (Section~\ref{sec:loop}).
  \item \emph{Train timeline}: the mean of each component of one train over the days, with the threshold $\tau$ and the marks of
    Section~\ref{sec:plotting}.
  \item \emph{Heatmap}: the margins on the grid of cells and days.
  \item \emph{Component comparison}: one component across the trains.
  \item \emph{Validation}: the checks of \texttt{validate} (Section~\ref{sec:validation_structural}).
  \item \emph{Raw data}: the table behind the tabs, with an export.
\end{enumerate}
The \emph{margin} of a cell on a day is the right side minus the left side of the reliability row of its model and bound, in its exact form,
with the $\tau$, $\varepsilon$ and bound of that cell and the saved tracked states (Sections~\ref{sec:dyn_gaussian}--\ref{sec:dyn_rainflow}).
A margin of at least $-\texttt{tolerance}$ means that the row holds. The exact form is used also when the solve used the tangent or
piecewise-linear form, so the margin shows the true slack.
The result carries every tracked quantity; the dashboard leaves out the entries that are \texttt{NaN}.
The margins of cells with different models and bounds are not comparable in size; only the sign is.

\emph{Mode ``Gamma diagnostic''} (Gamma cells). Tabs: an overview with the counts of repairs and replacements and the outcome of the
replay validator; the tracked mean of each Gamma cell with its cap $A_1(\beta^*\tau,\varepsilon)/\beta^*$ (Section~\ref{sec:dyn_gamma});
the exact tail probability of the true damage against $\varepsilon$ (Section~\ref{sec:validation_gamma}); a heatmap of this tail over the cells and days;
and the raw data.

\subsection{Validators of the Gamma Model}
\label{sec:validation_gamma}

The three validators of this section check the Gamma cells of a result. Each takes the normalised configuration \texttt{cfg} (Section~\ref{sec:inputs})
and the result dictionary, or an input path and a result path, and returns a report (a mapping of plain Python types).
Each raises \texttt{ValueError} when \texttt{cfg} has no Gamma cell, when a required array is missing, has a wrong shape or has a non-finite entry,
and when an array that must be binary is not within \texttt{tolerance} of $\{0,1\}$.
The argument \texttt{tolerance} (default $10^{-8}$) must be positive. With \texttt{raise\_on\_failure=True} a report with \texttt{valid = false} raises
\texttt{AssertionError}. The file versions take also \texttt{report\_path} (default \texttt{None}): the report is written there as YAML or JSON
(\texttt{ValueError} for another suffix), and the folder is created if it is missing.
The tracked state of a Gamma cell is the mean $\mu$ and, under $\ARDone$, its anchor $g^{\mu}$ (Section~\ref{sec:dyn_gamma}).

\emph{Replay validator.}
\begin{lstlisting}[language=Python]
validate_gamma_replay_schedule(cfg, result, *, tolerance=1e-8,
                               raise_on_failure=False) -> dict
validate_gamma_replay_files(input_path, result_path, report_path=None,
                            **kwargs) -> dict
\end{lstlisting}
It replays the bookkeeping of the saved schedule. It reads the keys \texttt{x}, \texttt{m}, \texttt{r} (optional), \texttt{mu}, \texttt{z}, and
the anchor \texttt{g\_mu} of the mean for the cells that use $\ARDone$.
Global violations (strings in \texttt{global\_violations}): a result horizon different from $T$; a train with more than one activity;
a mission not covered exactly once; repair and replacement of a cell on one day; an intervention on a day with a mission or without a depot day.
For each Gamma cell and day $k$ the validator computes the reference trajectory of the mean and of its anchor
(Section~\ref{sec:encoding}) from the saved schedule, with the initial values, the starting anchor and the replacement value of the cell.
The \emph{event} of the day is \texttt{"replacement"}, \texttt{"repair"}, \texttt{"mission"} or \texttt{"idle"}. The errors, by field, are:

\begingroup\small
\begin{longtable}{@{}p{3.6cm}p{11.4cm}@{}}
\toprule
Field & Error \\
\midrule
{\raggedright \texttt{tracked\_mean}\par} & $\max\{0,\ \text{reference mean} - \texttt{mu}_{i\ell k}\}$: a tracked mean below the reference breaks the upper-bound
  property of the one-sided rows (Section~\ref{sec:encoding}). A tracked mean above the reference is slack and is not an error. \\
{\raggedright \texttt{mean\_anchor}\par} & the same for the anchor \texttt{g\_mu}, under $\ARDone$ \\
{\raggedright \texttt{removed\_mean}\par} & $|z_{i\ell k} - x^m_{i\ell k}\,\rho_{i\ell}(\mu_{i\ell,k-1} - g_{i\ell,k-1})|$ computed from the saved states, with $g = 0$ under $\ARDinf$
  (Section~\ref{sec:repair_cost}); $z = 0$ on a day without repair \\
{\raggedright \texttt{cap}\par} & $\max\{0,\ \texttt{mu}_{i\ell k} - A_1(\beta^*\tau,\varepsilon)/\beta^*\}$ (Section~\ref{sec:dyn_gamma}) \\
{\raggedright \texttt{repeatability}\par} & $\max\{0,\ \texttt{mu}_{i\ell T} - \texttt{mu}_{i\ell H_1}\}$, and the same for the saved anchor \\
\bottomrule
\end{longtable}
\endgroup

An error above \texttt{tolerance} adds a mapping \texttt{i}, \texttt{l}, \texttt{k}, \texttt{event}, \texttt{field}, \texttt{expected}, \texttt{saved}, \texttt{error}
to \texttt{violations}. The report has the keys \texttt{validator}, \texttt{valid} (no global violation and no violation), \texttt{solver\_status},
\texttt{dimensions}, \texttt{gamma\_cells}, \texttt{gamma\_ard1\_cells}, \texttt{transitions\_checked}, \texttt{repairs}, \texttt{replacements},
\texttt{maximum\_errors} (the largest error of each field), \texttt{global\_violations}, \texttt{violations} and \texttt{timing}.

\emph{Exact validator} (module \texttt{degradation\_model.\allowbreak gamma\_utils.\allowbreak gamma\_tail\_validator}; the package does not export it).
\begin{lstlisting}[language=Python]
validate_gamma_tail_bound_schedule(cfg, result, *, tolerance=1e-8,
    convolution_tolerance=1e-12, max_series_terms=100000,
    include_steps=False, raise_on_failure=False) -> dict
validate_gamma_tail_bound_files(input_path, result_path, report_path=None,
                                **kwargs) -> dict
\end{lstlisting}
It checks the safety claim itself, for the saved schedule and without the tracked states.
For each Gamma cell it rebuilds the true damage on each day $k = 1,\ldots,T$ as a list of independent Gamma terms, as in the proof of the class of true states
(Section~\ref{sec:dyn_gamma}): every mission adds a term with the exact rate of its increment; a repair multiplies the rate of the repairable terms by $1/(1-\rho)$
(the shape is unchanged; for $\rho = 1$ the repairable terms disappear); under $\ARDone$ the repaired terms join the frozen part of the anchor;
a replacement starts a new list with the replacement term.
An initial or replacement damage without a rate (the inputs \texttt{gamma\_beta\_0} or \texttt{gamma\_beta\_new} are omitted) is a known constant $c$ and
is not a term: the tail is then $\Pr(\text{sum of the terms} > \tau - c)$, which is $1$ if $c > \tau$.
It computes the tail $\Pr(D_k > \tau)$ of the sum of the terms with the Moschopoulos series. The series writes the law of a sum of
independent Gamma variables with different rates as a mixture of Gamma laws with non-negative weights. When all rates are equal, the tail is the Gamma tail.
The partial sum of the series is a lower bound of the tail. The validator adds the mixture weight that remains, which is at most \texttt{convolution\_tolerance},
and so obtains a certified upper bound of the tail. If the series does not reach this tolerance within \texttt{max\_series\_terms} terms,
the validator raises \texttt{RuntimeError}. The argument \texttt{convolution\_tolerance} must lie in $(0,1)$ and \texttt{max\_series\_terms} must be positive.
A day passes when the upper bound is at most $\varepsilon_{i\ell} + \texttt{tolerance}$ and the saved mean is at least the exact mean $\mathbb{E}D_k$ minus \texttt{tolerance}.
Under $\ARDone$ the saved anchor must also be at least the exact mean of the frozen part, minus \texttt{tolerance}.
(The saved values may be larger: they are upper bounds.)
The report adds, to the keys of the replay report, \texttt{minimum\_reliability\_slack} ($\min (\varepsilon - \text{tail upper bound})$), \texttt{worst\_reliability}
(the step that attains it), \texttt{maximum\_remaining\_mass}, the numbers of series terms, and the timing of the tail computations.
With \texttt{include\_steps=True} the report has the list \texttt{steps} with the terms and values of every cell and day.

\emph{Monte Carlo validator.}
\begin{lstlisting}[language=Python]
validate_gamma_stochastic_schedule(cfg, result, *, repetitions=100000,
    random_seed=20260908, maximum_schedule_failure_rate=None,
    confidence_level=0.95, batch_size=20000, tolerance=1e-8,
    raise_on_failure=False) -> dict
validate_gamma_schedule(cfg, result, *, mode=None, repetitions=None,
    random_seed=None, maximum_schedule_failure_rate=None,
    confidence_level=None, batch_size=None, tolerance=1e-8,
    raise_on_failure=False) -> dict
validate_gamma_schedule_files(input_path, result_path, report_path=None,
                              **kwargs) -> dict
\end{lstlisting}
The schedule stays fixed. In each of \texttt{repetitions} runs, the validator draws the true damage of every Gamma cell:
the initial damage, each mission increment and each replacement draw are independent Gamma variables with the inputs' means and rates
(the initial damage under $\ARDone$ is the sum of two independent terms with the means $g^{\mu}_0$ and $\mu_0 - g^{\mu}_0$, the first one being the anchor).
An initial or replacement damage without a rate (the inputs \texttt{gamma\_beta\_0} or \texttt{gamma\_beta\_new} are omitted) is the known constant of its mean.
Each day it applies the repair law of the cell to the sampled path: $D \leftarrow (1-\rho)\,D$ under $\ARDinf$; $D \leftarrow g + (1-\rho)(D - g)$ then $g \leftarrow D$ under $\ARDone$;
a replacement draws the damage of the replacement part and, under $\ARDone$, sets $g$ to it. After each day it records whether $D > \tau$.
A \emph{schedule failure} is a run in which some Gamma cell exceeds its threshold on some day.
The draws use a generator seeded with \texttt{random\_seed} and are made in batches of \texttt{batch\_size} runs; the same arguments give the same report.
\texttt{confidence\_level} must lie in $(0,1)$, \texttt{repetitions} and \texttt{batch\_size} must be positive, and
\texttt{maximum\_schedule\_failure\_rate} (if given) must lie in $[0,1]$.

The report counts the failures of the fleet, of each cell and of the worst day of each cell. Each rate comes with its one-sided Clopper--Pearson
upper confidence bound at \texttt{confidence\_level}. For $f$ failures in $n$ runs the bound is the \texttt{confidence\_level}-quantile of the
Beta$(f+1,\ n-f)$ law, and $1$ if $f = n$. The keys are:

\begingroup\small
\begin{longtable}{@{}p{5.8cm}p{9.2cm}@{}}
\toprule
Key & Content \\
\midrule
{\raggedright \texttt{validator}, \texttt{mode}, \texttt{valid}, \texttt{timing}\par} & {\raggedright Name of the validator, the mode, the verdict, the time used.\par} \\
{\raggedright \texttt{scope}, \texttt{complete\_fleet\_validation}\par} & {\raggedright ``Gamma cells only''; false if the fleet has cells of other models.\par} \\
{\raggedright \texttt{excluded\_cells}\par} & {\raggedright For each cell that is not a Gamma cell: i, l, component, model, reason.\par} \\
{\raggedright \texttt{repetitions}, \texttt{random\_seed}, \texttt{batch\_size}, \texttt{confidence\_level}\par} & {\raggedright The arguments used.\par} \\
{\raggedright \texttt{schedule}\par} & {\raggedright failed\_replays, failure\_rate, failure\_rate\_upper\_confidence\_bound, maximum\_failure\_rate, observed\_target\_met, confidence\_qualified\_target\_met.\par} \\
{\raggedright \texttt{cells}\par} & {\raggedright For each Gamma cell: i, l, component, repair\_model, threshold, analytical\_epsilon, failed\_replays, failure\_rate and its upper confidence bound, worst\_time\_step with the failures, the rate and the bound of that day, maximum\_sampled\_damage, maximum\_threshold\_excess.\par} \\
\bottomrule
\end{longtable}
\endgroup

The report is valid when the observed schedule failure rate is at most the target \texttt{maximum\_}\allowbreak\texttt{schedule\_failure\_rate}, or when no target is given.
The result is an empirical stress test and not a proof: zero failures in $n$ runs does not show that the true rate is below $1/n$.

\texttt{validate\_gamma\_schedule} selects the validators by \texttt{mode}: \texttt{"deterministic"} (the replay validator; the default),
\texttt{"stochastic"} (the Monte Carlo validator), or \texttt{"both"} (a report with the keys \texttt{deterministic} and \texttt{stochastic} and \texttt{valid} true only if both are valid).
Another value raises \texttt{ValueError}. An argument that is \texttt{None} takes the default of the Monte Carlo validator.
The Monte Carlo validator samples Gamma cells only. It lists every other cell in \texttt{excluded\_cells} with the reason ``not a Gamma cell''.

\section{Input File Schema}\label{sec:schema}

\subsection{Format Rules}

The supported formats are \textbf{YAML} (\texttt{.yaml}, \texttt{.yml}), \textbf{JSON} (\texttt{.json}) and \textbf{HDF5} (\texttt{.h5}, \texttt{.hdf5}).
The format follows the extension of the path, in lower case.
Any other extension raises an error.
One reader serves every function that takes an input path.

\begin{itemize}[leftmargin=1.5em]
  \item A file holds one mapping from the keys of Section~\ref{sec:inputs} to values. The three formats hold the same keys with the same meaning.
  \item An array is a nested list in YAML and JSON. The outermost list runs over the first axis of the shape, in the order of Section~\ref{sec:inputs}: train, component, mission, day.
  \item A key that some cell reads is required if it has no default. A key that no cell reads may be omitted. For a key that only some cells read, the entries of the other cells are \texttt{null}.
  \item YAML is read in the safe mode: a file cannot create program objects.
\end{itemize}

\subsection{HDF5 Layout}\label{sec:schema_hdf5}

An HDF5 file has a flat layout: every key is stored at the root of the file, as an attribute or as a dataset.
The reader makes one entry of every root attribute and of every root dataset, so every key of Section~\ref{sec:inputs} can be read from HDF5.
A byte string is decoded as UTF-8.
A text attribute that begins with \texttt{\symbol{123}} or \texttt{[} is parsed as JSON.

\textbf{Input file.}
\begin{itemize}[leftmargin=1.5em]
  \item A scalar number, a flag or a single string is an attribute, or a dataset of shape~$()$.
  \item \texttt{H} is an integer, or two integers, as an attribute or as a dataset.
  \item An array of numbers is a dataset of type \texttt{float64} with the shape of the key. A \texttt{null} entry is \texttt{NaN}.
  \item A selector array and \texttt{component\_names} are datasets of strings, or JSON text. A \texttt{null} entry of a selector is the empty string.
  \item \texttt{gurobi\_params} is JSON text.
\end{itemize}

\textbf{Output file.}
The result keys of Section~\ref{sec:output} are stored by this rule: a number or a string is an attribute; an array of numbers is a dataset of type \texttt{float64}; a list of strings, a mapping or a list of records is JSON text in an attribute.
A key whose value is absent is not stored.
The table below lists the placement of every key.
The arrays \texttt{v}, \texttt{R}, \texttt{K}, \texttt{g\_mu}, \texttt{g\_v}, \texttt{g\_R}, \texttt{g\_K} and \texttt{z} hold \texttt{NaN} where the quantity is not tracked, because a \texttt{float64} dataset has no \texttt{null}.

\begingroup\small
\begin{longtable}{@{}p{4.2cm}p{10.0cm}@{}}
\toprule
Placement & Keys \\
\midrule
\endhead
\rejrow{Attribute: number or string}{\texttt{status}, \texttt{objective}, \texttt{mip\_gap}, \texttt{bound}, \texttt{F}, \texttt{M}, \texttt{L}, \texttt{H} (one or two integers), \texttt{H1}, \texttt{H2}, \texttt{T}, \texttt{objective\_mode}, \texttt{evaluation\_horizon}, \texttt{u}, \texttt{J\_initialization}, \texttt{J\_op}, \texttt{J\_op\_average}, \texttt{peak\_damage}, \texttt{damage\_penalty}, \texttt{operating\_objective}, \texttt{J\_total}, \texttt{projected\_evaluation\_cost}, \texttt{evaluation\_objective}; the selectors \texttt{bound\_method}, \texttt{repair\_model}, \texttt{reliability\_impl} and \texttt{encoding} when they are one string}
\rejrow{Dataset (\texttt{float64})}{\texttt{x}, \texttt{mu}, \texttt{v}, \texttt{R}, \texttt{K}, \texttt{g\_mu}, \texttt{g\_v}, \texttt{g\_R}, \texttt{g\_K}, \texttt{z}, \texttt{m}, \texttt{r}, \texttt{idle}, \texttt{step\_costs}, \texttt{tau}, \texttt{mu\_0}, \texttt{v\_0}}
\rejrow{JSON text in an attribute}{\texttt{model}, \texttt{component\_names}, \texttt{component\_costs}, \texttt{performance}, \texttt{optimization\_progress}, \texttt{warm\_start}, \texttt{relaxation\_warm\_start}; the selectors \texttt{bound\_method}, \texttt{repair\_model}, \texttt{reliability\_impl} and \texttt{encoding} when they differ between cells}
\bottomrule
\end{longtable}
\endgroup

\subsection{Keys by Cell Type}\label{sec:schema_fields}

The table below lists, for each cell type, the keys that the cell requires and the optional keys that it reads in addition to those of the first row.
A transitory profile (\texttt{*\_trans}) is optional for every profile key.
The optional keys of the first row are read at every cell; the others only at cells of the type.

\begingroup\small
\begin{longtable}{@{}p{2.6cm}p{4.7cm}p{6.9cm}@{}}
\toprule
Cell type & Required & Optional (default in Section~\ref{sec:inputs}) \\
\midrule
\endhead
\typerow{Every cell}{\texttt{F}, \texttt{M}, \texttt{H}, \texttt{model}, \texttt{tau}, \texttt{epsilon}, \texttt{rho}, \texttt{mu\_0}, \texttt{mu}, \texttt{C\_M}, \texttt{C\_R}, \texttt{C\_rep}, \texttt{C\_D}}{\texttt{L}, \texttt{repair\_model}, \texttt{reliability\_impl}, \texttt{encoding}, \texttt{tangent\_ref}, \texttt{pwl\_points}, \texttt{anchor\_mu\_0}, \texttt{replacement\_mu}, \texttt{replacement\_as\_new}, \texttt{allow\_replacement}, \texttt{penalty\_type}, \texttt{component\_names}, \texttt{mu\_trans}, the run options}
\typerow{G}{\texttt{v}}{\texttt{v\_0}, \texttt{replacement\_v}, \texttt{anchor\_v\_0}, \texttt{v\_max\_user}, \texttt{v\_trans}}
\typerow{IG}{\texttt{eta}}{none}
\typerow{W}{\texttt{sigma}}{\texttt{v\_0}, \texttt{replacement\_v}, \texttt{v\_max\_user}}
\typerow{Ga}{\texttt{gamma\_beta}}{\texttt{gamma\_beta\_0}, \texttt{gamma\_beta\_new}, \texttt{gamma\_beta\_trans}}
\typerow{RF, Mk}{\texttt{bound\_method} with the value \texttt{markov}}{none}
\typerow{RF, Ca}{\texttt{v}}{\texttt{v\_0}, \texttt{replacement\_v}, \texttt{anchor\_v\_0}, \texttt{v\_max\_user}, \texttt{v\_trans}}
\typerow{RF, Ho}{\texttt{support}}{\texttt{R\_0}, \texttt{R\_new}, \texttt{anchor\_R\_0}, \texttt{support\_trans}}
\typerow{RF, Be}{\texttt{v}, \texttt{support}}{\texttt{v\_0}, \texttt{support\_0}, \texttt{replacement\_v}, \texttt{support\_new}, \texttt{anchor\_v\_0}, \texttt{v\_max\_user}, \texttt{v\_trans}, \texttt{support\_trans}}
\typerow{RF, Ch}{\texttt{cgf}, \texttt{s\_chernoff}}{\texttt{cgf\_0}, \texttt{cgf\_new}, \texttt{anchor\_K\_0}, \texttt{cgf\_trans}}
\bottomrule
\end{longtable}
\endgroup

Every rainflow cell also needs the key \texttt{bound\_method} when its bound is not Cantelli.
\texttt{anchor\_mu\_0} and the other starting anchors are read only at cells with \texttt{repair\_model} equal to \texttt{ard1}.

\subsection{YAML Example (Mixed Models)}\label{sec:schema_example}

The example has $F = 3$ trains, $M = 2$ missions and $L = 2$ components.
The six cells use Gaussian, inverse Gaussian, Wiener, Gamma, rainflow with Bernstein, and rainflow with Chernoff.
It uses an array, a list and a scalar for cell values, a profile of shape $(L, M)$ and a transitory profile of shape $(L, M, H_1)$.
It passes every check of Section~\ref{sec:consistency_checks}.

\begin{lstlisting}
F: 3
M: 2
L: 2
H: [2, 3]               # H1 = 2, H2 = 3, so T = 5
component_names: [Wheelset, Brake pad]

model:
  - [gaussian, inverse_gaussian]
  - [wiener,   gamma]
  - [rainflow, rainflow]
repair_model:
  - [ard1,   ard1]
  - [ardinf, ard1]      # the default is ardinf
  - [ard1,   ard1]
bound_method:           # read at rainflow cells only
  - [null,      null]
  - [null,      null]
  - [bernstein, chernoff]
reliability_impl:       # linear-row cells (IG, Gamma, Chernoff): null
  - [tangent, null]
  - [pwl,     null]
  - [exact,   null]
encoding:
  - [bigm, bigm]
  - [bigm, bigm]
  - [bigm, indicator]

tau: 1.0
epsilon:
  - [0.05, 0.05]
  - [0.05, 0.05]
  - [0.10, 0.10]
tangent_ref: 0.4        # read at the tangent cell only
pwl_points: 6           # read at the pwl cell only
rho: [0.8, 0.7]         # per component, the same for every train

mu_0:
  - [0.12, 0.10]
  - [0.15, 0.08]
  - [0.09, 0.14]
v_0:
  - [0.0004, null]
  - [0.0005, null]
  - [0.0004, null]
support_0:              # Bernstein: width of the initial range
  - [null, null]
  - [null, null]
  - [0.05, null]
anchor_mu_0:            # ard1 cells only; at most mu_0
  - [0.10, 0.06]
  - [null, 0.05]
  - [0.00, 0.10]
anchor_v_0:             # ard1 cells with a variance state
  - [0.0002, null]
  - [null,   null]
  - [0.0002, null]
anchor_K_0:             # ard1 cell with Chernoff; K_0 = s * mu_0 = 0.70
  - [null, null]
  - [null, null]
  - [null, 0.50]

replacement_mu: 0.02
replacement_v:
  - [0.00004, null]
  - [0.00004, null]
  - [0.00004, null]
support_new:
  - [null, null]
  - [null, null]
  - [0.03, null]
replacement_as_new: false
allow_replacement: true

# profiles, axes (L, M): one row per component, one column per mission
mu:
  - [0.06, 0.05]
  - [0.04, 0.07]
mu_trans:               # axes (L, M, H1): days 1 and 2
  - [[0.07, 0.06], [0.06, 0.05]]
  - [[0.04, 0.04], [0.07, 0.07]]
v:
  - [0.0001, 0.0001]
  - [null,   null]
support:
  - [0.05, 0.04]
  - [null, null]
cgf:                    # at s_chernoff = 5; at least s * mu
  - [null, null]
  - [0.22, 0.38]
s_chernoff:
  - [null, null]
  - [null, null]
  - [null, 5.0]

sigma:
  - [null, null]
  - [0.03, null]
  - [null, null]
eta:
  - [null, 0.15]
  - [null, null]
  - [null, null]
v_max_user:
  - [null, null]
  - [0.05, null]
  - [null, null]
gamma_beta:
  - [null, null]
  - [null, 20.0]
  - [null, null]
gamma_beta_new:         # gamma_beta_0 is omitted: the initial damage is known
  - [null, null]
  - [null, 20.0]
  - [null, null]

C_M: 1.0
C_R: [2.0, 1.5]
C_rep: [10.0, 8.0]
C_D: 1.5
penalty_type: inf_norm
objective_mode: evaluation_total
evaluation_horizon: 30

verbose: 0
mip_gap: 0.05
time_limit: 600
fast: false
gurobi_params: {Threads: 4}
progress_interval_seconds: 10
relaxation_warm_start: false
\end{lstlisting}

In the example, the inverse Gaussian cell uses the law \texttt{ard1} with the starting anchor $0.06 \le \mu_0 = 0.10$. The Gamma cell has the starting anchor $0.05 \le \mu_0 = 0.08$ and no \texttt{gamma\_beta\_0}, so that its initial damage is known. The Chernoff cell has $K_0 = 5 \cdot 0.14 = 0.70$ by default, so that its starting anchor $0.50$ is allowed.
The Bernstein cell has the support constant $b = \max(0.05, 0.03, 0.05) = 0.05$, so that $2\kappa b/3 = 0.077 < \tau = 1$ with $\kappa = \ln 10$.

% ================================================================
\section{Open Questions}
\label{sec:open}
% ================================================================

This section lists the design questions that this specification leaves open.
The planned work of the project is listed in its to-do file,
\texttt{Things to change.txt}.

\begin{itemize}[leftmargin=1.5em]
  \item \textbf{Hierarchical failure-probability bounds.}
        Every reliability row bounds the failure probability of one component:
        $\Pr(D_{i\ell k} > \tau_{i\ell}) \leq \varepsilon_{i\ell}$ for all $i, \ell, k$.
        A user may want to state the requirement at a coarser level:
        \begin{itemize}
          \item \textbf{Train level:} $\Pr(\exists\,\ell : D_{i\ell k} > \tau_{i\ell}) \leq \varepsilon_i$
                for all $i, k$, the probability that \emph{any} component of train $i$
                fails. Under independence of the components this reads
                $1 - \prod_\ell (1 - \Pr(D_{i\ell k} > \tau_{i\ell})) \leq \varepsilon_i$,
                a nonlinear row that couples the components of a train.
          \item \textbf{Fleet level:} $\Pr(\exists\,i,\ell : D_{i\ell k} > \tau_{i\ell}) \leq \varepsilon$
                for all $k$, the probability that \emph{any} component of the fleet fails.
                Under independence this row couples all trains.
        \end{itemize}
        Both require either union-bound approximations (conservative and linear) or
        product-form rows (exact and nonlinear). The per-component levels give both
        through the union bound: $\varepsilon_{i\ell} \leq \varepsilon_i / L$ for every
        $\ell$ gives the train level, and $\varepsilon_{i\ell} \leq \varepsilon/(FL)$
        for every cell gives the fleet level.
  \item \textbf{Sampling error of the empirical inputs.}
        Rainflow counting yields \emph{estimates} of the per-mission mean, variance,
        support width and cumulant generating function, not exact values. A guaranteed
        reliability bound needs one-sided confidence margins on these inputs.
  \item \textbf{Gamma cap over the reachable states.}
        The cap of Section~\ref{sec:dyn_gamma} is the largest cap that is safe for every
        member of the class of true states defined there. The states that a schedule can
        reach form a smaller set, because the tracked mean also satisfies the cap on every
        earlier day. For a cell with one rate and $\beta^*\tau$ below the threshold
        $t_1(\varepsilon)$ of Section~\ref{sec:dyn_gamma}, a larger safe cap may exist.
        Its value is not known.
\end{itemize}

% ================================================================
\section{Changelog}
% ================================================================

\emph{Change log (v0.3, panel review).}
ARA1 maintenance removed from the Gaussian model (no literature basis for the
variance extension). Output \texttt{v} shape changed from $F \times V \times 2H$
to $F \times L \times 2H$ with NaN for non-variance components.
$\varepsilon$ bound tightened from $(0, 0.5)$ to $(0, 0.01]$.
IG FSD TODO replaced with verified Proposition (MLR proof).
Tightness remark refined to three-mechanism argument.

\emph{Change log (v0.4, rainflow extension).}
Rainflow degradation model added: distribution-free, with per-mission empirical
mean/variance/support inputs from rainflow counting; reliability certified by a
Cantelli or Bernstein tail bound (new \texttt{tail\_bound} option,
\texttt{"cantelli"} default); exact nonconvex-quadratic and conservative LP
tangent formulations both specified; loop constraint imposed on both tracked
moments, with the argument that moment domination chains across repetitions
while the scalar bound comparison $U_{2H} \leq U_H$ does not; ARD1 only,
mirroring the Gaussian precedent. \texttt{formulation} token renamed
\texttt{"socp"} $\to$ \texttt{"exact"} (\texttt{"socp"} retained as a deprecated
alias with a \texttt{UserWarning}); the previous identification of the
Gaussian/Wiener exact reliability constraint as a rotated SOC was corrected---with
the variance affine in the decisions, the squared form is a nonconvex quadratic,
handled by Gurobi's \texttt{NonConvex=2} support (the \texttt{"quadratic"}
damage penalty remains genuinely convex SOC). Wiener loop-constraint wording
corrected per accepted review item A-02 (moment bounds, not FSD). Version
reconciliation: the v0.3 change log above corresponds to the revision previously
titled ``v0.2 (revised)'' on the title page.
A post-draft panel pass (markers RF-042--RF-045 and in-place corrections)
additionally fixed the loop-induction base case ($m \geq 2$ with
the loop row as the $m=1$ case), the boundary wording of the
$\tau > 2\kappa b/3$ precondition, the a-fortiori centering step in the
Bernstein application, and the problem-class qualifier for IG/Gamma-only
fleets; extended the Bhatia--Davis consistency checks to the initial and
replacement states and the shape-validation list to $b^{\mathrm{inc}}, b^{0},
b^{\mathrm{new}}$; relocated \texttt{tail\_bound} to the model-specific
schema tables; added the fifth plot border style; and corrected the YAML
example's $\varepsilon$ to respect its own $(0, 0.01]$ bound.

\emph{Change log (v0.5, loop-penalty removal).}
The soft loop-constraint penalty $\mathcal{L}_{\mathrm{loop}}$ and its
coefficient $C_P$ were removed throughout (objective, per-model penalty
formulas, total penalty sum, input tables, YAML example, full-problem
statement); the hard loop constraints are unchanged for all five models. The
former three-mechanism big-M tightness remark was revised to its two remaining
cost mechanisms ($C_D\,u$ and $C_R\,z$), and a slack-robustness argument was
added: the tracked states are componentwise upper bounds on the true moments,
each one-sided dynamics row dominates the exact monotone one-step map, so the
hard loop constraint yields $\Phi(\cdot_H) \leq \cdot_H$ for the exact
second-half period map $\Phi$ --- the induction premise of the repeatability
argument --- and monotonicity of the reliability certificates keeps every
future repetition certified even when the tracked states carry slack; the
penalty's only role was a tie-break toward larger sustainability margins.
A subsequent review pass (markers REV-001--REV-022) closed the ARA1
loop-constraint anchor gap by adding the hard rows
$\mu^{\mathrm{last}}_{2H} \leq \mu^{\mathrm{last}}_{H}$ (Wiener--ARA1) and
$\alpha^{\mathrm{last}}_{2H} \leq \alpha^{\mathrm{last}}_{H}$ (Gamma--ARA1),
with a full-state monotonicity derivation showing they restore the chaining
premise $\Phi(s_H) \leq s_H$, and extended the slack-robustness remark
accordingly. The same pass made the idle-as-all-zeros assignment semantics
explicit (the $C_M$ term is a genuine decision cost, not a constant),
corrected the ARA1 repair-cost anchor index ($k \to k{-}1$), removed the
stale duplicate parameter \texttt{alpha\_target} (legacy name of $\tau$),
reconciled the optionality of \texttt{tail\_bound} and the defaulted schema
fields, added the derived constants to the notation table, expanded the two
rainflow FUTURE bullets with pointers to the scoping document's Route-I
refinements, and recorded the ARA1 anchor-convention questions (pre- vs.\
post-repair storage; no reset at replacement) as open items.

\emph{Change log (v1.0, merged package).}
The specification describes the package that merges the branches \texttt{lore}
and \texttt{main}, on \texttt{main}'s code base, and it is authoritative for the
formulation and for the tools. The five degradation models are cell builders of
one registry. The repair laws carry the names of the literature: $\ARDone$, the
reduction of the degradation accumulated since the last intervention
\cite[$\S$3.1, eq.~(4)]{mercier2019}, and $\ARDinf$, the reduction of the whole
degradation, named by analogy with ARA$_\infty$ \cite[eq.~(20), p.~11]{doyen2004}.
The horizon has a transitory and an operating phase, $H = [H_1, H_2]$, and the
loop row bounds the full tracked state, anchors included, at $k = T$ by its value
at $k = H_1$. Every dynamics row is one-sided, in big-M or indicator form. The
rainflow model offers five tail bounds (Markov, Cantelli, Hoeffding, Bernstein,
Chernoff); the quadratic reliability rows have an exact, a tangent and a
piecewise-linear form. The inverse Gaussian row is a cap from a conditional tail
mean, safe for any dependence between missions; the Gamma row is the sharpest
safe cap on the tracked mean over the class of true states, for any mix of rates.
The reliability level, the tangent point and the costs are set per cell. The
horizon sweeps, the warm starts, the schedule plot and the validators are
specified. The input keys follow one table, and unknown or old keys raise an
error. The change markers of earlier versions are removed; git history keeps
them.

% ================================================================
\section*{References}
\addcontentsline{toc}{section}{References}
% ================================================================

\begin{thebibliography}{9}

\bibitem{doyen2004}
L.~Doyen and O.~Gaudoin.
\newblock Classes of imperfect repair models based on reduction of failure
  intensity or virtual age.
\newblock \emph{Reliability Engineering \& System Safety}, 84(1):45--56, 2004.
\newblock doi:10.1016/S0951-8320(03)00173-X. Read in the authors' preprint
  submitted on 25 October 2002, whose page numbers are cited.

\bibitem{mercier2019}
S.~Mercier and I.~T. Castro.
\newblock Stochastic comparisons of imperfect maintenance models for a gamma
  deteriorating system.
\newblock \emph{European Journal of Operational Research}, 273(1):237--248,
  2019.
\newblock doi:10.1016/j.ejor.2018.06.020.

\bibitem{roosta2015}
F.~Roosta-Khorasani and G.~J. Sz\'ekely.
\newblock Schur properties of convolutions of gamma random variables.
\newblock \emph{Metrika}, 78:997--1014, 2015.
\newblock doi:10.1007/s00184-015-0537-9. Read in arXiv:1601.04731v1, whose page
  numbers are cited.

\bibitem{szekely2003}
G.~J. Sz\'ekely and N.~K. Bakirov.
\newblock Extremal probabilities for {G}aussian quadratic forms.
\newblock \emph{Probability Theory and Related Fields}, 126:184--202, 2003.
\newblock doi:10.1007/s00440-003-0262-6.

\end{thebibliography}

\end{document}
